% =====================================================================
%  Báo cáo Tiểu luận 02 — Knowledge Graph, RAG và AI Agent cho y tế và tài chính Việt Nam
%  Nguyễn Duy Nghĩa · D23CTPM01 · B23DCCN600
%
%  Dựng bằng một lệnh:  python build.py
%  build.py sinh macro số liệu, bảng, trích mã nguồn rồi gọi
%      latexmk -lualatex -shell-escape main.tex
%  Hình dữ liệu là PDF vector vẽ bằng matplotlib (figures/); sơ đồ lý thuyết là TikZ (tikz/),
%  được externalize vào tikzcache/.
% =====================================================================
% extreport chỉ có 8-12, 14, 17, 20pt: khai báo 13pt bị bỏ qua và cả bài rơi về 10pt. Dùng 14pt.
\documentclass[a4paper,14pt,oneside]{extreport}

% --------------------------- Ngôn ngữ & font -------------------------
\usepackage{fontspec}
\setmainfont{Times New Roman}
\setsansfont{Segoe UI}
\directlua{luaotfload.add_fallback("toanhoc", {"Cambria Math:mode=node;"})}
\setmonofont{Consolas}[Scale=0.9, RawFeature={fallback=toanhoc}]
\usepackage{polyglossia}
\setmainlanguage{vietnamese}
\gappto\captionsvietnamese{%
  \renewcommand{\listfigurename}{Danh mục Hình vẽ}%
  \renewcommand{\listtablename}{Danh mục Bảng biểu}%
  \renewcommand{\contentsname}{Mục lục}}

% --------------------------- Khổ trang -------------------------------
\usepackage[a4paper,left=3cm,right=2cm,top=2.5cm,bottom=2.5cm,headheight=20pt]{geometry}
\usepackage{setspace}
\onehalfspacing
\setlength{\emergencystretch}{3em}   % cho phép giãn khoảng trắng thêm thay vì để từ dài (tên lớp, \texttt) tràn lề

% --------------------------- Toán học --------------------------------
\usepackage{amsmath,amssymb,amsfonts}
\usepackage{bm}

% --------------------------- Hình & bảng -----------------------------
\usepackage{graphicx}
\graphicspath{{assets/}{figures/}}
\usepackage{booktabs,longtable,array,tabularx,multirow}
\usepackage{float}
\usepackage[table,dvipsnames]{xcolor}
\usepackage[font=normalsize,labelfont=bf,labelsep=period]{caption}
\usepackage{subcaption}
\usepackage{adjustbox}
\usepackage{makecell}   % tiêu đề cột hai dòng trong bảng do build.py sinh
\renewcommand{\cellset}{\renewcommand{\arraystretch}{0.95}}

% --------------------------- Danh sách -------------------------------
\usepackage{enumitem}
\setlist[itemize]{topsep=3pt,itemsep=2pt,parsep=0pt,leftmargin=1.5em}
\setlist[enumerate]{topsep=3pt,itemsep=2pt,parsep=0pt,leftmargin=1.7em}

% --------------------------- Màu & tiêu đề ---------------------------
\definecolor{ptitblue}{HTML}{123A75}
\definecolor{ptitred}{HTML}{C00000}
\definecolor{softgray}{HTML}{F2F4F7}

\usepackage{titlesec}
\titleformat{\chapter}[display]
  {\normalfont\sffamily\bfseries\color{ptitblue}}
  {\Large\chaptertitlename\ \thechapter}{6pt}{\LARGE}
\titlespacing*{\chapter}{0pt}{-18pt}{22pt}
\titleformat{\section}{\normalfont\large\sffamily\bfseries\color{ptitblue}}{\thesection}{0.6em}{}
\titleformat{\subsection}{\normalfont\sffamily\bfseries\color{ptitblue}}{\thesubsection}{0.6em}{}
\titleformat{\subsubsection}{\normalfont\bfseries\itshape}{\thesubsubsection}{0.6em}{}
\setcounter{secnumdepth}{3}
\setcounter{tocdepth}{2}

% --------------------------- Đầu trang / chân trang ------------------
\usepackage{fancyhdr}
\pagestyle{fancy}
\fancyhf{}
\renewcommand{\headrulewidth}{0.4pt}
\fancyhead[L]{\footnotesize\itshape Nhóm lớp: CT \textbar{} Nhóm: 07 -- Nguyễn Duy Nghĩa}
\fancyhead[R]{\footnotesize\itshape Tiểu luận 02: KG, RAG và AI Agent}
\fancyfoot[C]{\thepage}
\fancypagestyle{plain}{\fancyhf{}\fancyfoot[C]{\thepage}\renewcommand{\headrulewidth}{0pt}}

% --------------------------- Tài liệu tham khảo ----------------------
\usepackage[backend=biber,style=numeric,sorting=none,maxbibnames=6]{biblatex}
\DeclareLanguageMapping{vietnamese}{english}
\addbibresource{refs.bib}

% --------------------------- TikZ / pgfplots -------------------------
% =====================================================================
%  Style dùng chung cho mọi hình TikZ/pgfplots của báo cáo.
%  Bảng màu phong phú, khối có bóng đổ và chuyển sắc, macro khối hộp 3D.
% =====================================================================
\usepackage{tikz}
\usepackage{pgfplots}
\pgfplotsset{compat=1.18}
\usetikzlibrary{external,positioning,arrows.meta,matrix,fit,calc,shapes.geometric,shapes.misc,shapes.arrows,
                backgrounds,decorations.pathreplacing,decorations.pathmorphing,shadows,shadows.blur,patterns,3d,
                perspective,fadings,shadings,trees}
\usepgfplotslibrary{groupplots,colormaps,fillbetween}

% Externalize: mỗi hình biên dịch một lần thành tikzcache/<tên>.pdf.
\tikzexternalize[prefix=tikzcache/]
\newcommand{\hinh}[1]{\tikzsetnextfilename{#1}\input{tikz/#1}}

% Cỡ chữ trong hình: lớp 14pt cho \tiny = 6pt, \scriptsize = 8pt, \footnotesize = 10pt; hình lại bị co theo bề
% rộng trang nên chữ còn nhỏ hơn. Trong mọi tikzpicture nâng ba cỡ này lên để chữ hình gần cỡ chữ thân bài.
\makeatletter
\tikzset{every picture/.append style={execute at begin picture={%
  \def\tiny{\@setfontsize\tiny{9.5}{11.5}}%
  \def\scriptsize{\@setfontsize\scriptsize{11}{13}}%
  \def\footnotesize{\@setfontsize\footnotesize{12}{14}}}}}
\makeatother

% ----------------------------------------------------------------- bảng màu
\definecolor{mZero}{HTML}{7F7F7F}
\definecolor{mOne}{HTML}{0072B2}
\definecolor{mTwo}{HTML}{E69F00}
\definecolor{mThree}{HTML}{009E73}
\definecolor{accentRed}{HTML}{C0392B}
\definecolor{lightBlue}{HTML}{DCE8F5}
\definecolor{lightOrange}{HTML}{FCE9CF}
\definecolor{lightGreen}{HTML}{D5F0E6}
\definecolor{lightGray}{HTML}{EDEDED}
% màu tươi cho sơ đồ
\definecolor{cBlue}{HTML}{2F80ED}
\definecolor{cSky}{HTML}{56CCF2}
\definecolor{cTeal}{HTML}{1ABC9C}
\definecolor{cGreen}{HTML}{27AE60}
\definecolor{cLime}{HTML}{A3D977}
\definecolor{cYellow}{HTML}{F2C94C}
\definecolor{cOrange}{HTML}{F2994A}
\definecolor{cRed}{HTML}{EB5757}
\definecolor{cPink}{HTML}{E86FA8}
\definecolor{cPurple}{HTML}{9B51E0}
\definecolor{cIndigo}{HTML}{5B5FE0}
\definecolor{cBrown}{HTML}{A0522D}
\definecolor{cNavy}{HTML}{123A75}
\definecolor{cSlate}{HTML}{4F5D75}

\pgfplotsset{
  a06/.style={
    width=8cm, height=5.6cm, grid=major, grid style={gray!25},
    tick label style={font=\small}, label style={font=\normalsize},
    legend style={font=\small, draw=gray!50, fill=white, fill opacity=0.9, text opacity=1, rounded corners=2pt},
    title style={font=\normalsize\bfseries},
    /pgf/number format/use comma, /pgf/number format/1000 sep={.}},
  m0/.style={mZero, very thick, mark=none},
  m1/.style={mOne, very thick, mark=none},
  m2/.style={mTwo, very thick, mark=none},
  m3/.style={mThree, very thick, mark=none},
}

\tikzset{
  % khối chữ nhật bo góc có bóng mờ và chuyển sắc từ trên xuống
  hop/.style={rounded corners=3pt, draw=#1!75!black, thick, top color=#1!12, bottom color=#1!38,
              blur shadow={shadow blur steps=6, shadow xshift=0.8pt, shadow yshift=-0.8pt, shadow opacity=25},
              align=center, inner sep=4pt, font=\small},
  hop/.default=cBlue,
  khoi/.style={hop=cBlue, minimum height=8mm},
  khoiC/.style={hop=cOrange, minimum height=8mm},
  khoiG/.style={hop=cGreen, minimum height=8mm},
  khoiX/.style={hop=cSlate, minimum height=8mm},
  khoiP/.style={hop=cPurple, minimum height=8mm},
  khoiR/.style={hop=cRed, minimum height=8mm},
  khoiT/.style={hop=cTeal, minimum height=8mm},
  % quả cầu có bóng (nơ-ron, phép toán)
  cau/.style={circle, inner color=white, outer color=#1!55, draw=#1!65!black, thick, text=black!85,
              font=\small\bfseries, inner sep=1pt, minimum size=8mm,
              blur shadow={shadow blur steps=5, shadow opacity=30}},
  cau/.default=cBlue,
  mui/.style={-{Stealth[length=2.6mm]}, very thick, cNavy},
  muiM/.style={-{Stealth[length=2.6mm]}, very thick, #1},
  nhan/.style={font=\small, text=black!75},
  o/.style={draw=gray!70, minimum size=7mm, inner sep=0pt, font=\small},
  tieude/.style={font=\normalsize\bfseries\sffamily, text=cNavy},
}

% \khoibaD{x}{y}{rộng}{cao}{sâu}{màu}: khối hộp 3D, góc dưới-trái-trước tại (x,y),
% chiều sâu vẽ xiên lên phải (0,6 × sâu theo 30 độ). Mặt trên sáng, mặt bên tối.
\newcommand{\khoibaD}[6]{%
  \pgfmathsetmacro{\kdx}{0.6*#5*cos(30)}\pgfmathsetmacro{\kdy}{0.6*#5*sin(30)}%
  \filldraw[draw=#6!60!black, fill=#6!85!black, line join=round] (#1,#2) rectangle ++(#3,#4);
  \filldraw[draw=#6!60!black, fill=#6!45, line join=round] (#1,#2+#4) -- ++(\kdx,\kdy) -- ++(#3,0) -- ++(-\kdx,-\kdy) -- cycle;
  \filldraw[draw=#6!60!black, fill=#6!60!black, line join=round] (#1+#3,#2) -- ++(\kdx,\kdy) -- ++(0,#4) -- ++(-\kdx,-\kdy) -- cycle;
}

% build.py --kiem-tra: mỗi hình externalize là một lượt chạy cả tài liệu với -halt-on-error, nên một lỗi ở chương
% sau làm hỏng mọi hình mới. Lượt kiểm tra định nghĩa \kiemtranhanh để vẽ hình tại chỗ và liệt kê mọi lỗi một lần.
\ifdefined\kiemtranhanh\AtBeginDocument{\tikzexternaldisable\let\tikzexternalenable\relax}\fi

% --------------------------- Mã nguồn --------------------------------
\usepackage[most]{tcolorbox}
\tcbuselibrary{external}
\tcbset{shield externalize}
\usepackage{minted}
\setminted{fontsize=\small,linenos,numbersep=6pt,breaklines,frame=lines,
           framesep=2mm,baselinestretch=1.05,tabsize=4}
\usemintedstyle{friendly}
\renewcommand{\listingscaption}{Đoạn mã}
\renewcommand{\listoflistingscaption}{Danh mục Đoạn mã nguồn}
\counterwithin{listing}{chapter}
% \maNguon{module__ten}{chú thích}: đoạn mã build.py trích nguyên văn từ src/TieuLuan 02/{tl2,scripts,app}/<module>.py,
% hàm, lớp, phương thức (Lop.ham) hoặc hằng <ten>; mã trong báo cáo vì vậy luôn là mã đã chạy ra số liệu.
% Không dùng [H]: đoạn mã dài phải được ngắt sang trang sau thay vì tràn chân trang.
\newcommand{\maNguon}[2]{%
  \begin{listing}[!htbp]
  \inputminted{python}{code/#1.py}
  \caption{#2}
  \label{code:#1}
  \end{listing}}
% \maDai{module__ten}{chú thích}: đoạn mã dài hơn một trang, cho phép ngắt trang giữa chừng.
\newcommand{\maDai}[2]{%
  \par\medskip\noindent\begin{minipage}{\linewidth}\captionof{listing}{#2}\label{code:#1}\end{minipage}\nopagebreak
  \inputminted{python}{code/#1.py}\par\medskip}
% \maTruyVan{số}{chú thích}: truy vấn SPARQL thứ <số> do build.py ghi ra từ outputs/metrics/kg.json.
\newcommand{\maTruyVan}[2]{%
  \begin{listing}[H]
  \inputminted{sparql}{code/sparql_q#1.rq}
  \caption{#2}
  \label{code:sparql#1}
  \end{listing}}

% --------------------------- Liên kết --------------------------------
\usepackage{hyperref}
\hypersetup{
  colorlinks=true, linkcolor=ptitblue, citecolor=ptitblue,
  urlcolor=ptitblue, filecolor=ptitblue,
  pdftitle={Tiểu luận: Đồ thị tri thức, RAG và tác tử thông minh cho y tế và tài chính Việt Nam},
  pdfauthor={Nguyễn Duy Nghĩa — B23DCCN600 — Nhóm lớp CT — Nhóm 07},
  pdfsubject={Phát triển các Hệ thống Thông minh},
  bookmarksnumbered=true
}
\usepackage{bookmark}
\usepackage{microtype}

% --------------------------- Lệnh riêng ------------------------------
\usepackage{etoolbox}
\newcommand{\nhanxet}[1]{\par\smallskip\noindent\textbf{Nhận xét:} #1\par\smallskip}
\newtcolorbox{phathien}[1]{
  breakable,
  colback=softgray, colframe=ptitblue, boxrule=0.7pt, arc=2pt,
  left=8pt, right=8pt, top=6pt, bottom=6pt,
  fonttitle=\bfseries\sffamily, title={#1}
}
% Vết suy luận ReAct do build.py sinh (generated/trace_<miền>.tex): mỗi bước một nhãn màu và đoạn treo lề.
\newtcolorbox{vetsuyluan}[1]{
  breakable, colback=white, colframe=ptitblue, boxrule=0.8pt, arc=3pt,
  left=6pt, right=6pt, top=5pt, bottom=5pt, fontupper=\small,
  fonttitle=\bfseries\sffamily, title={#1}
}
\newcommand{\vetNhan}[2]{\colorbox{#1}{\makebox[1.9cm][c]{\textcolor{white}{\sffamily\bfseries\footnotesize #2}}}}
\newcommand{\vetBuoc}[3]{\par\noindent\hangindent=2.45cm\hangafter=1\makebox[2.45cm][l]{\vetNhan{#1}{#2}}#3\par\vspace{2pt}}
\newcommand{\vetCauHoi}[1]{\vetBuoc{ptitblue}{CÂU HỎI}{\textit{#1}}}
\newcommand{\vetThought}[1]{\vetBuoc{cIndigo}{THOUGHT}{#1}}
\newcommand{\vetAction}[4]{\vetBuoc{cOrange!90!black}{ACTION}{\texttt{#1}(\,#4\,) \textcolor{gray}{-- loại #2, #3 ms}}}
\newcommand{\vetObs}[1]{\vetBuoc{cTeal!85!black}{OBSERV.}{{\footnotesize #1}}}
\newcommand{\vetVerify}[4]{\vetBuoc{cRed}{VERIFY}{trích dẫn hợp lệ: #1; trung thực số học: #2; dùng \textbf{#3}; tổng thời gian #4 s}}
\newcommand{\vetDong}[1]{\par\noindent #1\par}
% \hinhpdf{tệp}{rộng}{chú thích}{nhãn}: hình dữ liệu matplotlib trong figures/.
\newcommand{\hinhpdf}[4]{%
  \begin{figure}[H]\centering
  \IfFileExists{figures/#1.pdf}{\includegraphics[width=#2\linewidth]{#1}}%
    {\fbox{\parbox[c][3cm][c]{0.8\linewidth}{\centering chưa vẽ: \detokenize{#1}}}\PackageWarning{solieu}{Thieu so lieu 'fig/#1'}}
  \caption{#3}\label{fig:#4}
  \end{figure}}
% \haihinh{tệp1}{chú thích 1}{tệp2}{chú thích 2}{chú thích chung}{nhãn}: hai hình nửa trang cạnh nhau (a), (b).
\newcommand{\haihinh}[6]{%
  \begin{figure}[H]\centering
  \begin{subfigure}[t]{0.49\linewidth}\centering\includegraphics[width=\linewidth]{#1}\caption{#2}\end{subfigure}\hfill
  \begin{subfigure}[t]{0.49\linewidth}\centering\includegraphics[width=\linewidth]{#3}\caption{#4}\end{subfigure}
  \caption{#5}\label{fig:#6}
  \end{figure}}
% \hinhtikz{tệp}{chú thích}{nhãn}: sơ đồ TikZ trong tikz/, co theo bề rộng trang nếu cần.
\newcommand{\hinhtikz}[3]{%
  \begin{figure}[H]\centering
  \begin{adjustbox}{max width=\linewidth}\hinh{#1}\end{adjustbox}
  \caption{#2}\label{fig:#3}
  \end{figure}}
% \anhUI{tệp}{chú thích}{nhãn}: ảnh chụp Web App (assets/ui/<tệp>.png) do app/screenshots.py sinh.
\newcommand{\anhUI}[3]{%
  \begin{figure}[H]\centering
  \IfFileExists{assets/ui/#1.png}{\fbox{\includegraphics[width=0.97\linewidth]{ui/#1.png}}}%
    {\fbox{\parbox[c][4cm][c]{0.9\linewidth}{\centering chưa chụp: \detokenize{#1}}}\PackageWarning{solieu}{Thieu so lieu 'ui/#1'}}
  \caption{#2}\label{fig:#3}
  \end{figure}}
% \bang{tệp}{chú thích}{nhãn}: bảng do build.py sinh trong generated/.
\newcommand{\bang}[3]{%
  \begin{table}[H]\centering\caption{#2}\label{tab:#3}
  \IfFileExists{generated/#1.tex}%
    {{\small\setlength{\tabcolsep}{4.5pt}\begin{adjustbox}{max width=\linewidth}\input{generated/#1}\end{adjustbox}}%
     \IfFileExists{generated/#1_note.tex}{\par\smallskip{\footnotesize\input{generated/#1_note}\par}}{}}%
    {\fbox{chưa sinh bảng: \detokenize{#1}}\PackageWarning{solieu}{Thieu so lieu 'tab/#1'}}
  \end{table}}

% Số liệu: \SL{tệp/khoá/.../lá}, sinh tự động từ outputs/metrics/*.json bởi build.py.
\newcommand{\defSL}[2]{\expandafter\def\csname sl@#1\endcsname{#2}}
\newcommand{\SL}[1]{\ifcsname sl@#1\endcsname\csname sl@#1\endcsname\else
  \PackageWarning{solieu}{Thieu so lieu '#1'}\textbf{??}\fi}
\input{generated/so_lieu}

\renewcommand{\arraystretch}{1.16}
\newcolumntype{L}[1]{>{\raggedright\arraybackslash}p{#1}}
\newcolumntype{C}[1]{>{\centering\arraybackslash}p{#1}}
\newcolumntype{R}[1]{>{\raggedleft\arraybackslash}p{#1}}

\begin{document}

% =====================================================================
%  TRANG BÌA (mẫu Assignment 04: khung lùi 6,5 mm, nội dung đặt tuyệt đối trong khung)
% =====================================================================
\begin{titlepage}
\thispagestyle{empty}
\tikzexternaldisable
\begin{tikzpicture}[remember picture,overlay]
  \node[anchor=north west, inner sep=0pt] at ([xshift=6.5mm,yshift=-6.5mm]current page.north west)
    {\includegraphics[width=197mm,height=284mm]{cover-frame.png}};
  \node[anchor=north, inner sep=0pt, text width=161mm, align=center]
    at ([yshift=-(6.5mm+20mm+6mm)]current page.north) {%
    {\bfseries\fontsize{15}{21}\selectfont HỌC VIỆN CÔNG NGHỆ BƯU CHÍNH VIỄN THÔNG\par}
    \vspace{4mm}
    {\bfseries\fontsize{14}{19}\selectfont KHOA CÔNG NGHỆ THÔNG TIN 1\par}
    \vspace{4mm}
    {\bfseries\fontsize{13.5}{18}\selectfont HỌC PHẦN PHÁT TRIỂN CÁC HỆ THỐNG THÔNG MINH\par}
    \vspace{12mm}
    \includegraphics[height=40mm]{ptit-logo.png}\par
    \vspace{10mm}
    {\bfseries\fontsize{22}{27}\selectfont\addfontfeature{LetterSpace=2}BÁO CÁO TIỂU LUẬN MÔN HỌC\par}
    \vspace{5mm}
    {\bfseries\fontsize{13}{19.5}\selectfont ĐỀ TÀI: ĐỒ THỊ TRI THỨC (KNOWLEDGE GRAPH) VÀ RAG,\\
     TRIỂN KHAI TÁC TỬ THÔNG MINH (AI AGENT) CHO\\
     Y TẾ VÀ TÀI CHÍNH VIỆT NAM\par}
    \vspace{11mm}
    {\fontsize{13}{16}\selectfont
     \renewcommand{\arraystretch}{1.8}
     \begin{tabular}{@{}p{52mm}l@{}}
       \textbf{Lớp:}                  & D23CTPM01 \\
       \textbf{Nhóm:}                 & Nhóm lớp CT -- Nhóm 07 \\
       \textbf{Họ và tên:}            & Nguyễn Duy Nghĩa \\
       \textbf{Mã sinh viên:}         & B23DCCN600 \\
       \textbf{Giảng viên hướng dẫn:} & PGS.TS Trần Đình Quế \\
       \textbf{Học kỳ:}               & Học kỳ 1 năm học 2026 -- 2027 \\
     \end{tabular}\par}};
  \node[anchor=south, inner sep=0pt] at ([yshift=6.5mm+10mm+14mm]current page.south)
    {\bfseries\itshape\fontsize{13}{18}\selectfont Hà Nội, Tháng 10/2026};
\end{tikzpicture}
\tikzexternalenable
\null
\end{titlepage}

% =====================================================================
%  DANH MỤC
% =====================================================================
\pagenumbering{roman}
\setcounter{page}{1}
\tableofcontents
\cleardoublepage
\phantomsection
\addcontentsline{toc}{chapter}{Danh mục Hình vẽ}
\listoffigures
\cleardoublepage
\phantomsection
\addcontentsline{toc}{chapter}{Danh mục Bảng biểu}
\listoftables
\cleardoublepage
\phantomsection
\addcontentsline{toc}{chapter}{Danh mục Đoạn mã nguồn}
\listoflistings

% =====================================================================
%  PHẦN THÂN
% =====================================================================
\cleardoublepage
\pagenumbering{arabic}
\setcounter{page}{1}
% =====================================================================
%  TỔNG QUAN BÁO CÁO
% =====================================================================
\chapter*{Tổng quan báo cáo}
\addcontentsline{toc}{chapter}{Tổng quan báo cáo}
\markboth{Tổng quan báo cáo}{Tổng quan báo cáo}

Mô hình ngôn ngữ lớn (LLM) viết trôi chảy nhưng không biết mình đang sai. Hỏi một trợ lý thông thường
\emph{``đang uống metformin thì tối nay uống bia được không''}, nó có thể trả lời bằng một đoạn văn nghe rất hợp
lý mà không nêu được căn cứ nào, hoặc nêu một con số không có trong bất kỳ tài liệu nào. Trong y tế và tài
chính, một câu trả lời như vậy còn nguy hiểm hơn là không trả lời. Tiểu luận này nghiên cứu ba kỹ thuật bổ sung
cho LLM chính phần ``biết'' và ``kiểm chứng'' mà nó thiếu: \textbf{đồ thị tri thức} (Knowledge Graph, KG) lưu tri
thức có cấu trúc và có nguồn; \textbf{sinh tăng cường truy xuất} (Retrieval-Augmented Generation, RAG) đưa văn
bản gốc vào ngữ cảnh trước khi sinh; và \textbf{tác tử thông minh} theo vòng lặp ReAct phối hợp KG, RAG với các
mô hình học sâu chuyên biệt để trả lời câu hỏi nhiều bước.

\section*{Mục tiêu và phạm vi}
Đề bài yêu cầu trình bày khái niệm và kiến trúc của KG, các thế hệ RAG, rồi áp dụng vào hai miền của Việt Nam:
quản lý \textbf{bệnh đái tháo đường típ 2} và phân tích \textbf{rổ cổ phiếu VN30 cùng giá vàng SJC}, cuối cùng
triển khai một AI Agent có giao diện web. Báo cáo này làm trọn chuỗi đó trên máy cá nhân, không gọi một API
LLM bên ngoài nào:

\begin{itemize}
  \item Viết kho tri thức \SL{retrieval_grid/n_docs} tài liệu Markdown (\SL{retrieval_grid/n_chars} ký tự) tổng
        hợp từ Quyết định 5481/QĐ-BYT, ADA, các Nghị định về vàng và báo cáo doanh nghiệp, kèm bộ
        \SL{retrieval_grid/n_questions} cặp hỏi đáp có đáp án nằm nguyên văn trong kho để chấm điểm truy xuất.
  \item Dựng hai đồ thị tri thức có bản thể học: \texttt{VietnamDiabetesKG} (\SL{kg/health/nodes} thực thể,
        \SL{kg/health/edges} quan hệ) và \texttt{VietnamFinanceKG} (\SL{kg/finance/nodes} thực thể,
        \SL{kg/finance/edges} quan hệ), xuất RDF, suy diễn RDFS và truy vấn SPARQL.
  \item So sánh định lượng 5 chiến lược chia đoạn $\times$ 5 bộ truy xuất, 3 mô hình nhúng, 7 cỡ đoạn,
        hai kỹ thuật tiền truy xuất bằng LLM (HyDE, Multi-Query) và bốn kiến trúc từ Naive RAG tới GraphRAG.
  \item Huấn luyện trên GPU hai mô hình học sâu làm công cụ cho tác tử: MLP chẩn đoán nguy cơ đái tháo đường
        (\SL{diabetes/data/n} hồ sơ) và BiLSTM dự báo lợi suất cho bốn cổ phiếu và vàng quy đổi.
  \item Ghép tất cả thành tác tử ReAct có bảy công cụ, bộ kiểm chứng trích dẫn và con số, đánh giá trên
        \SL{dem/agent_scen} kịch bản, và đóng gói thành Web App Flask.
\end{itemize}

\section*{Phân bổ CPU và GPU}
Máy thực nghiệm có một GPU laptop 8 GB và một CPU 32 luồng. Nguyên tắc phân bổ là: việc nào là phép nhân ma
trận lớn lặp lại nhiều lần thì chạy GPU, việc nào là tra cứu, duyệt đồ thị hay tính toán nhỏ thì để CPU, vì chi
phí chép dữ liệu qua PCIe lớn hơn phần tính toán. Bảng~\ref{tab:moitruong} ghi cụ thể từng thành phần.

\begin{table}[H]\centering
\caption{Môi trường thực nghiệm và phân bổ tính toán}\label{tab:moitruong}
{\small\begin{tabular}{L{4.6cm}L{2.0cm}L{7.6cm}}
\toprule
Thành phần & Thiết bị & Lý do \\
\midrule
Qwen2.5-1.5B-Instruct (bf16) & GPU & sinh văn bản tự hồi quy, mỗi token một lượt qua 28 tầng Transformer \\
bge-reranker-v2-m3 (bf16) & GPU & chấm 30 cặp (câu hỏi, đoạn) mỗi truy vấn bằng cross-encoder 568 triệu tham số \\
multilingual-e5 (fp16 khi chung tiến trình với LLM) & GPU & mã hóa câu hỏi và toàn bộ đoạn khi lập chỉ mục \\
Huấn luyện MLP và BiLSTM & GPU & lô 1024 và 128 mẫu, lan truyền ngược \\
BM25, TF-IDF, chia đoạn & CPU & tính trên chỉ mục thưa, dưới 1 ms mỗi truy vấn \\
Hai KG (NetworkX), RDF, SPARQL, suy diễn & CPU & duyệt đồ thị vài chục nút, rdflib chỉ chạy CPU \\
Suy luận MLP, BiLSTM trong tác tử & GPU & đã nạp sẵn trên GPU cùng LLM; MC Dropout chạy 200 mẫu một lô \\
Vẽ hình, dựng báo cáo & CPU & matplotlib và LuaLaTeX \\
\midrule
\multicolumn{3}{L{14.6cm}}{GPU: NVIDIA GeForce RTX 4060 Laptop 8 GB; CPU: AMD Ryzen 9 7945HX, 32 luồng; RAM
40 GB. Python 3.13, PyTorch 2.13 + CUDA 12.6, sentence-transformers, NetworkX, rdflib, Flask.} \\
\bottomrule
\end{tabular}}
\end{table}

Một ràng buộc thực tế đã định hình kiến trúc: trên Windows, trình điều khiển WDDM tính cả bộ nhớ GPU đã cấp
phát vào hạn mức \emph{commit} của hệ điều hành. Nạp LLM sau embedder và reranker khiến tiến trình vượt hạn mức và
dừng đột ngột khi đọc tệp trọng số. Cách khắc phục được giữ lại trong mã: nạp LLM trước, lưu sẵn reranker ở dạng
bf16 trên đĩa và dùng embedder fp16 khi chung tiến trình với LLM; cả ba mô hình khi đó cùng chạy trong khoảng 8 GB.

\section*{Cấu trúc báo cáo}
Báo cáo bám theo bảy phần của đề bài. \textbf{Chương 1} nêu khái niệm KG, RAG và vai trò của cơ sở tri thức.
\textbf{Chương 2} đi vào cấu tạo KG: thành phần, mô hình bộ ba RDF, bản thể học, suy diễn và so sánh Graph DB
với Vector DB. \textbf{Chương 3} trình bày ba thế hệ RAG, từ Naive RAG, các chiến lược chia đoạn, Advanced RAG tới
GraphRAG, mỗi phần đều có thực nghiệm trên kho tri thức của bài. \textbf{Chương 4} xây hệ thống y tế cho bệnh
đái tháo đường típ 2, \textbf{Chương 5} xây hệ thống tài chính cho VN30 và vàng SJC. \textbf{Chương 6} ghép thành
tác tử ReAct, đánh giá và triển khai Web App. \textbf{Chương 7} kết luận và nêu hướng phát triển.

\begin{phathien}{Quy ước về số liệu và mã nguồn}
Không con số thực nghiệm nào trong báo cáo được gõ tay. Mỗi kịch bản trong \texttt{scripts/} ghi kết quả ra
\texttt{outputs/metrics/*.json}; \texttt{build.py} đọc các tệp đó để sinh macro số liệu, bảng và vết suy luận, rồi
trích nguyên văn từng hàm được trích dẫn từ mã nguồn trong \texttt{src/TieuLuan 02}. Hình dữ liệu vẽ bằng
matplotlib từ cùng các tệp JSON; sơ đồ lý thuyết vẽ bằng TikZ. Chạy lại thực nghiệm rồi dựng lại báo cáo thì mọi
con số, bảng và hình đổi theo.
\end{phathien}

% =====================================================================
%  CHƯƠNG 1 — MỞ ĐẦU VỀ KNOWLEDGE GRAPH VÀ RAG
% =====================================================================
\chapter{Mở đầu về Knowledge Graph và RAG}\label{chap:mo-dau}

Chương này đặt nền cho toàn bài: đồ thị tri thức là gì, RAG là gì, vì sao hai kỹ thuật này bổ sung đúng chỗ yếu
của LLM, và chúng được ghép với nhau như thế nào trong hệ thống của tiểu luận. Hình~\ref{fig:hesinhthai} là bản đồ
của cả báo cáo; mỗi khối trên hình tương ứng một chương phía sau.

\hinhtikz{t1_ecosystem}{Tổng quan hệ sinh thái tri thức đa miền của tiểu luận: nguồn tri thức thô được tổ chức
thành hai kho bổ sung cho nhau là đồ thị tri thức và chỉ mục vector; tác tử ReAct phối hợp hai kho với các mô hình
học sâu để trả lời người dùng}{hesinhthai}

\section{Khái niệm đồ thị tri thức và RAG}

\subsection{Khái niệm đồ thị tri thức (Knowledge Graph, KG)}
Đồ thị tri thức là một đồ thị dữ liệu dùng để tích lũy và truyền đạt tri thức về thế giới thực, trong đó
\textbf{nút} biểu diễn thực thể và \textbf{cạnh} biểu diễn quan hệ giữa các thực thể~\cite{hogan2021kg}. Khác một
cơ sở dữ liệu quan hệ, KG không cố định lược đồ thành các bảng: muốn thêm một loại quan hệ mới chỉ cần thêm một
nhãn cạnh mới. Khác một văn bản, mỗi mẩu tri thức trong KG là một đơn vị nguyên tử, máy đọc được và có thể truy
vấn chính xác.

Về hình thức, gọi $\mathcal{E}$ là tập thực thể, $\mathcal{R}$ là tập quan hệ (vị từ). Một KG là tập các
\textbf{bộ ba} (chủ thể, vị từ, đối tượng):
\begin{equation}
  \mathcal{G} = \{ (s, p, o) \mid s, o \in \mathcal{E},\ p \in \mathcal{R} \} \subseteq \mathcal{E} \times \mathcal{R} \times \mathcal{E}.
  \label{eq:kg}
\end{equation}
Ví dụ, sự kiện ``metformin tương tác nguy hiểm với rượu bia'' trong \texttt{VietnamDiabetesKG} được lưu thành bộ
ba
\[
  (\textsf{Metformin},\ \textsf{TUONG\_TAC\_NGUY\_HIEM},\ \textsf{Ruou\_Bia}).
\]
Đi kèm tập bộ ba là một
\textbf{bản thể học} (ontology) quy định thực thể thuộc lớp nào, vị từ nào được nối những lớp nào; bản thể học
biến một đồ thị nhãn tự do thành tri thức có ngữ nghĩa, để máy kiểm tra được và suy ra được tri thức mới
(Chương~\ref{chap:kg}).

Ba tính chất làm KG có giá trị với hệ thống thông minh:
\begin{itemize}
  \item \textbf{Tường minh và truy vết được.} Mỗi cạnh trong hai KG của bài mang thuộc tính \texttt{source} trỏ về
        tài liệu gốc; câu trả lời dựa trên KG luôn nêu được căn cứ.
  \item \textbf{Đa bước.} Câu hỏi ``Fed tăng lãi suất thì vàng SJC ra sao'' không có câu trả lời trong bất kỳ
        đoạn văn đơn lẻ nào, nhưng là một đường đi qua 3--4 cạnh trên đồ thị nhân quả tài chính.
  \item \textbf{Có ràng buộc.} Bản thể học chặn được bộ ba vô nghĩa như ``rau xanh là thuốc đầu tay của metformin''
        ngay khi nạp dữ liệu, trước khi nó lan vào câu trả lời.
\end{itemize}

\subsection{Khái niệm sinh tăng cường truy xuất (RAG)}
RAG được Lewis và cộng sự đề xuất năm 2020~\cite{lewis2020rag} để giải bài toán tri thức chuyên sâu: thay vì bắt
mô hình nhớ mọi thứ trong tham số, ta cho nó \emph{tra cứu} một kho văn bản ngoài trước khi sinh. Một hệ RAG có
hai thành phần: bộ truy xuất $p_\eta(z \mid x)$ chọn các đoạn văn $z$ liên quan tới câu hỏi $x$, và bộ sinh
$p_\theta(y \mid x, z)$ viết câu trả lời $y$ dựa trên cả câu hỏi lẫn đoạn văn. Xác suất câu trả lời là tổng theo
các đoạn văn được truy xuất:
\begin{equation}
  p(y \mid x) \approx \sum_{z \in \operatorname{top-}k\, p_\eta(\cdot \mid x)} p_\eta(z \mid x)\, p_\theta(y \mid x, z).
  \label{eq:rag}
\end{equation}
Trong thực hành hiện nay, các đoạn top-$k$ được nối vào prompt và LLM sinh một lần; công thức~\eqref{eq:rag} vẫn
là cách hiểu đúng: chất lượng câu trả lời bị chặn trên bởi chất lượng truy xuất. Nếu đoạn văn chứa đáp án không
nằm trong top-$k$, bộ sinh dù mạnh tới đâu cũng chỉ còn cách đoán.

Hình~\ref{fig:bonho} đặt RAG trong khung hai loại bộ nhớ. \textbf{Bộ nhớ tham số} là trọng số LLM: rộng, nén, học
một lần, không sửa được từng mẩu và không chỉ ra được nguồn. \textbf{Bộ nhớ phi tham số} là kho tri thức ngoài:
cập nhật bằng cách thêm hay sửa tài liệu, mỗi câu trả lời trỏ được về đoạn văn cụ thể. RAG nối hai bộ nhớ qua bộ
truy xuất.

\hinhtikz{t1_memory}{Hai loại bộ nhớ của một hệ RAG: bộ nhớ tham số nằm trong trọng số LLM, bộ nhớ phi tham số là
kho tri thức ngoài; bộ truy xuất là cầu nối}{bonho}

\subsection{Quá trình phát triển}
Biểu diễn tri thức bằng đồ thị không mới. Mạng ngữ nghĩa có từ thập niên 1960; Web ngữ nghĩa của Berners-Lee năm
2001~\cite{bernerslee2001semantic} chuẩn hóa nó thành RDF, OWL và SPARQL; năm 2012 Google đưa thuật ngữ
``Knowledge Graph'' vào sản phẩm tìm kiếm. Nhánh thứ hai là mô hình ngôn ngữ: kiến trúc Transformer
năm 2017~\cite{vaswani2017attention} mở đường cho LLM, rồi RAG năm 2020 nối LLM với kho văn bản. Hai nhánh gặp
nhau từ 2023~\cite{pan2024unifying}: ReAct~\cite{yao2023react} cho LLM gọi công cụ trong vòng suy luận, GraphRAG năm
2024~\cite{edge2024graphrag} dùng đồ thị và cộng đồng để trả lời câu hỏi tổng hợp (Hình~\ref{fig:dongthoigian}).

\hinhtikz{t1_timeline}{Dòng thời gian của biểu diễn tri thức và RAG: từ mạng ngữ nghĩa, Web ngữ nghĩa, Google
Knowledge Graph tới Transformer, RAG, ReAct và GraphRAG}{dongthoigian}

\section{Vai trò và ý nghĩa của Knowledge Base và RAG}
Cơ sở tri thức (Knowledge Base, KB) là tập tri thức của một miền được tổ chức để máy truy cập được; trong bài này
KB có hai dạng song song: kho văn bản Markdown (dạng phi cấu trúc, phục vụ RAG) và hai KG (dạng có cấu trúc).
Giá trị của KB với một hệ thống dùng LLM nằm ở năm điểm.

\begin{enumerate}
  \item \textbf{Giảm ảo giác.} Ảo giác (hallucination) là khi mô hình sinh nội dung trôi chảy nhưng sai hoặc không
        có căn cứ~\cite{ji2023hallucination}. Ở Mục~\ref{sec:agent-ketqua}, cùng mô hình Qwen2.5-1.5B trả lời
        \SL{dem/agent_scen} câu hỏi kiểm thử: khi không có công cụ, câu trả lời chỉ chứa
        \SL{agent/agg/all/coverage_llm_onlyPct}\% dữ kiện bắt buộc; khi có KB qua tác tử, tỉ lệ này là
        \SL{agent/agg/all/coveragePct}\%.
  \item \textbf{Cập nhật không cần huấn luyện lại.} Nghị định 232/2025/NĐ-CP sửa Nghị định 24/2012 về vàng; với KB,
        chỉ cần thêm một thực thể và một cạnh \texttt{SUA\_DOI}, không phải tinh chỉnh mô hình.
  \item \textbf{Truy vết nguồn.} Trong y tế, một khuyến nghị không có nguồn thì không dùng được. Mọi cạnh trong
        hai KG đều mang mã tài liệu nguồn (\SL{kg/health/with_source} trên \SL{kg/health/edges} cạnh y tế,
        \SL{kg/finance/with_source} trên \SL{kg/finance/edges} cạnh tài chính).
  \item \textbf{Bảo mật dữ liệu.} Hồ sơ xét nghiệm của người bệnh không rời máy: LLM, embedder, reranker và cả hai
        KG chạy cục bộ.
  \item \textbf{Suy luận có kiểm soát.} Quy tắc nghiệp vụ (chống chỉ định, dấu tác động vĩ mô) nằm trong KG dưới
        dạng tường minh; LLM chỉ diễn đạt lại, không phải tự suy ra.
\end{enumerate}

Bảng~\ref{tab:sosanh-kb} so sánh ba cách đưa tri thức miền vào một hệ thống dùng LLM.

\begin{table}[H]\centering
\caption{So sánh LLM đơn lẻ, tinh chỉnh (fine-tuning) và hệ thống tích hợp KB + RAG}\label{tab:sosanh-kb}
{\small\begin{tabular}{L{3.3cm}L{3.5cm}L{3.5cm}L{3.9cm}}
\toprule
Tiêu chí & LLM đơn lẻ & Tinh chỉnh & KB + RAG (+ KG) \\
\midrule
Nguồn tri thức & tham số, cố định lúc huấn luyện & tham số, học thêm dữ liệu miền & kho ngoài, đọc lúc trả lời \\
Cập nhật & không & huấn luyện lại, tốn GPU & thêm, sửa tài liệu hoặc bộ ba \\
Truy vết nguồn & không & không & có, tới từng đoạn hoặc cạnh \\
Ảo giác & cao với tri thức hiếm & giảm nhưng vẫn còn & thấp nếu có bộ kiểm chứng \\
Câu hỏi đa bước & tự suy, khó kiểm soát & tự suy & duyệt đường trên KG \\
Chi phí triển khai & thấp & cao & trung bình: chỉ mục + truy xuất \\
Độ trễ & một lượt sinh & một lượt sinh & thêm truy xuất, rerank \\
Phù hợp & hội thoại chung & đổi văn phong, định dạng & y tế, pháp lý, tài chính \\
\bottomrule
\end{tabular}}
\end{table}

Tinh chỉnh và RAG không loại trừ nhau: tinh chỉnh dạy mô hình \emph{cách} trả lời (văn phong, định dạng trích dẫn),
RAG cung cấp \emph{nội dung} để trả lời. Với một mô hình nhỏ như Qwen2.5-1.5B, Chương~\ref{chap:agent} cho thấy mô
hình hiểu nội dung bằng chứng tốt nhưng thường quên ghi số hiệu trích dẫn; đó chính là phần việc mà tinh chỉnh
có thể bù.

\section{Bài toán và cách tiếp cận của tiểu luận}
Đề bài đặt hai bài toán ứng dụng có tính chất rất khác nhau, và chính sự khác nhau đó làm rõ khi nào cần KG,
khi nào cần RAG, khi nào cần mô hình học sâu.

\begin{itemize}
  \item \textbf{Y tế} là miền của \emph{ràng buộc}: tiêu chuẩn chẩn đoán có ngưỡng cụ thể, thuốc có chống chỉ định,
        thực phẩm có chỉ số đường huyết. Câu trả lời phải đúng từng con số và có nguồn. KG lưu ràng buộc, RAG cung
        cấp lời giải thích, MLP ước lượng nguy cơ từ hồ sơ xét nghiệm.
  \item \textbf{Tài chính} là miền của \emph{quan hệ nhân quả và bất định}: tỷ giá tác động lên cổ phiếu qua nhiều
        kênh, có kênh trái dấu nhau; giá tương lai không biết trước. KG lưu chiều tác động, dữ liệu thị trường kiểm
        chứng lại chiều đó, BiLSTM dự báo kèm khoảng bất định.
\end{itemize}

Cả hai miền dùng chung một lõi: lớp \texttt{KnowledgeGraph} với bản thể học, cùng một kho văn bản và đường ống
truy xuất, cùng một LLM cục bộ và cùng một tác tử ReAct. Mỗi thành phần được đánh giá riêng trước khi ghép:
truy xuất ở Chương~\ref{chap:rag}, KG ở Chương~\ref{chap:kg}, \ref{chap:yte}, \ref{chap:taichinh}, mô hình học
sâu ở Chương~\ref{chap:yte}, \ref{chap:taichinh}, và toàn hệ thống ở Chương~\ref{chap:agent}.

% =====================================================================
%  CHƯƠNG 2 — KIẾN TRÚC KNOWLEDGE GRAPH VÀ CƠ SỞ TRI THỨC
% =====================================================================
\chapter{Kiến trúc Knowledge Graph và cơ sở tri thức}\label{chap:kg}

Chương~\ref{chap:mo-dau} định nghĩa KG là một tập bộ ba. Chương này mở tập bộ ba đó ra: một KG được cấu tạo từ
những gì, được viết ra theo chuẩn RDF thế nào, bản thể học thêm vào đó những ràng buộc và luật suy diễn ra sao,
và khi nào nên lưu tri thức trong một cơ sở dữ liệu đồ thị thay vì một cơ sở dữ liệu vector. Mọi ví dụ đều lấy
từ hai KG thật của bài, cài trong \texttt{tl2/kg\_base.py}.

\section{Các thành phần cấu tạo của Knowledge Graph}
Một KG có ba loại thành phần (Hình~\ref{fig:thanhphan}):

\begin{itemize}
  \item \textbf{Thực thể (entity)} là nút của đồ thị, đại diện cho một đối tượng có danh tính: một bệnh
        (\textsf{DTD\_Tip2}), một thuốc (\textsf{Metformin}), một món ăn (\textsf{Com\_Trang}), một chính sách
        (\textsf{ND\_232\_2025}), một cổ phiếu (\textsf{FPT}). Mỗi thực thể có mã định danh duy nhất, một
        \textbf{kiểu} là một lớp trong bản thể học, nhãn tiếng Việt để hiển thị và danh sách bí danh để liên kết
        với câu hỏi tự nhiên (``tiểu đường'', ``đtđ'', ``type 2'' đều trỏ về \textsf{DTD\_Tip2}).
  \item \textbf{Quan hệ (relation)} là cạnh có hướng, gán nhãn bằng một vị từ:
        $\textsf{Metformin} \xrightarrow{\textsf{TUONG\_TAC\_NGUY\_HIEM}} \textsf{Ruou\_Bia}$. Giữa hai thực thể có
        thể có nhiều cạnh với vị từ khác nhau, nên cấu trúc dữ liệu phù hợp là \emph{đa đồ thị có hướng}
        (\texttt{networkx.MultiDiGraph}~\cite{hagberg2008networkx}).
  \item \textbf{Thuộc tính (attribute)} gắn lên nút (mã ICD-10 \texttt{E11}, chỉ số GI 73 của cơm trắng, ngưỡng
        126 mg/dL) hoặc lên cạnh (mức độ \texttt{cao}, hậu quả, nguồn trích dẫn \texttt{h03-ruou-bia}).
\end{itemize}

\hinhtikz{t2_components}{Ba thành phần của một KG, minh họa trên \texttt{VietnamDiabetesKG}: thực thể là nút
có kiểu và thuộc tính, quan hệ là cạnh có nhãn vị từ, thuộc tính có thể nằm trên cả nút lẫn cạnh}{thanhphan}

Thuộc tính trên cạnh là chỗ hai mô hình dữ liệu đồ thị phổ biến khác nhau. \textbf{Đồ thị thuộc tính có nhãn}
(Labeled Property Graph, LPG), mô hình của Neo4j, cho phép gắn trực tiếp cặp khóa--giá trị lên cạnh.
\textbf{RDF} chỉ có bộ ba, nên muốn nói ``cạnh này có mức độ cao'' phải biến chính cạnh thành một thực thể, kỹ
thuật gọi là \emph{reification} (Hình~\ref{fig:lpgrdf}). Lớp \texttt{KnowledgeGraph} của bài làm việc theo LPG
khi xây dựng và truy vấn, rồi xuất sang RDF có reification khi cần suy diễn và SPARQL; nhờ vậy dùng được thế mạnh
của cả hai.

\hinhtikz{t2_lpg_rdf}{Cùng một sự kiện ``metformin tương tác nguy hiểm với rượu bia, mức độ cao'' trong mô hình
LPG (thuộc tính nằm trên cạnh) và mô hình RDF (cạnh được reification thành một \texttt{rdf:Statement})}{lpgrdf}

Bản thể học được cài bằng hai lớp dữ liệu nhỏ. \texttt{Predicate} mô tả một vị từ: miền (domain) là lớp của chủ
thể, miền giá trị (range) là lớp của đối tượng, cùng ba tính chất logic đối xứng, bắc cầu, nghịch đảo.
\texttt{Ontology} giữ cây lớp dưới dạng từ điển \emph{lớp $\to$ lớp cha} và trả lời câu hỏi ``lớp $A$ có phải
là một loại $B$ không'' bằng cách đi ngược lên gốc (Đoạn mã~\ref{code:kg_base__Ontology}).

\maNguon{kg_base__Predicate}{Lớp \texttt{Predicate}: vị từ kèm miền, miền giá trị và tính chất logic}
\maNguon{kg_base__Ontology}{Lớp \texttt{Ontology}: cây phân cấp lớp và phép kiểm tra \texttt{is\_a}}

Điểm quan trọng nhất của lõi KG là mọi bộ ba đều được \textbf{kiểm tra ràng buộc trước khi nạp}
(Đoạn mã~\ref{code:kg_base__KnowledgeGraph.validate}): chủ thể phải thuộc miền của vị từ, đối tượng phải thuộc
miền giá trị, xét theo cả lớp con. Thêm $(\textsf{Rau\_Xanh}, \textsf{CHI\_DINH\_DAU\_TAY}, \textsf{Metformin})$
bị từ chối ngay vì rau xanh không phải một bệnh. Mục~\ref{sec:rangbuoc} thử bốn bộ ba cố tình sai như vậy.

\maNguon{kg_base__KnowledgeGraph.validate}{Kiểm tra ràng buộc miền và miền giá trị khi thêm một bộ ba}

Để LLM đọc được tri thức, mỗi bộ ba được \textbf{diễn đạt thành câu} tiếng Việt kèm các thuộc tính cạnh
(Đoạn mã~\ref{code:kg_base__KnowledgeGraph.verbalize}). Đây là định dạng mà tác tử đưa vào prompt ở
Chương~\ref{chap:agent}, ví dụ ``Metformin —[tương tác nguy hiểm với]→ Rượu bia (ethanol) (severity: cao; effect:
tăng nguy cơ nhiễm toan lactic\dots)''.

\maNguon{kg_base__KnowledgeGraph.verbalize}{Diễn đạt bộ ba thành câu tiếng Việt để đưa vào ngữ cảnh LLM}

\section{Mô hình biểu diễn bộ ba RDF}
RDF (Resource Description Framework) là chuẩn W3C để mô tả tài nguyên trên Web bằng bộ
ba~\cite{w3c2014rdf}. Mỗi thành phần của bộ ba là một \textbf{IRI} (định danh tài nguyên toàn cục), một
\textbf{literal} (giá trị có kiểu như số, chuỗi có nhãn ngôn ngữ) hoặc một \textbf{nút trống}. Chủ thể là IRI hoặc
nút trống, vị từ luôn là IRI, đối tượng là bất kỳ loại nào. Một đồ thị RDF đơn giản là \emph{tập hợp} các bộ ba
(Hình~\ref{fig:bobardf}); hai đồ thị gộp lại bằng phép hợp tập hợp, điều làm RDF đặc biệt hợp để tích hợp dữ liệu
từ nhiều nguồn.

\hinhtikz{t2_triple}{Bộ ba RDF (chủ thể, vị từ, đối tượng) viết bằng IRI đầy đủ trong không gian tên
\texttt{ex:} của bài, và đồ thị RDF như một tập các bộ ba}{bobardf}

Phép ánh xạ từ KG của bài sang RDF (Đoạn mã~\ref{code:kg_base__KnowledgeGraph.to_rdf}) đi theo bốn bước:

\begin{enumerate}
  \item Mỗi lớp thành một \texttt{rdfs:Class}, quan hệ lớp cha thành \texttt{rdfs:subClassOf}.
  \item Mỗi vị từ thành một \texttt{rdf:Property} kèm \texttt{rdfs:domain}, \texttt{rdfs:range}, nhãn tiếng Việt và
        các tiên đề OWL \texttt{owl:SymmetricProperty}, \texttt{owl:TransitiveProperty}, \texttt{owl:inverseOf}.
  \item Mỗi thực thể có bộ ba \texttt{rdf:type} trỏ tới lớp, \texttt{rdfs:label} và một bộ ba cho mỗi thuộc tính nút.
  \item Mỗi cạnh thành một bộ ba; nếu cạnh có thuộc tính thì thêm một \texttt{rdf:Statement} mang
        \texttt{rdf:subject}, \texttt{rdf:predicate}, \texttt{rdf:object} và từng thuộc tính.
\end{enumerate}

\maDai{kg_base__KnowledgeGraph.to_rdf}{Xuất KG sang RDF bằng rdflib, giữ thuộc tính cạnh bằng reification}

Bảng~\ref{tab:kgstats} tổng hợp kích thước hai KG. Số bộ ba RDF lớn hơn nhiều số cạnh vì mỗi thực thể sinh thêm
bộ ba kiểu và nhãn, mỗi cạnh có thuộc tính sinh thêm khối reification.

\bang{tab_kg_stats}{Quy mô hai đồ thị tri thức của tiểu luận}{kgstats}

\section{Bản thể học và biểu diễn tri thức hình thức}
Theo định nghĩa kinh điển của Gruber, bản thể học là \emph{``một đặc tả tường minh của một khái niệm
hóa''}~\cite{gruber1993ontology}: nó liệt kê các loại sự vật tồn tại trong một miền và các quan hệ được phép giữa
chúng. Trong thực hành Web ngữ nghĩa, bản thể học được viết bằng RDFS và OWL~\cite{w3c2012owl} và gồm ba lớp
tri thức:

\begin{itemize}
  \item \textbf{Phân cấp lớp}: \textsf{SafeFood} $\sqsubseteq$ \textsf{Food} $\sqsubseteq$ \textsf{Substance}
        $\sqsubseteq$ \textsf{MedicalEntity}. Hình~\ref{fig:ontology} là cây lớp của \texttt{VietnamDiabetesKG}
        với đúng \SL{kg/health/classes} lớp khai báo trong mã.
  \item \textbf{Ràng buộc vị từ}: \textsf{TUONG\_TAC\_NGUY\_HIEM} có miền \textsf{Drug} và miền giá trị
        \textsf{RiskSubstance}; \textsf{KHUYEN\_NGHI\_AN} nối \textsf{Disease} với \textsf{Food}.
  \item \textbf{Tiên đề tính chất}: vị từ có thể đối xứng, bắc cầu hoặc là nghịch đảo của vị từ khác
        (Hình~\ref{fig:tiende}).
\end{itemize}

\hinhtikz{t2_ontology}{Cây phân cấp lớp (\texttt{rdfs:subClassOf}) của bản thể học y tế: đúng
\SL{kg/health/classes} lớp khai báo trong \texttt{health\_ontology()}}{ontology}

\hinhtikz{t2_axioms}{Ba tiên đề tính chất của vị từ: bắc cầu, đối xứng, nghịch đảo; cạnh suy ra vẽ nét đứt}{tiende}

\subsection{Suy diễn tiến trên RDF}
Bản thể học không chỉ để kiểm tra mà còn để \textbf{suy ra} tri thức không được ghi tường minh. Hàm
\texttt{rdfs\_closure} cài suy diễn tiến (forward chaining): áp lặp lại năm luật cho tới khi không sinh thêm bộ ba
mới nào (điểm bất động). Với $\sqsubseteq$ là \texttt{rdfs:subClassOf} và \texttt{a} là \texttt{rdf:type}:
\begin{align}
  &(A \sqsubseteq B) \wedge (B \sqsubseteq C) \;\Rightarrow\; A \sqsubseteq C &&\text{(phân cấp bắc cầu)}\\
  &(x\ \texttt{a}\ A) \wedge (A \sqsubseteq B) \;\Rightarrow\; x\ \texttt{a}\ B &&\text{(kế thừa kiểu)}\\
  &p \text{ đối xứng} \wedge (x\, p\, y) \;\Rightarrow\; (y\, p\, x) &&\text{(đối xứng)}\\
  &p \text{ bắc cầu} \wedge (x\, p\, y) \wedge (y\, p\, z) \;\Rightarrow\; (x\, p\, z) &&\text{(bắc cầu)}\\
  &p \equiv q^{-1} \wedge (x\, p\, y) \;\Rightarrow\; (y\, q\, x) &&\text{(nghịch đảo)}
\end{align}
Năm luật này là tập con của ngữ nghĩa RDFS và OWL 2 RL đủ cho hai KG của bài. Luật \texttt{rdfs:domain} và
\texttt{rdfs:range} (suy ra kiểu của chủ thể từ vị từ) cố ý không bật: ở đây domain và range được dùng như
\emph{ràng buộc} để chặn dữ liệu sai, không phải để suy ra kiểu mới.

\maDai{kg_base__rdfs_closure}{Suy diễn tiến năm luật RDFS/OWL tới điểm bất động}

Trên đồ thị y tế, suy diễn thêm \SL{kg/health/inferred} bộ ba vào \SL{kg/health/rdf_triples} bộ ba gốc; trên
đồ thị tài chính thêm \SL{kg/finance/inferred} bộ ba vào \SL{kg/finance/rdf_triples}. Hình~\ref{fig:suydien}
minh họa trên một phần của \texttt{VietnamFinanceKG}: vị từ \textsf{THUOC\_VE} bắc cầu nên từ
``FPT thuộc VN30'' và ``VN30 thuộc HOSE'' suy ra ``FPT thuộc HOSE''.

\hinhtikz{t2_inference}{Suy diễn tiến trên \texttt{VietnamFinanceKG}: cạnh tường minh nét liền, bộ ba suy ra nét
đứt đỏ}{suydien}

\subsection{Truy vấn SPARQL}
SPARQL là ngôn ngữ truy vấn chuẩn của RDF~\cite{w3c2013sparql}: một truy vấn là một \emph{mẫu đồ thị} chứa biến
(\texttt{?thuoc}), và kết quả là mọi cách gán biến làm mẫu khớp với đồ thị. Bài chạy \SL{dem/sparql} truy vấn
trên cả đồ thị gốc lẫn đồ thị đã suy diễn (Bảng~\ref{tab:sparql}). Truy vấn Q1 đọc mức độ tương tác qua khối
reification; Q5 dùng đường thuộc tính \texttt{(p|q)+} để tìm mọi biến chịu tác động từ lãi suất Fed sau bất kỳ số
bước nào.

\maTruyVan{1}{Q1: thuốc tương tác nguy hiểm với rượu bia, đọc mức độ và hậu quả qua reification}
\maTruyVan{5}{Q5: mọi biến chịu tác động trực tiếp hay gián tiếp của lãi suất Fed (đường thuộc tính)}
\maTruyVan{2}{Q2: món ăn và chỉ số GI, truy vấn theo lớp cha \texttt{ex:Food}}

\bang{tab_sparql}{Kết quả \SL{dem/sparql} truy vấn SPARQL trước và sau suy diễn}{sparql}

Cột ``trước suy diễn'' cho thấy vai trò thực sự của bản thể học. Q2, Q4 và Q6 hỏi theo lớp cha (\textsf{Food},
\textsf{Drug}, \textsf{Stock}) trong khi dữ liệu chỉ ghi lớp lá (\textsf{ModerateFood}, \textsf{Biguanide}\dots),
nên trên đồ thị gốc chúng trả về rỗng; sau suy diễn kế thừa kiểu chúng trả về đầy đủ. Không có suy diễn, người
viết truy vấn phải tự liệt kê mọi lớp con, và truy vấn hỏng mỗi khi bản thể học thêm một lớp mới.
Hình~\ref{fig:kg_inference} tổng hợp hai hiệu ứng: số bộ ba tăng thêm và số dòng kết quả của từng truy vấn.

\hinhpdf{kg_inference}{1.0}{Ảnh hưởng của suy diễn: số bộ ba RDF trước, sau suy diễn (trái) và số dòng kết quả
của từng truy vấn SPARQL (phải)}{kg_inference}

\subsection{Kiểm tra ràng buộc}\label{sec:rangbuoc}
Để chứng minh ràng buộc domain/range thực sự hoạt động, kịch bản \texttt{build\_kgs.py} cố tình thêm
\SL{dem/violations} bộ ba sai. Cả bốn đều bị từ chối với thông báo nêu đúng lớp mong đợi và lớp thực tế
(Bảng~\ref{tab:vipham}). Đây là lớp phòng thủ đầu tiên chống lại dữ liệu bẩn, đặt ngay ở thời điểm nạp, trước
khi một bộ ba sai kịp lan vào câu trả lời của tác tử.

\bang{tab_violations}{Bốn bộ ba cố tình sai và thông báo của bộ kiểm tra ràng buộc}{vipham}

\section{Xây KB và KG từ dữ liệu bảng}
Không phải tri thức nào cũng có sẵn dạng văn bản. Bộ dữ liệu chẩn đoán đái tháo đường~\cite{mustafa2023diabetes}
là một bảng CSV \SL{kg/tabular/diabetes/records} dòng; dữ liệu giá cổ phiếu cũng là bảng. Tiếp cận của bài
tham khảo dungvt, đưa một bảng thành hai dạng tri thức song song (Hình~\ref{fig:kbpipeline}):

\begin{itemize}
  \item \textbf{KB văn bản}: mỗi dòng thành một câu mô tả, ví dụ ``Hồ sơ Person\_0: giới tính nữ, tuổi 80, tăng
        huyết áp không, bệnh tim có\dots''. Dạng này đưa vào RAG được.
  \item \textbf{KG bộ ba}: mỗi ô thành một bộ ba $(\textsf{Person\_0}, \textsf{has\_age}, 80)$. Dạng này truy vấn
        chính xác theo điều kiện số được.
\end{itemize}

\hinhtikz{t2_kb_pipeline}{Xây KB và KG từ dữ liệu bảng: mỗi dòng CSV được văn bản hóa thành một tài liệu (KB) và
tách thành các bộ ba (KG), phục vụ hai kiểu truy vấn khác nhau}{kbpipeline}

\maNguon{kg_tabular__diabetes_documents}{Văn bản hóa mỗi hồ sơ thành một câu tiếng Việt}
\maNguon{kg_tabular__diabetes_triples}{Tách mỗi hồ sơ thành các bộ ba thuộc tính}

\bang{tab_tabular}{Chuyển bảng chẩn đoán đái tháo đường thành KB văn bản và KG}{tabular}

Hai dạng tri thức trả lời hai loại câu hỏi khác nhau. Câu hỏi ``có bao nhiêu người trên 70 tuổi, tăng huyết áp,
HbA1c từ 8,8 trở lên'' là một phép lọc trên thuộc tính số: KG lọc ba điều kiện, trả đúng
\SL{kg/tabular/numeric_query/kg_hits} trên \SL{kg/tabular/numeric_query/truth} hồ sơ trong
\SL{kg/tabular/numeric_query/kg_ms} ms. Tìm kiếm văn bản bằng TF-IDF trên 20.000 tài liệu mẫu chỉ đưa được
\SL{kg/tabular/numeric_query/tfidf_top20_correct} hồ sơ đúng vào top-20 (Hình~\ref{fig:kg_tabular}), vì độ tương
đồng từ vựng không hiểu phép so sánh ``$\ge 8{,}8$''. Bài học: RAG mạnh với câu hỏi diễn giải, KG mạnh với câu
hỏi có điều kiện và số liệu; một hệ thống tốt cần cả hai.

\hinhpdf{kg_tabular}{1.0}{Câu hỏi có điều kiện số trên dữ liệu bảng: truy vấn KG trả đúng toàn bộ, tìm văn bản
TF-IDF gần như bỏ sót}{kg_tabular}

\section{So sánh chi tiết: Graph Database và Vector Database}
Hai loại kho dữ liệu chuyên dụng cho hệ thống dùng LLM là cơ sở dữ liệu đồ thị (Neo4j, GraphDB, Amazon
Neptune) và cơ sở dữ liệu vector (FAISS, Milvus, Qdrant, Chroma). Chúng trả lời hai câu hỏi khác nhau
(Hình~\ref{fig:graphvector}): Graph DB trả lời \emph{``cái gì nối với cái gì, bằng quan hệ gì''}; Vector DB trả
lời \emph{``cái gì giống cái này nhất''}.

\hinhtikz{t2_graph_vector}{Graph DB duyệt quan hệ tường minh giữa các thực thể; Vector DB tìm láng giềng gần nhất
theo khoảng cách trong không gian nhúng nhiều chiều}{graphvector}

\begin{table}[H]\centering
\caption{Bảng đối chuẩn kỹ thuật Graph Database và Vector Database}\label{tab:graphvector}
{\small\begin{tabular}{L{3.2cm}L{5.4cm}L{5.4cm}}
\toprule
Tiêu chí & Graph Database & Vector Database \\
\midrule
Đơn vị dữ liệu & nút, cạnh có nhãn, thuộc tính & vector dày $\mathbb{R}^d$ kèm metadata \\
Mô hình truy vấn & khớp mẫu đồ thị (Cypher, SPARQL), duyệt đường & láng giềng gần nhất (k-NN, ANN) \\
Chỉ mục & danh sách kề, chỉ mục theo nhãn & HNSW, IVF, lượng tử hóa PQ \\
Câu hỏi đa bước & tự nhiên: mỗi bước là một bước duyệt & không có khái niệm bước \\
Câu hỏi ngữ nghĩa mờ & kém: cần đúng tên thực thể & mạnh: hiểu diễn đạt khác nhau \\
Điều kiện số, phép đếm & chính xác & không hỗ trợ trực tiếp \\
Khả năng giải thích & đường đi và nguồn của từng cạnh & điểm tương đồng, khó giải thích \\
Chi phí xây dựng & cao: cần trích xuất, chuẩn hóa thực thể & thấp: nhúng văn bản thô \\
Cập nhật & thêm, sửa bộ ba cục bộ & nhúng lại đoạn văn đổi \\
Vai trò trong bài & \texttt{VietnamDiabetesKG}, \texttt{VietnamFinanceKG} & chỉ mục e5 trên \SL{retrieval_grid/n_docs} tài liệu \\
\bottomrule
\end{tabular}}
\end{table}

Hai loại kho không cạnh tranh mà bổ sung: Vector DB tìm \emph{đoạn văn} giải thích, Graph DB cung cấp
\emph{quan hệ} chính xác và đường suy luận. GraphRAG ở Chương~\ref{chap:rag} và tác tử ở Chương~\ref{chap:agent}
dùng cả hai cùng lúc.

% =====================================================================
%  CHƯƠNG 3 — CÁC THẾ HỆ KIẾN TRÚC RAG
% =====================================================================
\chapter[Các thế hệ kiến trúc Retrieval-Augmented Generation]{Các thế hệ kiến trúc\\
Retrieval-Augmented Generation}\label{chap:rag}

RAG không phải một kiến trúc mà là một họ kiến trúc tiến hóa qua ba thế hệ~\cite{gao2023survey}
(Hình~\ref{fig:thehe}). \textbf{Naive RAG} chia văn bản, nhúng, tìm top-$k$ rồi sinh. \textbf{Advanced RAG} thêm
các bước tối ưu trước truy xuất (viết lại câu hỏi) và sau truy xuất (xếp lại, nén ngữ cảnh). \textbf{GraphRAG}
thay đơn vị tri thức từ đoạn văn rời rạc sang đồ thị thực thể và cộng đồng, để trả lời được câu hỏi đa bước và
câu hỏi tổng hợp. Chương này đi qua cả ba thế hệ, và với mỗi thế hệ đo trên chính kho tri thức của bài xem nó cải
thiện được gì, tốn thêm bao nhiêu.

\hinhtikz{t3_generations}{Ba thế hệ RAG xếp tầng: Naive RAG, Advanced RAG với tối ưu tiền và hậu truy xuất,
GraphRAG với đồ thị thực thể và tóm tắt cộng đồng}{thehe}

\section{Naive RAG và các điểm nghẽn}
\subsection{Quy trình hoạt động của Naive RAG}
Naive RAG có ba pha (Hình~\ref{fig:naive}). Pha \textbf{lập chỉ mục} chạy ngoại tuyến: tài liệu được chia thành
đoạn (chunk), mỗi đoạn được mã hóa thành vector $\mathbf{d}_j = f(\text{đoạn}_j) \in \mathbb{R}^{768}$ và lưu vào
chỉ mục. Pha \textbf{truy xuất} chạy khi có câu hỏi: câu hỏi được mã hóa thành $\mathbf{q} = f(\text{câu hỏi})$ và
$k$ đoạn có độ tương đồng cosin lớn nhất được chọn:
\begin{equation}
  \operatorname{sim}(\mathbf{q}, \mathbf{d}_j) = \frac{\mathbf{q}^\top \mathbf{d}_j}{\lVert\mathbf{q}\rVert\,\lVert\mathbf{d}_j\rVert},
  \qquad \mathcal{Z}_k(\mathbf{q}) = \operatorname*{arg\,top-}k_{j}\ \operatorname{sim}(\mathbf{q}, \mathbf{d}_j).
\end{equation}
Pha \textbf{tạo sinh} nối $k$ đoạn vào prompt theo khuôn ``BẰNG CHỨNG: [1]\dots[k] — CÂU HỎI — TRẢ LỜI'' và để LLM
sinh câu trả lời.

\hinhtikz{t3_naive}{Naive RAG ba pha: lập chỉ mục ngoại tuyến, truy xuất và tạo sinh trực tuyến, kèm ba điểm
nghẽn chính}{naive}

\subsection{Các điểm nghẽn của Naive RAG}
Mỗi pha có một điểm nghẽn, và lỗi ở pha trước không pha sau nào sửa được:

\begin{enumerate}
  \item \textbf{Chia đoạn cắt rời ngữ cảnh.} Cửa sổ 400 ký tự có thể cắt giữa câu ``chống chỉ định khi eGFR
        $< 30$''. Đoạn chứa nửa đầu không còn con số, đoạn chứa nửa sau không còn tên thuốc. Mục~\ref{sec:chunking}
        đo hiện tượng này.
  \item \textbf{Truy xuất trượt từ khóa hoặc trượt ngữ nghĩa.} Truy xuất dày hiểu ``tiểu đường'' gần nghĩa ``đái
        tháo đường'' nhưng dễ bỏ qua một mã cụ thể như ``Nghị định 232''; truy xuất thưa thì ngược lại. Với cách
        chia đoạn cố định và truy xuất dense thuần, chỉ \SL{retrieval_grid/grid/fixed/dense/hit@1Pct}\% câu hỏi có đoạn đúng ở vị trí
        đầu.
  \item \textbf{Câu hỏi đa bước không nằm trong đoạn nào.} ``Fed tăng lãi suất thì vàng SJC ra sao'' cần nối ít
        nhất ba mẩu tri thức nằm ở ba tài liệu khác nhau; top-$k$ theo độ tương đồng với câu hỏi không bảo đảm lấy
        đủ cả ba.
\end{enumerate}

\section{Thiết lập thực nghiệm truy xuất}\label{sec:thietlap}
\subsection{Kho tri thức và bộ câu hỏi}
Kho tri thức gồm \SL{retrieval_grid/n_docs} tài liệu Markdown, tám tài liệu y tế và chín tài liệu tài chính
(Bảng~\ref{tab:corpus}). Mỗi tài liệu tổng hợp từ một văn bản gốc có số hiệu, viết lại bằng lời của tác giả và có
cấu trúc tiêu đề hai cấp, mô phỏng cách một phác đồ hay một báo cáo phân tích được tổ chức trong thực tế.

\bang{tab_corpus}{Kho tri thức \SL{retrieval_grid/n_docs} tài liệu của tiểu luận}{corpus}

Bộ đánh giá có \SL{retrieval_grid/n_questions} câu hỏi (\SL{dem/qa_health} y tế, \SL{dem/qa_finance} tài chính).
Mỗi câu hỏi đi kèm một \emph{chuỗi đáp án} nằm nguyên văn trong kho (Bảng~\ref{tab:qa}); câu hỏi được viết bằng
cách diễn đạt khác với văn bản (``tiểu đường'' thay cho ``đái tháo đường'', ``uống bia'' thay cho ``rượu bia'') để
không chấm điểm sự trùng từ đơn thuần.

\bang{tab_qa_sample}{Trích bộ câu hỏi đánh giá truy xuất}{qa}

\subsection{Độ đo}
Một đoạn được coi là \textbf{liên quan} nếu nó chứa trọn chuỗi đáp án; định nghĩa này không cần gán nhãn tay và
công bằng giữa các chiến lược chia đoạn, vì đoạn bị cắt ngang đáp án không được tính. Với $r$ là thứ hạng của đoạn
liên quan đầu tiên trong danh sách trả về, ba độ đo được dùng (Hình~\ref{fig:dodo}):
\begin{align}
  \text{Hit@}k &= \mathbb{1}[r \le k], \qquad \text{MRR} = \frac{1}{|Q|}\sum_{q \in Q} \frac{1}{r_q},\\
  \text{nDCG@}10 &= \frac{1}{\text{IDCG@}10}\sum_{i=1}^{10}\frac{\text{rel}_i}{\log_2(i+1)},
\end{align}
trong đó $\text{rel}_i \in \{0,1\}$ cho biết đoạn ở vị trí $i$ có liên quan không và IDCG là DCG của thứ tự lý
tưởng. Hit@1 trả lời ``đoạn đầu tiên có đúng không'', câu hỏi quan trọng nhất khi chỉ đưa ít đoạn vào LLM; MRR
phạt nhẹ khi đoạn đúng ở vị trí 2--3; nDCG thưởng việc đặt \emph{mọi} đoạn liên quan lên trên.

\hinhtikz{t3_metrics}{Cách tính Hit@k, MRR và nDCG@10 trên một danh sách xếp hạng}{dodo}

\maNguon{retrieval__rank_metrics}{Tính Hit@k, MRR và nDCG@10 cho một câu hỏi}

\section{Các chiến lược phân mảnh (chunking)}\label{sec:chunking}
\subsection{Bài toán đánh đổi: độ dài đoạn và khoảng chồng lấn}
Chọn cỡ đoạn là chọn giữa hai loại lỗi. Đoạn \emph{ngắn} chính xác nhưng dễ cắt rời đáp án khỏi ngữ cảnh của nó;
đoạn \emph{dài} giữ trọn ngữ cảnh nhưng vector của nó là trung bình của nhiều ý, nên khó khớp với một câu hỏi cụ
thể, và mỗi đoạn đưa vào LLM tốn nhiều token hơn. Chồng lấn (overlap) $o$ giữa hai đoạn kề giảm xác suất cắt
ngang câu: với cửa sổ $w$, bước trượt là $w-o$, số đoạn của tài liệu dài $L$ là $\lceil (L-o)/(w-o) \rceil$.

Thí nghiệm quét bảy cỡ đoạn từ 120 tới 1300 ký tự (chia đệ quy, chồng lấn 15\%) cho thấy cả hai mặt của đánh đổi
(Hình~\ref{fig:ret_sweep}, Bảng~\ref{tab:sweep}). Ở cỡ 120 ký tự, chỉ
\SL{retrieval_sweep/0/intactPct}\% câu hỏi còn đoạn chứa trọn đáp án, và dense Hit@1 chỉ đạt
\SL{retrieval_sweep/0/dense_hit@1Pct}\%. Tăng cỡ đoạn thì Hit@1 tăng, nhưng ngữ cảnh top-5 phình từ
\SL{retrieval_sweep/0/ctx_tokens_top5} lên \SL{retrieval_sweep/6/ctx_tokens_top5} token. Lưu ý một hiệu ứng
của cách chấm: ở cỡ 1300 ký tự kho chỉ còn \SL{retrieval_sweep/6/n_chunks} đoạn, gần với truy xuất cả tài liệu,
nên tìm đúng ``đoạn'' dễ hơn nhiều; Hit@1 cao ở đó không có nghĩa LLM nhận được ngữ cảnh gọn. Vùng 400--900 ký
tự là điểm cân bằng, và là cỡ được chọn cho các thí nghiệm sau.

\hinhpdf{ret_sweep}{1.0}{Quét cỡ đoạn: tỉ lệ trúng tăng theo cỡ đoạn trong khi số ký tự ngữ cảnh top-5 tăng tuyến
tính}{ret_sweep}

\bang{tab_sweep}{Ảnh hưởng của cỡ đoạn tới truy xuất BM25 và dense}{sweep}

\subsection{Năm chiến lược phân mảnh}
Năm chiến lược được cài trong \texttt{tl2/chunking.py}, mọi chiến lược đều giữ vị trí ký tự gốc của đoạn để chấm
điểm được (Hình~\ref{fig:chunking}):

\begin{enumerate}
  \item \textbf{Cố định (fixed-size)}: cửa sổ trượt 400 ký tự, chồng lấn 60, không quan tâm ranh giới câu.
  \item \textbf{Đệ quy (recursive)}: thử cắt ở dấu phân cách cấp cao nhất (đoạn trống, xuống dòng, dấu chấm,
        chấm phẩy, dấu phẩy, khoảng trắng) và chỉ xuống cấp thấp hơn khi mảnh vẫn quá dài; các mảnh nhỏ liền nhau
        được gộp tới giới hạn.
  \item \textbf{Theo cấu trúc (structure-aware)}: mỗi mục Markdown là một đoạn, gắn chuỗi tiêu đề làm ngữ cảnh
        (``Thuốc điều trị > Metformin > Chống chỉ định''); mục quá 900 ký tự mới tách đệ quy bên trong.
  \item \textbf{Ngữ nghĩa (semantic)}: nhúng từng câu, cắt khi khoảng cách cosin giữa hai câu liền nhau vượt phân
        vị thứ 80 của chuỗi khoảng cách trong tài liệu.
  \item \textbf{Cha -- con (parent-child, small-to-big)}: lập chỉ mục các đoạn con khoảng 220 ký tự để khớp chính
        xác, nhưng trả cho LLM đoạn cha là cả mục (Hình~\ref{fig:chacon}).
\end{enumerate}

\hinhtikz{t3_chunking}{Năm chiến lược chia đoạn trên cùng một tài liệu: dải ngang là tài liệu, mỗi khối màu là
một đoạn}{chunking}

\hinhtikz{t3_parent_child}{Truy xuất cha -- con: khớp trên đoạn con nhỏ, đưa đoạn cha đủ ngữ cảnh cho LLM}{chacon}

\maNguon{chunking___recursive_spans}{Lõi của chia đoạn đệ quy: cắt ở dấu phân cách cấp cao nhất có thể}
\maNguon{chunking__semantic}{Chia đoạn ngữ nghĩa theo phân vị khoảng cách cosin giữa các câu liền nhau}
\maNguon{chunking__parent_child}{Chia đoạn cha -- con: cha là mục, con là mảnh đệ quy nhỏ bên trong}

Bảng~\ref{tab:chunkstats} và Hình~\ref{fig:ret_chunk_lengths} cho thấy hình dạng các đoạn sinh ra. Cố định cho độ
dài gần như đều; theo cấu trúc dao động theo độ dài mục; ngữ nghĩa dao động mạnh nhất, có đoạn chỉ
\SL{retrieval_grid/chunk_stats/semantic/min_len} ký tự (một dòng tiêu đề bị tách riêng) và có đoạn tới
\SL{retrieval_grid/chunk_stats/semantic/max_len} ký tự. Chia cố định là chiến lược duy nhất làm mất đáp án: chỉ
\SL{retrieval_grid/chunk_stats/fixed/intactPct}\% câu hỏi còn đoạn chứa trọn chuỗi đáp án.

\bang{tab_chunk_stats}{Thống kê đoạn của năm chiến lược trên toàn kho}{chunkstats}

\hinhpdf{ret_chunk_lengths}{1.0}{Phân bố độ dài đoạn của năm chiến lược}{ret_chunk_lengths}

Hình~\ref{fig:ret_semantic} mở hộp chiến lược ngữ nghĩa trên tài liệu \texttt{h03-ruou-bia}: mỗi cột là khoảng
cách giữa hai câu liền nhau, đường ngang là ngưỡng phân vị 80, cột đỏ là điểm cắt. Các điểm cắt rơi đúng vào chỗ
văn bản chuyển chủ đề (từ cơ chế nhiễm toan lactic sang nguy cơ hạ glucose huyết), nhưng cũng có điểm cắt giữa
một danh sách gạch đầu dòng vì các mục của danh sách ngắn và khác từ vựng.

\hinhpdf{ret_semantic}{1.0}{Chia đoạn ngữ nghĩa: khoảng cách cosin giữa các câu liền nhau và điểm cắt}{ret_semantic}

\subsection{Ý nghĩa của chunking trong dữ liệu y tế và tài chính Việt Nam}
Văn bản y tế và tài chính có hai đặc điểm làm chunking quan trọng hơn văn bản thông thường. Thứ nhất, \emph{con
số gắn chặt với ngữ cảnh}: ``$\ge 6{,}5\%$'' chỉ có nghĩa khi đi cùng ``HbA1c'' và ``tiêu chuẩn chẩn đoán''; một
đoạn chứa con số nhưng mất chủ ngữ là một đoạn nguy hiểm. Thứ hai, \emph{tài liệu có cấu trúc chuẩn}: phác đồ Bộ
Y tế, nghị định, báo cáo thường niên đều chia mục rõ ràng. Chiến lược theo cấu trúc khai thác đúng đặc điểm thứ
hai để bảo vệ đặc điểm thứ nhất: tiêu đề mục được gắn vào đoạn, nên đoạn ``Chống chỉ định'' luôn biết nó thuộc
mục ``Metformin''. Tiếng Việt còn có thêm một khó khăn: từ được ghép từ nhiều âm tiết cách nhau bởi khoảng trắng,
nên cắt theo khoảng trắng có thể tách ``đái tháo'' khỏi ``đường''; chia đệ quy ưu tiên dấu câu trước khoảng trắng
để tránh điều này.

\section{Truy xuất: thưa, dày, lai và xếp lại}
\subsection{Truy xuất thưa: TF-IDF và BM25}
Truy xuất thưa biểu diễn văn bản bằng tần suất từ. BM25~\cite{robertson2009bm25} cho điểm tài liệu $d$ với câu hỏi
$q$ gồm các từ $t$ bằng
\begin{equation}
  \operatorname{BM25}(q,d) = \sum_{t \in q} \operatorname{IDF}(t)\,
  \frac{f(t,d)\,(k_1+1)}{f(t,d) + k_1\left(1 - b + b\,\frac{|d|}{\operatorname{avgdl}}\right)},
\end{equation}
với $f(t,d)$ là tần suất từ, $|d|$ là độ dài đoạn, $k_1 = 1{,}5$ điều tiết bão hòa tần suất và $b = 0{,}75$ chuẩn
hóa độ dài. Tách từ ở mức âm tiết viết thường, giữ số và đơn vị như ``mg'', ``dl''.

\maNguon{retrieval__BM25Retriever}{Bộ truy xuất BM25 trên các âm tiết tiếng Việt}

\subsection{Truy xuất dày: mô hình nhúng đa ngôn ngữ}
Truy xuất dày~\cite{karpukhin2020dpr} dùng một bộ mã hóa Transformer ánh xạ câu hỏi và đoạn văn vào cùng một không
gian vector. Bài dùng \texttt{multilingual-e5-base}~\cite{wang2024me5}, mô hình được huấn luyện tương phản trên dữ
liệu đa ngôn ngữ, yêu cầu tiền tố \texttt{query:} cho câu hỏi và \texttt{passage:} cho đoạn văn. Đây là kiến trúc
\textbf{bi-encoder}~\cite{reimers2019sbert}: câu hỏi và đoạn văn được mã hóa \emph{riêng}, nên vector đoạn văn tính trước một lần khi lập
chỉ mục (Hình~\ref{fig:bicross}).

\maNguon{retrieval__Embedder}{Bọc SentenceTransformer cho mô hình e5: tiền tố, chuẩn hóa L2, chạy GPU}

So sánh ba mô hình nhúng (Bảng~\ref{tab:embed}, Hình~\ref{fig:ret_embedders}) cho thấy họ e5 vượt xa
\texttt{paraphrase-multilingual-MiniLM}: Hit@1 \SL{retrieval_embed/e5-base/hit@1Pct}\% so với
\SL{retrieval_embed/minilm/hit@1Pct}\%. Bản e5-small và e5-base có Hit@1 bằng nhau trên bộ câu hỏi này nhưng phân
bổ khác theo miền: e5-base tốt hơn ở y tế, e5-small tốt hơn ở tài chính. Bài chọn e5-base vì y tế là miền đòi hỏi
độ chính xác cao hơn.

\bang{tab_embed}{So sánh ba mô hình nhúng đa ngôn ngữ (truy xuất dense, đoạn theo cấu trúc)}{embed}

\hinhpdf{ret_embedders}{1.0}{Ba mô hình nhúng đa ngôn ngữ trên cùng bộ câu hỏi}{ret_embedders}

Hình~\ref{fig:ret_pca} chiếu \SL{retrieval_grid/chunk_stats/structure/n_chunks} đoạn theo cấu trúc xuống hai
chiều bằng PCA: đoạn y tế và đoạn tài chính tách thành hai vùng rõ rệt, và trong mỗi vùng các đoạn cùng tài liệu
nằm gần nhau. Hai thành phần đầu chỉ giữ \SL{retrieval_viz/embedding_pca/var_ratio/0Pct}\% và
\SL{retrieval_viz/embedding_pca/var_ratio/1Pct}\% phương sai, nhắc rằng không gian 768 chiều phức tạp hơn nhiều so
với những gì một hình hai chiều thể hiện. Hình~\ref{fig:ret_simmatrix} là ma trận cosin giữa mười câu hỏi và các
đoạn: ô sáng nhất mỗi hàng thường là đoạn đúng, nhưng khoảng cách giữa ô đúng và các ô kế bên nhỏ, lý do cần thêm
bước xếp lại.

\hinhpdf{ret_pca}{0.92}{Các đoạn trong không gian nhúng e5-base chiếu PCA hai chiều}{ret_pca}
\hinhpdf{ret_simmatrix}{1.0}{Độ tương đồng cosin giữa câu hỏi và đoạn văn}{ret_simmatrix}

\subsection{Truy xuất lai bằng Reciprocal Rank Fusion}
BM25 và dense mắc lỗi ở những câu hỏi khác nhau, nên gộp chúng lại có lợi. Điểm số của hai bộ truy xuất không
cùng thang đo (BM25 không giới hạn, cosin trong $[-1,1]$), nên thay vì cộng điểm, Reciprocal Rank Fusion
(RRF)~\cite{cormack2009rrf} cộng nghịch đảo thứ hạng:
\begin{equation}
  \operatorname{RRF}(d) = \sum_{r \in \{\text{BM25},\ \text{dense}\}} \frac{1}{k + \operatorname{rank}_r(d)}, \qquad k = 60.
\end{equation}
Hằng số $k$ làm phẳng đóng góp của các vị trí đầu, để một đoạn đứng hạng 1 ở một bộ truy xuất nhưng hạng 50 ở bộ
kia không áp đảo một đoạn đứng hạng 3 ở cả hai.

\maNguon{retrieval__HybridRRF}{Truy xuất lai: hợp nhất thứ hạng BM25 và dense bằng RRF}

\subsection{Xếp lại bằng cross-encoder}
Bi-encoder nhanh nhưng không bao giờ nhìn câu hỏi và đoạn văn \emph{cùng lúc}. Cross-encoder~\cite{nogueira2019rerank}
nối hai chuỗi thành một đầu vào và cho Transformer chú ý chéo giữa chúng, đầu ra là một điểm liên quan trực tiếp
$s(q,d) = \sigma\bigl(\mathbf{w}^\top \mathrm{Enc}([q; d])_{\texttt{[CLS]}}\bigr)$ (Hình~\ref{fig:bicross}). Cái
giá là phải chạy mô hình một lần cho \emph{mỗi} cặp, nên chỉ dùng được ở giai đoạn hai trên một tập ứng viên nhỏ.
Bài dùng \texttt{bge-reranker-v2-m3}~\cite{chen2024bgem3}, xếp lại 30 ứng viên của truy xuất lai
(Hình~\ref{fig:hybrid}).

\hinhtikz{t3_bi_cross}{Bi-encoder mã hóa riêng câu hỏi và đoạn văn rồi so cosin; cross-encoder mã hóa chung cả cặp
và cho điểm trực tiếp}{bicross}

\hinhtikz{t3_hybrid}{Phễu truy xuất lai: BM25 và dense chạy song song, hợp nhất RRF, cross-encoder xếp lại 30 ứng
viên và giữ top-$k$}{hybrid}

\maNguon{retrieval__Reranker}{Cross-encoder bge-reranker-v2-m3 chạy bf16 trên GPU}

\subsection{Kết quả: chiến lược chia đoạn $\times$ bộ truy xuất}
Bảng~\ref{tab:gridhit1} và Hình~\ref{fig:ret_grid_hit1} là kết quả chính của chương: 25 tổ hợp chia đoạn và truy
xuất trên \SL{retrieval_grid/n_questions} câu hỏi. Tổ hợp tốt nhất là \textbf{chia theo cấu trúc + truy xuất lai +
xếp lại}, Hit@1 \SL{retrieval_grid/grid/structure/hybrid_rerank/hit@1Pct}\% và MRR
\SL{retrieval_grid/grid/structure/hybrid_rerank/mrr}; tổ hợp Naive RAG (cố định + dense) chỉ đạt
\SL{retrieval_grid/grid/fixed/dense/hit@1Pct}\%.

\bang{tab_grid_hit1}{Hit@1 của 25 tổ hợp chiến lược chia đoạn và bộ truy xuất}{gridhit1}
\hinhpdf{ret_grid_hit1}{1.0}{Bản đồ nhiệt Hit@1: chiến lược chia đoạn (hàng) $\times$ bộ truy xuất (cột)}{ret_grid_hit1}
\bang{tab_grid_mrr}{MRR của 25 tổ hợp chiến lược chia đoạn và bộ truy xuất}{gridmrr}

Ba quan sát rút ra từ lưới:
\begin{itemize}
  \item \textbf{Chiến lược chia đoạn quan trọng ngang với bộ truy xuất.} Cùng bộ hybrid + rerank, đổi từ chia cố
        định sang chia theo cấu trúc nâng Hit@1 từ \SL{retrieval_grid/grid/fixed/hybrid_rerank/hit@1Pct}\% lên
        \SL{retrieval_grid/grid/structure/hybrid_rerank/hit@1Pct}\%.
  \item \textbf{Mỗi tầng của phễu đều có ích.} Với chia theo cấu trúc: BM25
        \SL{retrieval_grid/grid/structure/bm25/hit@1Pct}\%, dense \SL{retrieval_grid/grid/structure/dense/hit@1Pct}\%,
        lai \SL{retrieval_grid/grid/structure/hybrid/hit@1Pct}\%, lai + xếp lại
        \SL{retrieval_grid/grid/structure/hybrid_rerank/hit@1Pct}\%.
  \item \textbf{Cha -- con mạnh nhất khi không có reranker.} Khớp trên đoạn con nhỏ cho dense Hit@1
        \SL{retrieval_grid/grid/parent_child/dense/hit@1Pct}\%, cao nhất cột dense; khi đã có reranker thì ưu thế
        này thu hẹp.
\end{itemize}

Bảng~\ref{tab:retdetail} đi sâu vào chiến lược theo cấu trúc. Hình~\ref{fig:ret_hitk} là đường Hit@$k$: mọi bộ
truy xuất đều tìm được đoạn đúng trong top-10 ở gần như mọi câu hỏi; khác biệt nằm ở \emph{thứ hạng}, đúng điều
mà reranker cải thiện. Hình~\ref{fig:ret_domain} tách theo miền: câu hỏi y tế khó hơn tài chính với mọi bộ truy
xuất, vì nhiều tài liệu y tế cùng nói về một bệnh và dùng chung từ vựng (metformin xuất hiện trong bốn tài liệu).

\bang{tab_ret_detail}{Chi tiết năm bộ truy xuất trên đoạn chia theo cấu trúc}{retdetail}
\hinhpdf{ret_hitk}{0.95}{Đường Hit@$k$ của năm bộ truy xuất}{ret_hitk}
\hinhpdf{ret_domain}{0.95}{Hit@1 theo miền y tế và tài chính}{ret_domain}

Chất lượng có giá. Hình~\ref{fig:ret_latency} đặt Hit@1 cạnh độ trễ: BM25 trả lời trong
\SL{retrieval_grid/grid/structure/bm25/latency_ms} ms, dense trong \SL{retrieval_grid/grid/structure/dense/latency_ms}
ms (gồm mã hóa câu hỏi trên GPU), còn hybrid + rerank mất \SL{retrieval_grid/grid/structure/hybrid_rerank/latency_ms}
ms vì phải chạy cross-encoder 30 lần. Với một tác tử hội thoại, nửa giây là chấp nhận được; với tìm kiếm thời gian
thực trên hàng triệu đoạn thì cần giảm số ứng viên hoặc dùng reranker nhỏ hơn.

\hinhpdf{ret_latency}{0.95}{Chất lượng và độ trễ của 25 tổ hợp: xếp lại đắt hơn nhiều nhưng chính xác hơn}{ret_latency}

\section{Advanced RAG: tối ưu đa chặng}
Advanced RAG giữ khung truy xuất -- sinh của Naive RAG nhưng thêm bước tối ưu ở hai đầu:

\begin{itemize}
  \item \textbf{Tiền truy xuất} sửa \emph{câu hỏi}: viết lại câu hỏi cho rõ, tách câu hỏi phức thành câu hỏi con,
        sinh nhiều cách hỏi (Multi-Query), hoặc sinh một câu trả lời giả định rồi dùng nó để tìm (HyDE).
  \item \textbf{Truy xuất} cải tiến bằng chia đoạn tốt hơn, truy xuất lai, lọc theo metadata (miền, tài liệu).
  \item \textbf{Hậu truy xuất} sửa \emph{ngữ cảnh}: xếp lại bằng cross-encoder, nén hoặc cắt bớt đoạn, sắp xếp lại
        để thông tin quan trọng nằm ở đầu và cuối prompt.
\end{itemize}

Phần truy xuất lai và xếp lại đã được đo ở mục trước. Mục này đo hai kỹ thuật tiền truy xuất dùng LLM cục bộ
(Hình~\ref{fig:hyde}).

\textbf{HyDE} (Hypothetical Document Embeddings)~\cite{gao2023hyde} dựa trên nhận xét: câu hỏi và đoạn trả lời
thường không giống nhau về hình thức, nhưng hai đoạn trả lời thì giống nhau. Thay vì nhúng câu hỏi $q$, HyDE cho
LLM viết một đoạn trả lời giả định $\tilde{d} = \mathrm{LLM}(q)$ rồi tìm bằng $f(\tilde{d})$. Biến thể
\emph{HyDE + câu hỏi} hợp nhất bằng RRF hai danh sách tìm bằng $f(q)$ và $f(\tilde d)$.
\textbf{Multi-Query} cho LLM viết ba cách hỏi khác $q_1, q_2, q_3$, truy xuất cho từng câu và hợp nhất bốn danh
sách bằng RRF.

\hinhtikz{t3_hyde}{Hai kỹ thuật tối ưu tiền truy xuất bằng LLM cục bộ: HyDE tìm bằng câu trả lời giả định,
Multi-Query hợp nhất nhiều cách hỏi bằng RRF}{hyde}

\maNguon{run_advanced__rrf}{Hợp nhất nhiều danh sách xếp hạng bằng RRF trong thí nghiệm Advanced RAG}

Kết quả (Bảng~\ref{tab:advanced}, Hình~\ref{fig:adv_rag}) là một kết quả âm có giá trị. Với Qwen2.5-1.5B, HyDE
\emph{làm giảm} Hit@1 từ \SL{advanced/agg/dense/hit@1Pct}\% xuống \SL{advanced/agg/hyde/hit@1Pct}\%, thua dense gốc
ở \SL{advanced/agg/hyde/losses_vs_dense} câu và chỉ thắng ở \SL{advanced/agg/hyde/wins_vs_dense} câu.
Multi-Query giảm ít hơn (\SL{advanced/agg/multiquery/hit@1Pct}\%) nhưng vẫn giảm. Bảng~\ref{tab:hydeex} giải thích
lý do: đoạn giả định của mô hình nhỏ chứa các chi tiết \emph{bịa}, ví dụ nói metformin cùng rượu gây ``tăng huyết
áp'', và vector của chi tiết bịa kéo truy xuất về các đoạn sai. HyDE chỉ có lợi khi LLM sinh đoạn giả định đủ đúng
để trông giống đoạn thật; với một mô hình 1,5 tỉ tham số trên miền chuyên sâu tiếng Việt, điều kiện đó không thỏa.
Mỗi câu hỏi còn tốn thêm \SL{advanced/gen_s/hyde} s sinh văn bản. Vì vậy tác tử ở Chương~\ref{chap:agent} không
dùng HyDE hay Multi-Query mà dùng truy xuất lai + xếp lại.

\bang{tab_advanced}{Tối ưu tiền truy xuất bằng LLM so với truy xuất dense gốc}{advanced}
\hinhpdf{adv_rag}{1.0}{Ảnh hưởng của HyDE và Multi-Query: Hit@1, MRR và số câu thắng, thua so với dense gốc}{adv_rag}
\bang{tab_hyde_ex}{Ví dụ văn bản do Qwen2.5-1.5B sinh trong HyDE và Multi-Query}{hydeex}

\section{GraphRAG: tổng hợp tri thức trên đồ thị}
\subsection{Ý tưởng và quy trình}
Advanced RAG vẫn coi tri thức là một tập đoạn văn rời rạc. GraphRAG~\cite{edge2024graphrag} thay đơn vị tri thức
bằng \textbf{đồ thị thực thể}, rồi thêm một tầng trừu tượng là \textbf{cộng đồng}: nhóm thực thể liên kết chặt
được tóm tắt trước thành một ``báo cáo cộng đồng''. Quy trình gồm bốn bước (Hình~\ref{fig:graphrag}):

\begin{enumerate}
  \item \textbf{Đồ thị thực thể}: ở bài báo gốc, LLM trích xuất thực thể và quan hệ từ văn bản; ở đây hai KG đã có
        sẵn và được hợp nhất thành một đồ thị vô hướng có trọng số (cạnh nhiều quan hệ được cộng trọng số).
  \item \textbf{Phát hiện cộng đồng}: thuật toán Louvain~\cite{blondel2008louvain} tìm phân hoạch $\{c_i\}$ cực đại
        độ đo modularity
        \begin{equation}
          Q = \frac{1}{2m}\sum_{i,j}\Bigl(A_{ij} - \frac{k_i k_j}{2m}\Bigr)\delta(c_i, c_j),
        \end{equation}
        với $A_{ij}$ là trọng số cạnh, $k_i$ là bậc có trọng số, $m$ là tổng trọng số. $Q$ đo mật độ cạnh trong cộng
        đồng vượt bao nhiêu so với một đồ thị ngẫu nhiên cùng dãy bậc.
  \item \textbf{Báo cáo cộng đồng}: các bộ ba nằm trọn trong một cộng đồng được diễn đạt thành câu và LLM viết một
        đoạn tóm tắt 3--4 câu.
  \item \textbf{Tìm kiếm}: \emph{cục bộ} liên kết thực thể trong câu hỏi rồi lấy các bộ ba quanh chúng và đoạn
        văn của tài liệu nguồn; \emph{toàn cục} xếp hạng các báo cáo cộng đồng theo độ tương đồng với câu hỏi.
\end{enumerate}

\hinhtikz{t3_graphrag}{Đường ống GraphRAG: đồ thị thực thể, cộng đồng Louvain, báo cáo cộng đồng do LLM viết,
tìm kiếm cục bộ và toàn cục}{graphrag}

Bước khó nhất trong thực tế là \textbf{liên kết thực thể}: nhận ra ``đang uống glucophage'' nói về
\textsf{Metformin}. Bộ liên kết của bài (Đoạn mã~\ref{code:graphrag__EntityLinker}) bỏ dấu tiếng Việt, khớp nguyên
từ với bảng bí danh và ưu tiên bí danh dài nhất trước, để ``gà rán'' không bị hiểu thành ``gà'' và một đoạn đã
khớp không bị khớp lại.

\maNguon{graphrag__EntityLinker}{Liên kết thực thể: khớp bí danh dài nhất trước, không dấu, nguyên từ}
\maNguon{graphrag__summarize_communities}{Tóm tắt cộng đồng bằng LLM cục bộ từ các bộ ba nằm trong cộng đồng}
\maDai{graphrag__GraphRAG}{Lớp \texttt{GraphRAG}: tìm kiếm cục bộ trên KG và tìm kiếm toàn cục trên báo cáo cộng đồng}

\subsection{Cộng đồng tri thức}
Louvain tìm được \SL{graphrag/n_comms} cộng đồng với modularity \SL{graphrag/modularity}
(Hình~\ref{fig:grag_communities}, Bảng~\ref{tab:communities}). Không cộng đồng nào chứa thực thể của cả hai miền,
đúng như mong đợi vì hai KG không có cạnh nối nhau. Các cộng đồng có nghĩa rõ ràng mà không cần ai gán nhãn: C0
gom đái tháo đường típ 2 với tiêu chuẩn chẩn đoán, yếu tố nguy cơ, các món tinh bột và lối sống; C2 gom ba thuốc
chính (metformin, insulin, gliclazide) với rượu bia, nhiễm toan lactic, hạ glucose huyết và các chống chỉ định,
tức đúng phân hệ ``an toàn thuốc''; C6 là bộ ba empagliflozin -- suy tim -- bệnh thận mạn, đúng chỉ định ưu tiên
của nhóm SGLT2. Bên tài chính, C1 là kênh tỷ giá -- lãi suất -- cổ phiếu và ngành, C3 là kênh vàng thế giới (Fed,
DXY, XAU/USD), C4 là chính sách vàng miếng trong nước. LLM viết \SL{dem/comms} bản tóm tắt trong
\SL{graphrag/summary_s} s.

\hinhpdf{grag_communities}{1.0}{Các cộng đồng Louvain trên đồ thị hợp nhất hai KG; nhãn ghi kích thước và thực thể
trung tâm, cạnh đứt nối hai cộng đồng khác nhau}{grag_communities}

\bang{tab_communities}{Các cộng đồng tri thức và báo cáo cộng đồng}{communities}

\subsection{So sánh bốn kiến trúc trên câu hỏi đa bước}
Để đo đúng thế mạnh của GraphRAG, bài dùng một bộ \SL{dem/grag_q} câu hỏi đa bước riêng, mỗi câu kèm các nhóm từ
khóa ứng với các dữ kiện bắt buộc phải có (ví dụ câu Fed -- vàng SJC cần cả kênh DXY, kênh XAU, kênh tỷ giá và
phần bù trong nước). Bốn kiến trúc dùng cùng LLM và cùng prompt có trích dẫn, chỉ khác ngữ cảnh:
Naive (cố định + dense, top-5), Advanced (cấu trúc + lai + xếp lại, top-5), GraphRAG cục bộ (bộ ba quanh thực thể
liên kết + 3 đoạn văn của tài liệu nguồn) và GraphRAG cục bộ + toàn cục (thêm 2 báo cáo cộng đồng gần nhất).
\textbf{Độ phủ ngữ cảnh} là tỉ lệ dữ kiện có mặt trong ngữ cảnh đưa vào LLM; \textbf{độ phủ câu trả lời} là tỉ lệ dữ
kiện có trong câu trả lời cuối.

\bang{tab_graphrag}{Bốn kiến trúc RAG trên \SL{dem/grag_q} câu hỏi đa bước}{graphrag}
\hinhpdf{grag_coverage}{1.0}{Độ phủ dữ kiện trong ngữ cảnh và trong câu trả lời, cùng độ dài ngữ cảnh của bốn kiến
trúc}{grag_coverage}

GraphRAG cục bộ cho câu trả lời đầy đủ nhất: độ phủ câu trả lời \SL{graphrag/agg/graph_local/ans_covPct}\% so với
\SL{graphrag/agg/naive/ans_covPct}\% của Naive và \SL{graphrag/agg/advanced/ans_covPct}\% của Advanced. Điều đáng
chú ý là Advanced có độ phủ \emph{ngữ cảnh} cao nhất (\SL{graphrag/agg/advanced/ctx_covPct}\%) nhưng độ phủ câu
trả lời thấp hơn GraphRAG: dữ kiện có trong năm đoạn văn dài, nhưng LLM nhỏ không gom đủ chúng lại. Bộ ba KG đã
diễn đạt thành câu ngắn, mỗi câu một quan hệ, nên LLM dễ chép đúng vào câu trả lời hơn. Thêm báo cáo toàn cục đưa
độ phủ ngữ cảnh lên \SL{graphrag/agg/graph_local_global/ctx_covPct}\% nhưng \emph{giảm} độ phủ câu trả lời xuống
\SL{graphrag/agg/graph_local_global/ans_covPct}\%, vì ngữ cảnh dài thêm gần gấp rưỡi
(\SL{graphrag/agg/graph_local_global/ctx_chars} ký tự) và mô hình nhỏ bị loãng sự chú ý. Tìm kiếm toàn cục hợp
với câu hỏi tổng hợp kiểu ``chủ đề chính của kho là gì'', không phải câu hỏi đa bước có đáp án cụ thể.

\bang{tab_grag_perq}{Độ phủ ngữ cảnh / câu trả lời theo từng câu hỏi đa bước}{gragperq}

\section{Chiến lược lựa chọn kiến trúc trong hệ thống thực tế}
Không có kiến trúc RAG tốt nhất cho mọi trường hợp. Bảng~\ref{tab:chonkientruc} tóm tắt hướng dẫn chọn rút ra từ
các thí nghiệm của chương.

\begin{table}[H]\centering
\caption{Hướng dẫn chọn kiến trúc RAG theo loại câu hỏi và điều kiện triển khai}\label{tab:chonkientruc}
{\small\begin{tabular}{L{4.4cm}L{4.4cm}L{5.4cm}}
\toprule
Tình huống & Kiến trúc nên dùng & Căn cứ trong bài \\
\midrule
Hỏi đáp tra cứu một dữ kiện & Advanced: cấu trúc + lai + xếp lại & Hit@1 cao nhất trên \SL{retrieval_grid/n_questions} câu \\
Không có GPU, cần độ trễ thấp & BM25 hoặc lai không xếp lại & nhanh hơn reranker hàng trăm lần \\
Câu hỏi đa bước, nhân quả & GraphRAG cục bộ & độ phủ câu trả lời cao nhất \\
Câu hỏi tổng hợp toàn kho & GraphRAG toàn cục & báo cáo cộng đồng tóm tắt sẵn \\
Câu hỏi có điều kiện số & truy vấn KG / SPARQL & TF-IDF bỏ sót gần hết (Mục 2.4) \\
LLM nhỏ ($\le$ 2 tỉ tham số) & tránh HyDE, Multi-Query & cả hai làm giảm Hit@1 \\
\bottomrule
\end{tabular}}
\end{table}

Tác tử của bài kết hợp ba hàng của bảng: truy xuất lai + xếp lại cho đoạn văn, KG cho quan hệ và ràng buộc, tính
toán trực tiếp cho số liệu thị trường, và để bộ lập kế hoạch chọn công cụ theo loại câu hỏi (Chương~\ref{chap:agent}).

% =====================================================================
%  CHƯƠNG 4 — HỆ THỐNG Y TẾ: ĐÁI THÁO ĐƯỜNG TÍP 2 TẠI VIỆT NAM
% =====================================================================
\chapter{Hệ thống y tế: bệnh đái tháo đường típ 2 tại Việt Nam}\label{chap:yte}

Chương này dựng phân hệ y tế của tác tử gồm ba lớp: kho tri thức văn bản về chẩn đoán, điều trị, dinh dưỡng;
\texttt{VietnamDiabetesKG} lưu các ràng buộc lâm sàng có nguồn; và một mạng MLP ước lượng nguy cơ đái tháo đường
từ hồ sơ xét nghiệm. Ba lớp được kiểm tra trên ba tình huống thường gặp ở người bệnh Việt Nam: uống rượu bia
khi đang dùng metformin, ăn thịt gà, và ăn cơm, bún, xôi.

\begin{phathien}{Lưu ý về phạm vi}
Hệ thống được xây cho mục đích học tập. Nội dung y khoa tổng hợp từ Quyết định 5481/QĐ-BYT, tiêu chuẩn của ADA và
các bảng chỉ số đường huyết đã công bố; mọi câu trả lời của tác tử đều kèm lời nhắc ``không thay thế chỉ định của
bác sĩ''.
\end{phathien}

\section{Bối cảnh dịch tễ và phác đồ của Bộ Y tế}
Đái tháo đường đang tăng nhanh tại Việt Nam cùng đô thị hóa, thay đổi chế độ ăn và giảm vận động. Theo Liên đoàn
Đái tháo đường Quốc tế, Việt Nam có khoảng 4 đến 5 triệu người trưởng thành mắc bệnh, tức khoảng 6 đến 7\% dân số
trưởng thành~\cite{idf2021atlas}, và các điều tra quốc gia ghi nhận tới khoảng một nửa số người mắc chưa biết mình
bị bệnh. Đái tháo đường típ 2 chiếm 90 đến 95\% số ca; cơ chế chính là đề kháng insulin ở cơ, gan và mô mỡ kèm suy
giảm dần chức năng tế bào beta của tụy.

Hướng dẫn chẩn đoán và điều trị đái tháo đường típ 2 hiện hành là tài liệu ban hành kèm Quyết định 5481/QĐ-BYT ngày
30/12/2020~\cite{byt5481}, thống nhất với tiêu chuẩn của Hiệp hội Đái tháo đường Hoa Kỳ~\cite{ada2024standards}.
Bảng~\ref{tab:chandoan} ghi bốn tiêu chuẩn chẩn đoán; đây chính là các thực thể loại \textsf{Biomarker} trong KG.

\begin{table}[H]\centering
\caption{Tiêu chuẩn chẩn đoán đái tháo đường và tiền đái tháo đường theo Quyết định 5481/QĐ-BYT}\label{tab:chandoan}
{\small\begin{tabular}{L{6.2cm}L{3.8cm}L{4.2cm}}
\toprule
Xét nghiệm & Đái tháo đường & Tiền đái tháo đường \\
\midrule
Glucose huyết tương lúc đói (nhịn ăn $\ge$ 8 giờ) & $\ge$ 126 mg/dL (7,0 mmol/L) & 100 -- 125 mg/dL \\
Glucose 2 giờ sau nghiệm pháp 75 g glucose & $\ge$ 200 mg/dL (11,1 mmol/L) & 140 -- 199 mg/dL \\
HbA1c (phòng xét nghiệm chuẩn hóa) & $\ge$ 6,5\% (48 mmol/mol) & 5,7 -- 6,4\% \\
Glucose bất kỳ kèm triệu chứng kinh điển & $\ge$ 200 mg/dL & -- \\
\bottomrule
\multicolumn{3}{l}{\footnotesize Không có triệu chứng kinh điển thì cần hai kết quả bất thường, cùng mẫu hoặc hai thời điểm.}
\end{tabular}}
\end{table}

Về điều trị, metformin là thuốc đầu tay; nhóm ức chế SGLT2 (empagliflozin) được ưu tiên khi người bệnh có suy tim
hoặc bệnh thận mạn; insulin bắt buộc với típ 1 và dùng khi típ 2 không kiểm soát được bằng thuốc uống. Mục tiêu
HbA1c cho đa số người trưởng thành là dưới 7\%. Thay đổi lối sống đi cùng mọi bậc điều trị: vận động ít nhất 150
phút mỗi tuần và giảm 5 đến 7\% cân nặng; ở người tiền đái tháo đường, chương trình DPP cho thấy can thiệp lối
sống giảm 58\% nguy cơ tiến triển thành bệnh~\cite{dpp2002}.

Ba đặc điểm của người bệnh Việt Nam khiến tri thức chung chưa đủ. Bữa ăn lấy \emph{cơm trắng} làm nền, chỉ số
đường huyết GI khoảng 73; các món xôi, bánh, bún, phở phổ biến; và \emph{rượu bia} gắn với tiệc tùng, đám cưới,
giỗ chạp, nên câu hỏi ``đang dùng thuốc thì uống được không'' là câu hỏi thật, thường xuyên. KG của bài vì vậy có
một phân hệ dinh dưỡng với các món ăn Việt và một phân hệ tương tác thuốc -- rượu.

\section{Thiết kế bản thể học và đồ thị tri thức y tế}
\subsection{Hệ thống thực thể y tế}
Bản thể học y tế có \SL{kg/health/classes} lớp (Hình~\ref{fig:ontology} ở Chương~\ref{chap:kg}), xếp quanh lớp gốc
\textsf{MedicalEntity}: bệnh, tình trạng và biến chứng, chỉ số sinh học, yếu tố nguy cơ, sáu nhóm thuốc, chất và
thực phẩm phân loại theo mức an toàn (\textsf{SafeFood}, \textsf{ModerateFood}, \textsf{RiskFood}), lối sống và
hướng dẫn chuyên môn. Hình~\ref{fig:schema} là lược đồ: các lớp chính và vị từ được phép nối chúng.

\hinhtikz{t4_schema}{Lược đồ \texttt{VietnamDiabetesKG}: các lớp chính và vị từ nối chúng theo ràng buộc miền
$\to$ miền giá trị của \texttt{health\_ontology()}}{schema}

\maNguon{kg_health__health_ontology}{Bản thể học y tế: \SL{kg/health/classes} lớp và \SL{kg/health/predicates}
vị từ kèm ràng buộc}

Mỗi thực thể có nhãn tiếng Việt, bí danh và thuộc tính lâm sàng: \textsf{HbA1c\_Cao} có \texttt{threshold=6.5},
\textsf{Suy\_Than\_Nang} có \texttt{egfr\_max=30}, \textsf{Com\_Trang} có \texttt{gi=73, gl=29}. Bảng~\ref{tab:kgtypes}
đếm thực thể theo lớp của cả hai KG, Hình~\ref{fig:kg_health_graph} vẽ toàn bộ đồ thị y tế.

\bang{tab_kg_node_types}{Số thực thể theo lớp của hai đồ thị tri thức}{kgtypes}

\hinhpdf{kg_health_graph}{1.0}{Toàn cảnh \texttt{VietnamDiabetesKG}: \SL{kg/health/nodes} thực thể xếp vòng theo
lớp, \SL{kg/health/edges} quan hệ}{kg_health_graph}

\subsection{Mối quan hệ tri thức và ràng buộc lâm sàng}
\texttt{VietnamDiabetesKG} có \SL{kg/health/predicates} vị từ, chia thành bốn phân hệ: \emph{chẩn đoán}
(\textsf{TIEU\_CHUAN\_CHAN\_DOAN}, \textsf{TIEN\_TRIEN\_THANH}, \textsf{YEU\_TO\_NGUY\_CO}), \emph{điều trị}
(\textsf{CHI\_DINH\_DAU\_TAY} và nghịch đảo \textsf{DIEU\_TRI}, \textsf{UU\_TIEN\_KHI},
\textsf{CHONG\_CHI\_DINH\_KHI}), \emph{an toàn thuốc} (\textsf{TUONG\_TAC\_NGUY\_HIEM}, \textsf{TANG\_NGUY\_CO},
\textsf{BIEN\_CHUNG\_CUA}, \textsf{XU\_TRI\_BANG}) và \emph{dinh dưỡng, lối sống} (\textsf{KHUYEN\_NGHI\_AN},
\textsf{HAN\_CHE\_AN}, \textsf{AN\_DINH\_LUONG}, \textsf{LA\_MOT\_DANG\_CUA} bắc cầu, \textsf{KHUYEN\_NGHI\_LOI\_SONG}).
Bảng~\ref{tab:kgrel} và Hình~\ref{fig:kg_relations} đếm số cạnh mỗi vị từ.

\bang{tab_kg_relations}{Số cạnh theo vị từ của hai đồ thị tri thức}{kgrel}
\hinhpdf{kg_relations}{1.0}{Phân bố quan hệ theo vị từ của hai KG}{kg_relations}

Đoạn mã~\ref{code:kg_health__build_health_kg} là toàn bộ dữ liệu của KG y tế. Mỗi lời gọi \texttt{F(\dots)} thêm
một bộ ba kèm \texttt{source} là mã tài liệu trong kho và các thuộc tính lâm sàng (mức độ, hậu quả, quy tắc, lý
do). Mọi bộ ba đều qua bộ kiểm tra ràng buộc của Chương~\ref{chap:kg} trước khi vào đồ thị.

\maDai{kg_health__build_health_kg}{Dữ liệu của \texttt{VietnamDiabetesKG}: thực thể có thuộc tính và bộ ba có nguồn}

\section{Mô hình học sâu chẩn đoán nguy cơ}
\subsection{Dữ liệu}
Công cụ \texttt{chan\_doan\_nguy\_co} của tác tử là một mạng MLP huấn luyện trên bộ dữ liệu dự báo đái tháo
đường~\cite{mustafa2023diabetes}: hồ sơ gồm giới tính, tuổi, tăng huyết áp, bệnh tim, tiền sử hút thuốc, BMI, HbA1c,
glucose máu và nhãn mắc bệnh. Sau khi bỏ trùng còn \SL{diabetes/data/n} hồ sơ, tỉ lệ mắc bệnh
\SL{diabetes/data/pos_ratePct}\%, chia phân tầng thành ba tập (Bảng~\ref{tab:diadata}). Hình~\ref{fig:dia_dist}
cho thấy HbA1c và glucose tách hai lớp rõ nhất, còn tuổi và BMI chồng lấn nhiều.

\bang{tab_dia_data}{Bộ dữ liệu chẩn đoán đái tháo đường sau tiền xử lý}{diadata}
\hinhpdf{dia_dist}{1.0}{Phân bố các đặc trưng chính theo nhãn mắc bệnh}{dia_dist}

Hồ sơ được mã hóa thành \SL{diabetes/data/n_features} đặc trưng: bốn cột số (tuổi, BMI, HbA1c, glucose), giới
tính, hai cờ bệnh nền và sáu cột one-hot của tiền sử hút thuốc; mọi đặc trưng được chuẩn hóa z-score theo trung
bình và độ lệch chuẩn của tập huấn luyện.

\maNguon{data__diabetes_features}{Mã hóa hồ sơ thành \SL{diabetes/data/n_features} đặc trưng}

\subsection{Kiến trúc và huấn luyện}
Mạng \texttt{DiabetesMLP} có hai tầng ẩn 128 và 64 nơ-ron (Hình~\ref{fig:mlp3d}). Với $\mathbf{x} \in
\mathbb{R}^{13}$:
\begin{align}
  \mathbf{h}_1 &= \operatorname{Dropout}_{0{,}2}\bigl(\operatorname{ReLU}(\operatorname{BN}(W_1\mathbf{x} + \mathbf{b}_1))\bigr) \in \mathbb{R}^{128},\\
  \mathbf{h}_2 &= \operatorname{Dropout}_{0{,}2}\bigl(\operatorname{ReLU}(\operatorname{BN}(W_2\mathbf{h}_1 + \mathbf{b}_2))\bigr) \in \mathbb{R}^{64},\\
  z &= \mathbf{w}_3^\top \mathbf{h}_2 + b_3, \qquad \hat{p} = \sigma(z).
\end{align}
Tổng cộng \SL{diabetes/mlp/params} tham số. Vì lớp dương chỉ chiếm chưa tới một phần mười, hàm mất mát là
entropy chéo nhị phân có trọng số lớp dương $w_+ = N_-/N_+$:
\begin{equation}
  \mathcal{L} = -\frac{1}{N}\sum_{i=1}^{N}\Bigl[w_+\, y_i \log \sigma(z_i) + (1-y_i)\log\bigl(1-\sigma(z_i)\bigr)\Bigr].
\end{equation}
Mạng được tối ưu bằng AdamW (tốc độ học $10^{-3}$, suy giảm trọng số $10^{-4}$), lô 1024 mẫu, giảm tốc độ học một
nửa khi PR-AUC trên tập kiểm định không tăng sau 3 epoch, dừng sớm nếu không cải thiện sau 8 epoch, tối đa 60
epoch. Huấn luyện chạy trên GPU trong \SL{diabetes/mlp/train_s} s.

\hinhtikz{t4_mlp3d}{Kiến trúc \texttt{DiabetesMLP}: 13 đặc trưng $\to$ 128 $\to$ 64 $\to$ 1 logit, mỗi tầng ẩn
gồm BatchNorm, ReLU và Dropout}{mlp3d}

\maNguon{dl_models__DiabetesMLP}{Mạng \texttt{DiabetesMLP} trong PyTorch}
\maDai{train_diabetes__fit_mlp}{Vòng huấn luyện MLP: BCE có \texttt{pos\_weight}, AdamW, giảm tốc độ học và dừng
sớm theo PR-AUC}

Đường học (Hình~\ref{fig:dia_history}) hội tụ êm: mất mát huấn luyện và kiểm định giảm cùng nhịp, không có dấu
hiệu quá khớp; mạng chạy đủ \SL{diabetes/mlp/epochs} epoch mà không kích hoạt dừng sớm, tốc độ học được giảm
nhiều lần ở nửa sau khi PR-AUC kiểm định chững lại.

\hinhpdf{dia_history}{1.0}{Đường học của MLP: mất mát BCE và chỉ số trên tập kiểm định qua các epoch}{dia_history}

\subsection{Kết quả và phân tích}
Ngưỡng quyết định được chọn để F1 lớn nhất trên tập kiểm định, ra \SL{diabetes/mlp/threshold}. Ngưỡng cao như vậy
là hệ quả trực tiếp của $w_+$: mất mát có trọng số đẩy mọi xác suất lên, nên 0,5 không còn là điểm cân bằng.
Trên tập kiểm thử, MLP đạt F1 \SL{diabetes/mlp/f1} và ROC-AUC \SL{diabetes/mlp/roc_auc}, vượt Logistic Regression
ở mọi chỉ số (Bảng~\ref{tab:diamodels}, Hình~\ref{fig:dia_compare}).

\bang{tab_dia_models}{MLP và Logistic Regression trên tập kiểm thử}{diamodels}

\hinhpdf{dia_roc_pr}{1.0}{Đường ROC và Precision -- Recall của hai mô hình trên tập kiểm thử}{dia_roc_pr}
\hinhpdf{dia_confusion}{1.0}{Ma trận nhầm lẫn tại ngưỡng đã chọn}{dia_confusion}
\hinhpdf{dia_compare}{1.0}{So sánh MLP và Logistic Regression theo từng chỉ số}{dia_compare}

Ma trận nhầm lẫn (Hình~\ref{fig:dia_confusion}) cho thấy tính cách của mô hình: chỉ
\SL{diabetes/mlp/confusion/0/1} ca dương tính giả trên \SL{diabetes/data/n_test} hồ sơ, precision
\SL{diabetes/mlp/precision}, nhưng bỏ sót \SL{diabetes/mlp/confusion/1/0} ca bệnh, recall
\SL{diabetes/mlp/recall}. Với một công cụ \emph{sàng lọc}, ngưỡng nên hạ thấp để đổi precision lấy recall; với một
công cụ trong tác tử trả lời người bệnh, ngưỡng cao giảm báo động giả. Bài giữ ngưỡng F1 và để tác tử trình bày
xác suất cùng ngưỡng thay vì một kết luận có/không.

Hai phân tích sau giải thích mô hình đã học gì. Độ quan trọng hoán vị (Hình~\ref{fig:dia_importance}) đo PR-AUC
giảm bao nhiêu khi xáo trộn một đặc trưng: HbA1c (\SL{diabetes/importance/HbA1c}) và glucose
(\SL{diabetes/importance/Glucose}) chiếm gần như toàn bộ, tuổi và BMI đóng góp nhỏ, các cờ còn lại gần bằng
không. Đường đáp ứng theo HbA1c (Hình~\ref{fig:dia_hba1c_curve}) cố định các đặc trưng khác và quét HbA1c: xác suất
vượt ngưỡng quyết định ở HbA1c \SL{dem/dia_cross_140}\% khi glucose 140 mg/dL và \SL{dem/dia_cross_100}\% khi glucose
100 mg/dL, sát ngưỡng chẩn đoán 6,5\% của Bộ Y tế. Mô hình tự học lại được tiêu chuẩn lâm sàng từ dữ liệu.

\hinhpdf{dia_importance}{1.0}{Độ quan trọng hoán vị của các đặc trưng}{dia_importance}
\hinhpdf{dia_hba1c_curve}{1.0}{Xác suất mắc bệnh theo HbA1c ở ba mức glucose}{dia_hba1c_curve}

Biểu đồ tin cậy (Hình~\ref{fig:dia_calibration}) cho thấy giới hạn quan trọng nhất: xác suất của MLP \emph{không
hiệu chuẩn}. Ở ngăn có xác suất dự đoán trung bình \SL{diabetes/calibration/5/pred}, tỉ lệ mắc thực tế chỉ
\SL{diabetes/calibration/5/obs}; chỉ ngăn trên cùng là khớp. Đây vẫn là hệ quả của $w_+$. Vì vậy con số mà tác tử
và giao diện hiển thị phải đọc như một \emph{điểm nguy cơ} so với ngưỡng, không phải xác suất mắc bệnh; muốn dùng
như xác suất cần hiệu chuẩn thêm bằng Platt scaling hoặc hồi quy isotonic trên tập kiểm định. Bộ dữ liệu cũng là
dữ liệu mô phỏng có quan hệ rất sắc giữa HbA1c và nhãn, nên hồ sơ có HbA1c trên 7\% gần như luôn được chấm sát
100\%; mô hình không nên dùng cho quyết định lâm sàng thật.

\hinhpdf{dia_calibration}{1.0}{Biểu đồ tin cậy: xác suất dự đoán bị đẩy cao so với tần suất thực tế}{dia_calibration}

Khi chạy trong tác tử, lớp \texttt{DiabetesRiskModel} nạp trọng số và tham số chuẩn hóa, điền giá trị mặc định
cho trường người dùng không cung cấp và trả mức nguy cơ ba bậc (Đoạn mã~\ref{code:dl_models__DiabetesRiskModel.predict}).

\maNguon{dl_models__DiabetesRiskModel.predict}{Suy luận nguy cơ từ hồ sơ trong tác tử}

\section{Phân tích các tình huống lâm sàng thực tế}
Ba bài toán dưới đây đều là câu hỏi người bệnh Việt Nam thường hỏi, và mỗi bài cần một loại suy luận khác nhau
trên KG. Bảng~\ref{tab:demohealth} là kết quả nguyên văn của các truy vấn KG tương ứng.

\bang{tab_demo_health}{Kết quả truy vấn \texttt{VietnamDiabetesKG} cho ba bài toán lâm sàng}{demohealth}

\subsection{Bài toán 1: tương tác giữa metformin và rượu bia}
Câu hỏi ``đang uống metformin, tối nay đi tiệc uống bia được không'' cần hai mẩu tri thức. Mẩu thứ nhất là một
cạnh trực tiếp $\textsf{Metformin} \xrightarrow{\textsf{TUONG\_TAC\_NGUY\_HIEM}} \textsf{Ruou\_Bia}$ với mức độ
\texttt{cao}. Mẩu thứ hai không được ghi tường minh mà là một \emph{đường hai bước}: cả metformin và rượu bia đều có
cạnh \textsf{TANG\_NGUY\_CO} tới \textsf{Nhiem\_Toan\_Lactic}. Hàm \texttt{drug\_warnings} lấy cạnh trực tiếp và
giao hai tập nguy cơ để tìm nguy cơ chung (Đoạn mã~\ref{code:kg_health__drug_warnings}). Cơ chế sinh học đằng sau
được vẽ ở Hình~\ref{fig:metformin}: ethanol tiêu thụ NAD$^+$ ở gan và ức chế tân tạo glucose từ lactat, metformin
cũng giảm tân tạo glucose, hai tác động cộng lại làm lactat ứ đọng.

\hinhtikz{t4_metformin}{Cơ chế tương tác metformin -- rượu bia: hai con đường cùng làm giảm sử dụng lactat ở gan,
dẫn tới nguy cơ nhiễm toan lactic}{metformin}

\maNguon{kg_health__drug_warnings}{Cảnh báo tương tác: cạnh trực tiếp và nguy cơ chung qua đường hai bước}

Câu trả lời đúng không phải ``tuyệt đối cấm'' mà là cảnh báo có mức độ: tránh uống quá mức và tránh say, không
uống khi đói, theo dõi dấu hiệu bất thường; với người dùng sulfonylurea hoặc insulin, rượu còn làm tăng mạnh nguy cơ
hạ glucose huyết tới 24 giờ sau khi uống. Khung dưới đây là vết suy luận thật của tác tử cho câu hỏi này, do
\texttt{build.py} sinh từ \texttt{agent.json}.

\begin{vetsuyluan}{Vết suy luận ReAct thật cho câu hỏi y tế (kịch bản 1)}
\input{generated/trace_health}
\end{vetsuyluan}

Bộ kiểm chứng đã từ chối bản nháp của LLM và trả bản tổng hợp theo khuôn. Bản nháp (dưới đây) đúng về nội dung,
nhưng không ghi số hiệu bằng chứng [n] nào, nên không truy vết được từng ý về nguồn; Mục~\ref{sec:agent-ketqua}
phân tích hiện tượng này trên toàn bộ kịch bản.

\begin{vetsuyluan}{Bản nháp của Qwen2.5-1.5B (bị bộ kiểm chứng từ chối vì thiếu trích dẫn)}
\input{generated/llmdraft_health}
\end{vetsuyluan}

\subsection{Bài toán 2: dinh dưỡng thịt gà}
``Người tiểu đường ăn thịt gà được không'' không có một câu trả lời: ức gà bỏ da luộc là đạm nạc 31 g/100 g, không
carbohydrate, được khuyến nghị; gà rán và da gà nhiều chất béo bão hòa, làm tăng đề kháng insulin, nên hạn chế.
Trong KG, \textsf{Thit\_Ga} là một \textsf{FoodGroup}, các dạng cụ thể nối tới nó bằng \textsf{LA\_MOT\_DANG\_CUA}.
Hàm \texttt{food\_advice} nhận ra câu hỏi về \emph{nhóm} thực phẩm và duyệt xuống các dạng con trước khi tra lời
khuyên (Đoạn mã~\ref{code:kg_health__food_advice}), nên trả cả hai chiều khuyến nghị thay vì một câu ``được'' hay
``không'' sai một nửa.

\maNguon{kg_health__food_advice}{Lời khuyên dinh dưỡng: duyệt nhóm thực phẩm xuống các dạng cụ thể}

\subsection{Bài toán 3: định lượng tinh bột và ẩm thực truyền thống}
Chỉ số đường huyết GI so sánh tốc độ một thực phẩm làm tăng glucose máu với glucose chuẩn (GI = 100); GI dưới 55
là thấp, trên 70 là cao~\cite{atkinson2008gi}. Tải lượng đường huyết GL $=$ GI $\times$ lượng carbohydrate (g) $/
100$ tính cả khẩu phần. Các món tinh bột Việt Nam trải khắp thang GI (Hình~\ref{fig:gi}): bún, phở GI khoảng 53, gạo
lứt 68, cơm trắng 73, khoai tây luộc 78, xôi nếp rất cao. KG không cấm cơm mà mã hóa ba mức: \textsf{AN\_DINH\_LUONG}
cho cơm trắng (tối đa khoảng một bát con mỗi bữa, ăn rau trước), \textsf{HAN\_CHE\_AN} cho khoai tây và xôi,
\textsf{KHUYEN\_NGHI\_AN} cho rau xanh (ít nhất nửa đĩa). Truy vấn SPARQL Q2 ở Chương~\ref{chap:kg} trả đúng danh
sách món có GI, sắp theo GI giảm dần, nhờ suy diễn kiểu từ các lớp \textsf{ModerateFood}, \textsf{RiskFood} lên
\textsf{Food}.

\hinhtikz{t4_gi}{Thang chỉ số đường huyết GI của các món trong KG và mô hình đĩa ăn: nửa đĩa rau, một phần tư
đạm, một phần tư tinh bột}{gi}

\subsection{Cơ chế phản ứng đa tầng trước tình huống phức tạp}
Câu hỏi thật hiếm khi chỉ một ý. Người bệnh nam 52 tuổi, HbA1c 7,3\%, glucose lúc đói 138 mg/dL, BMI 26,5 hỏi
``tôi đang uống metformin 500 mg, tối nay đi liên hoan có uống bia và ăn gà rán được không''. Tác tử phản ứng theo
ba tầng (Hình~\ref{fig:ui_health}):

\begin{enumerate}
  \item \textbf{Tầng học sâu}: hồ sơ xét nghiệm được gửi kèm, nên bộ lập kế hoạch gọi \texttt{chan\_doan\_nguy\_co}
        trước; MLP chấm điểm nguy cơ vượt ngưỡng, mức ``cao''.
  \item \textbf{Tầng đồ thị}: liên kết thực thể nhận ra \textsf{Metformin}, \textsf{Ruou\_Bia}, \textsf{Ga\_Ran};
        \texttt{kg\_y\_te} trả cảnh báo tương tác, các chống chỉ định của metformin và lời khuyên hạn chế gà rán.
  \item \textbf{Tầng văn bản}: \texttt{truy\_xuat\_rag} lấy các đoạn về rượu bia trong bữa tiệc và cơ chế nhiễm toan
        lactic từ tài liệu \texttt{h03-ruou-bia} để giải thích.
\end{enumerate}

Mỗi tầng trả lời một phần mà tầng khác không trả lời được: MLP không biết gì về bia, KG không biết người này có
nguy cơ cao, văn bản không chấm điểm được hồ sơ. Hình~\ref{fig:ui_health_food} là câu hỏi về cơm trắng trên cùng
giao diện, khi không gửi kèm hồ sơ: bộ lập kế hoạch bỏ qua MLP và chỉ dùng KG cùng RAG.

\anhUI{ui_health}{Trợ lý y tế trên Web App: hồ sơ xét nghiệm, vết suy luận ReAct, câu trả lời có trích dẫn và danh
sách bằng chứng}{ui_health}
\anhUI{ui_health_food}{Câu hỏi về cơm trắng không kèm hồ sơ: tác tử chỉ gọi KG và RAG}{ui_health_food}

% =====================================================================
%  CHƯƠNG 5 — HỆ THỐNG TÀI CHÍNH: VN30 VÀ GIÁ VÀNG SJC
% =====================================================================
\chapter{Hệ thống tài chính: VN30 và giá vàng SJC}\label{chap:taichinh}

Phân hệ tài chính trả lời loại câu hỏi mà nhà đầu tư cá nhân Việt Nam hay đặt ra: tỷ giá USD/VND tăng thì cổ phiếu
nào hưởng lợi, Fed tăng lãi suất thì vàng SJC đi đâu, phiên tới FPT có thể ở mức giá nào. Ba câu hỏi đó cần ba loại
tri thức khác nhau. Câu đầu cần \emph{cơ chế} (doanh thu bằng ngoại tệ, chi phí bằng đồng); câu thứ hai cần \emph{suy
luận nhiều bước} qua các biến trung gian (DXY, giá vàng thế giới, giá quy đổi); câu thứ ba cần \emph{dữ liệu chuỗi
thời gian}. Chương này dựng ba lớp tương ứng: \texttt{VietnamFinanceKG} lưu đồ thị nhân quả có dấu và có nguồn,
bộ suy luận đa bước tính chiều tác động ròng trên đồ thị đó, và mạng BiLSTM dự báo lợi suất phiên kế tiếp. Lớp thứ
tư, ít hiển nhiên hơn, là bước \emph{kiểm chứng KG bằng dữ liệu}: dấu của mỗi cạnh lý thuyết được đối chiếu với
tương quan thực nghiệm trên gần 12 năm giao dịch.

\begin{phathien}{Lưu ý về phạm vi}
Hệ thống phục vụ mục đích học tập, không phải công cụ khuyến nghị đầu tư. Kho văn bản tài chính tổng hợp từ văn bản
pháp luật (Nghị định 24/2012/NĐ-CP, Nghị định 232/2025/NĐ-CP), thông cáo của Ngân hàng Nhà nước và báo cáo thường
niên doanh nghiệp; mọi câu trả lời của tác tử đều kết thúc bằng câu ``Thông tin chỉ mang tính tham khảo, không phải
khuyến nghị đầu tư''.
\end{phathien}

% ---------------------------------------------------------------------
\section{Bối cảnh vĩ mô và dữ liệu thị trường}\label{sec:tc-boicanh}
\subsection{Bốn cổ phiếu, bốn kênh truyền dẫn}
Rổ VN30 gồm 30 cổ phiếu vốn hóa và thanh khoản lớn nhất sàn HOSE. Bài chọn bốn mã không phải vì vốn hóa mà vì mỗi
mã đại diện một kênh truyền dẫn vĩ mô khác nhau, để KG có đủ dấu dương và dấu âm cho cùng một biến nguồn:

\begin{itemize}
  \item \textbf{FPT} (công nghệ thông tin): doanh thu xuất khẩu phần mềm thanh toán bằng USD, JPY, EUR, còn chi phí
        chủ yếu là lương kỹ sư trả bằng đồng. Tỷ giá tăng làm doanh thu quy đổi tăng nhanh hơn chi phí, nên về mặt
        cơ bản FPT \emph{hưởng lợi} khi USD/VND tăng; Nhật Bản là thị trường xuất khẩu lớn nhất nên FPT còn nhạy với
        đồng yên.
  \item \textbf{HPG} (thép): nhập quặng sắt và than cốc bằng USD và có nợ vay ngoại tệ, nên tỷ giá tăng làm giá vốn
        tăng và phát sinh lỗ tỷ giá; ngược lại, giải ngân đầu tư công làm tăng cầu thép xây dựng.
  \item \textbf{VCB} (ngân hàng): lợi nhuận phụ thuộc tăng trưởng tín dụng, mà tín dụng nhạy với lãi suất điều hành
        của Ngân hàng Nhà nước (NHNN).
  \item \textbf{VIC} (bất động sản, ô tô): chi phí tài chính lớn và có nợ vay ngoại tệ, chịu cả hai kênh lãi suất
        và tỷ giá.
\end{itemize}

Chỉ số VN30 được đại diện bằng giá chứng chỉ quỹ ETF mô phỏng VN30 (chuỗi \texttt{VN30ETF}). Ngoài bốn cổ phiếu, dữ liệu gồm giá vàng thế giới XAU/USD, tỷ giá USD/VND và chỉ số đồng
đô la DXY. Hàm \texttt{market\_frame} lấy lịch giao dịch của HOSE làm trục thời gian và điền tiến các chuỗi nước
ngoài vào ngày giao dịch Việt Nam, nên mọi tài sản cùng có \SL{market/n_days} phiên chung từ \SL{market/range/0} tới
\SL{market/range/1}.

\bang{tab_market_desc}{Thống kê mô tả các tài sản trong giai đoạn nghiên cứu}{marketdesc}

\hinhpdf{mk_prices}{1.0}{Diễn biến giá điều chỉnh của bốn cổ phiếu, ETF VN30 và vàng quy đổi, chuẩn hóa về 100
ở đầu kỳ (thang log)}{mk_prices}

Bảng~\ref{tab:marketdesc} và Hình~\ref{fig:mk_risk} cho thấy lợi suất cao luôn đi cùng rủi ro cao. VIC có lợi suất
năm hóa cao nhất (\SL{market/describe/VIC/ann_retPct}\%) nhưng từng sụt \SL{market/describe/VIC/max_ddPct}\% từ
đỉnh; HPG sụt sâu nhất, \SL{market/describe/HPG/max_ddPct}\%, đúng tính chu kỳ của ngành thép. Vàng quy đổi ra
đồng có lợi suất \SL{market/describe/GOLD_VND/ann_retPct}\% mỗi năm, ngang VCB, nhưng mức sụt tối đa chỉ
\SL{market/describe/GOLD_VND/max_ddPct}\%, vì vậy vàng đóng vai trò tài sản phòng thủ trong danh mục của người Việt.
Tỷ giá USD/VND biến động chỉ \SL{market/describe/USDVND/ann_volPct}\% mỗi năm do NHNN điều hành trong biên độ, nhưng
chính vì ít biến động nên mỗi lần tỷ giá tăng nhanh đều là tin lớn với thị trường.

\hinhpdf{mk_risk}{0.92}{Lợi suất -- rủi ro năm hóa; diện tích bong bóng tỷ lệ với mức sụt giảm tối đa}{mk_risk}

\subsection{Tương quan giữa các tài sản}
Ma trận tương quan (Hình~\ref{fig:mk_corr}) được tính trên hai thang: lợi suất ngày và lợi suất cộng dồn 20 phiên.
Thang 20 phiên ít nhiễu hơn và gần với tầm nhìn của nhà đầu tư cá nhân. Ba nhóm quan hệ hiện ra rõ:

\begin{enumerate}
  \item Bốn cổ phiếu cùng tương quan dương với ETF VN30 (từ \SL{market/corr_20d/VIC/VN30ETF} với VIC tới
        \SL{market/corr_20d/VCB/VN30ETF} với VCB): rủi ro chung của thị trường chiếm phần lớn biến động.
  \item USD/VND tương quan \emph{âm} với cả bốn cổ phiếu (FPT \SL{market/corr_20d/USDVND/FPT}, HPG
        \SL{market/corr_20d/USDVND/HPG}): tỷ giá tăng thường đi cùng khối ngoại bán ròng và tâm lý rủi ro.
  \item Vàng thế giới gần như quyết định vàng quy đổi (\SL{market/corr_20d/XAU/GOLD_VND}), và ngược chiều với DXY
        (\SL{market/corr_20d/DXY/XAU}) vì vàng được định giá bằng USD.
\end{enumerate}

\hinhpdf{mk_corr}{1.0}{Ma trận tương quan Pearson của lợi suất ngày (trái) và lợi suất 20 phiên (phải)}{mk_corr}

Tương quan không đứng yên. Hình~\ref{fig:mk_rolling} vẽ tương quan trượt 252 phiên (khoảng một năm) của các cặp
tương ứng với cạnh trong KG: cặp DXY -- vàng luôn âm, còn cặp tỷ giá -- cổ phiếu dao động quanh 0 và đổi dấu theo
giai đoạn. Đây là cảnh báo đầu tiên rằng một cạnh có dấu cố định trong KG chỉ mô tả \emph{xu hướng cơ bản}, không
phải quy luật đúng ở mọi thời điểm.

\hinhpdf{mk_rolling}{1.0}{Tương quan trượt một năm của các cặp tương ứng với cạnh trong \texttt{VietnamFinanceKG}}{mk_rolling}

\subsection{Thị trường vàng Việt Nam: giá quy đổi và phần bù SJC}\label{sec:vang}
Giá vàng trong nước niêm yết theo \emph{lượng} (37,5 g), còn giá thế giới theo \emph{ounce troy} (31,1035 g). Giá
vàng thế giới quy đổi ra đồng, chưa gồm thuế phí, là

\begin{equation}\label{eq:vangquydoi}
  G_t \;=\; \frac{\mathrm{XAU}_t \times (\mathrm{USD/VND})_t \times 1{,}2057}{10^6}
  \quad\text{(triệu đồng/lượng)},
\end{equation}

\noindent với hệ số $1{,}2057 \approx 37{,}5 / 31{,}1035$ đổi giá mỗi ounce sang giá mỗi lượng. Giá vàng miếng SJC
bằng giá quy đổi cộng một \emph{phần bù} do cung cầu trong nước quyết định:
$P^{\mathrm{SJC}}_t = G_t + \Delta_t$. Ngày \SL{market/gold/last/date}, XAU/USD $=$ \SL{market/gold/last/xau}
USD/ounce và USD/VND $=$ \SL{market/gold/last/usdvnd}, nên $G_t \approx$ \SL{market/gold/last/gold_vnd} triệu
đồng/lượng. Công thức~\eqref{eq:vangquydoi} chính là công cụ \texttt{quy\_doi\_vang} của tác tử, và là lý do hai cạnh
\textsf{XAU\_USD} $\to$ \textsf{Vang\_Quy\_Doi}, \textsf{USD\_VND} $\to$ \textsf{Vang\_Quy\_Doi} trong KG đều mang
dấu dương.

\hinhtikz{t5_gold}{Cấu thành giá vàng miếng SJC: giá thế giới, tỷ giá và hệ số đổi đơn vị tạo nên giá quy đổi; phần
bù trong nước phụ thuộc chính sách}{t5_gold}

\hinhpdf{mk_gold}{1.0}{XAU/USD, USD/VND và giá vàng quy đổi theo~\eqref{eq:vangquydoi}}{mk_gold}

Phần bù $\Delta_t$ là đại lượng chịu tác động của chính sách, và là nơi KG cần tri thức pháp luật:

\begin{itemize}
  \item \textbf{Nghị định 24/2012/NĐ-CP}~\cite{nd24_2012} quy định Nhà nước độc quyền sản xuất vàng miếng, độc quyền
        xuất nhập khẩu vàng nguyên liệu để sản xuất vàng miếng; NHNN chọn SJC làm thương hiệu vàng miếng quốc gia.
        Từ năm 2014 NHNN gần như không nhập vàng nguyên liệu, nguồn cung bị giới hạn trong khi nhu cầu tích trữ tăng,
        nên phần bù có lúc lên tới 15 -- 18 triệu đồng/lượng đầu năm 2024.
  \item \textbf{Can thiệp năm 2024}: NHNN đấu thầu vàng miếng từ tháng 4, rồi bán trực tiếp qua bốn ngân hàng thương
        mại nhà nước và Công ty SJC từ tháng 6; phần bù thu hẹp còn khoảng 3 -- 5 triệu đồng/lượng.
  \item \textbf{Nghị định 232/2025/NĐ-CP}~\cite{nd232_2025} sửa đổi Nghị định 24, có hiệu lực từ tháng 10/2025,
        bãi bỏ cơ chế độc quyền sản xuất vàng miếng và cho phép doanh nghiệp, ngân hàng đủ điều kiện sản xuất vàng
        miếng, nhập khẩu vàng nguyên liệu theo hạn mức do NHNN quản lý.
\end{itemize}

Trong KG, ba sự kiện này là ba cạnh vào nút \textsf{Chenh\_Lech\_SJC}: Nghị định 24 làm phần bù tăng (dấu $+$),
hai biện pháp còn lại làm phần bù giảm (dấu $-$). Vàng nhẫn 9999 không chịu cơ chế độc quyền nên chỉ nối với giá
quy đổi, không nối với phần bù.

% ---------------------------------------------------------------------
\section{Đồ thị tri thức tài chính \texttt{VietnamFinanceKG}}\label{sec:kgtc}
\subsection{Bản thể học}
Bản thể học tài chính nhỏ hơn bản thể học y tế nhưng có cấu trúc khác: \SL{kg/finance/classes} lớp và chỉ
\SL{kg/finance/predicates} vị từ. Lớp gốc \textsf{FinancialEntity} tách thành yếu tố vĩ mô (\textsf{MacroFactor},
với ba lớp con \textsf{PolicyTool}, \textsf{GlobalFactor}, \textsf{MarketFlow}), tài sản (\textsf{Asset} $\to$
\textsf{Commodity} $\to$ \textsf{GoldAsset}, \textsf{Stock}, \textsf{Index}), văn bản pháp luật, sự kiện, chỉ báo,
ngành và sàn giao dịch. Hai vị từ quan trọng nhất mang \emph{dấu}: \textsf{TAC\_DONG\_CUNG\_CHIEU} ($+1$, biến
nguồn tăng thì biến đích có xu hướng tăng) và \textsf{TAC\_DONG\_NGUOC\_CHIEU} ($-1$). Bốn vị từ còn lại là quan hệ
cấu trúc (thành viên chỉ số, ngành, sửa đổi văn bản, phạm vi điều chỉnh) và không tham gia lan truyền tác động.
\textsf{THUOC\_VE} được khai báo bắc cầu, nên suy diễn RDFS (Hình~\ref{fig:suydien}) tự sinh ``FPT thuộc về HOSE'' từ
hai cạnh ``FPT thuộc VN30'' và ``VN30 thuộc HOSE''.

\maNguon{kg_finance__finance_ontology}{Bản thể học tài chính: \SL{kg/finance/classes} lớp và
\SL{kg/finance/predicates} vị từ, hai vị từ nhân quả mang dấu}

\subsection{Thực thể và quan hệ nhân quả}
KG có \SL{kg/finance/nodes} thực thể và \SL{kg/finance/edges} cạnh, trong đó
\SL{kg/finance/relations/TAC_DONG_CUNG_CHIEU} cạnh cùng chiều và \SL{kg/finance/relations/TAC_DONG_NGUOC_CHIEU}
cạnh ngược chiều. Cả \SL{kg/finance/with_source} cạnh đều có thuộc tính \texttt{source} trỏ về tài liệu trong kho
văn bản (\texttt{f01} -- \texttt{f09}) và thuộc tính \texttt{mechanism} giải thích cơ chế bằng một câu, nên mỗi kết
luận của tác tử đều truy được về văn bản gốc. Thực thể có chuỗi giá thật mang thuộc tính \texttt{series}, nối KG với
dữ liệu thị trường; nhờ đó Mục~\ref{sec:kiemchung} kiểm chứng được dấu cạnh. Hình~\ref{fig:t5_causal} tách đồ thị
nhân quả thành hai kênh để dễ đọc, Hình~\ref{fig:kg_finance_graph} vẽ toàn bộ KG.

\hinhtikz{t5_causal}{Đồ thị nhân quả của \texttt{VietnamFinanceKG}: (a) kênh vàng, (b) kênh tỷ giá -- lãi suất --
cổ phiếu; xanh liền là cùng chiều, đỏ đứt là ngược chiều}{t5_causal}

\maDai{kg_finance__build_finance_kg}{Xây \texttt{VietnamFinanceKG}: thực thể kèm bí danh, chuỗi giá, và các cạnh
nhân quả có nguồn, có cơ chế}

\hinhpdf{kg_finance_graph}{1.0}{Toàn cảnh \texttt{VietnamFinanceKG}: \SL{kg/finance/nodes} thực thể,
\SL{kg/finance/edges} quan hệ; màu nút theo lớp, màu cạnh theo dấu tác động}{kg_finance_graph}

Khi xuất ra RDF, KG có \SL{kg/finance/rdf_triples} bộ ba; bộ suy diễn RDFS thêm \SL{kg/finance/inferred} bộ ba
(kiểu kế thừa từ lớp cha, quan hệ bắc cầu) để tổng số đạt \SL{kg/finance/rdf_after}. Chẳng hạn ``FPT là một
\textsf{Stock}'' kéo theo ``FPT là một \textsf{Asset}'' và ``FPT là một \textsf{FinancialEntity}''. Nhờ vậy truy vấn SPARQL hỏi ``mọi tài sản chịu tác động ngược chiều của tỷ giá'' không
cần liệt kê từng lớp con.

% ---------------------------------------------------------------------
\section{Suy luận tác động đa bước}\label{sec:multihop}
\subsection{Dấu của đường đi và tác động ròng}
Với một biến nguồn $u$ (chẳng hạn lãi suất Fed) và một biến đích $v$ (vàng SJC), hàm \texttt{impact\_paths} liệt kê
mọi đường đơn $\pi = (u = x_0 \to x_1 \to \dots \to x_k = v)$ dài không quá 4 bước trên đồ thị con nhân quả. Dấu của
đường là tích dấu các cạnh, và tác động ròng là tổng dấu các đường có trọng số nghịch đảo độ dài:

\begin{equation}\label{eq:rong}
  s(\pi) \;=\; \prod_{j=1}^{k} s(x_{j-1}, x_j), \qquad
  N(u, v) \;=\; \sum_{\pi:\, u \rightsquigarrow v} \frac{s(\pi)}{|\pi|}.
\end{equation}

Trọng số $1/|\pi|$ thể hiện giả định rằng tác động suy giảm qua mỗi bước trung gian: một kênh trực tiếp đáng tin hơn
một chuỗi bốn mắt xích. Khi mọi đường cùng dấu, kết luận là ``tích cực'' hoặc ``tiêu cực''; khi các đường mâu thuẫn,
bộ suy luận \emph{không} gộp thành một chiều mà trả về ``hỗn hợp'' kèm số kênh mỗi dấu, để tác tử nói rõ kết quả phụ
thuộc kênh truyền dẫn nào mạnh hơn.

\maNguon{kg_finance__impact_paths}{Suy luận đa bước: liệt kê đường nhân quả, tính dấu từng đường và tác động
ròng~\eqref{eq:rong}}

\hinhtikz{t5_multihop}{Ba đường nhân quả từ lãi suất Fed tới vàng SJC mà \texttt{impact\_paths} tìm được, dấu của
từng đường và tác động ròng}{t5_multihop}

\bang{tab_multihop}{Kết quả suy luận đa bước cho bảy cặp nguồn -- đích}{multihop}

Bảng~\ref{tab:multihop} có hai trường hợp ``hỗn hợp'' đáng phân tích:

\begin{itemize}
  \item \textbf{Fed $\to$ vàng SJC}: hai đường âm (Fed $\to$ XAU $\to$ quy đổi $\to$ SJC, và Fed $\to$ DXY $\to$ XAU
        $\to$ quy đổi $\to$ SJC) và một đường dương (Fed $\to$ DXY $\to$ USD/VND $\to$ quy đổi $\to$ SJC). Đường dương
        có thật về kinh tế: Fed tăng lãi suất làm USD mạnh, tỷ giá tăng, và giá vàng tính bằng đồng được đỡ một phần.
        Tác động ròng $N = -1/3 - 1/4 + 1/4 =$ \SL{kg/demo_finance/Lai_Suat_Fed->Vang_SJC/net}: chiều chi phí cơ hội
        của vàng thắng, nhưng nhà đầu tư trong nước chịu giảm ít hơn nhà đầu tư nắm vàng tính bằng USD.
  \item \textbf{Lãi suất NHNN $\to$ HPG}: một đường trực tiếp âm (lãi suất cao làm giảm cầu xây dựng) và một đường
        hai bước dương (lãi suất VND cao làm giảm áp lực tỷ giá, mà tỷ giá thấp có lợi cho doanh nghiệp nhập nguyên
        liệu). Ròng $-1 + 1/2 =$ \SL{kg/demo_finance/Lai_Suat_NHNN->HPG/net}.
\end{itemize}

\subsection{Lan truyền cú sốc}
Câu hỏi ngược lại, ``một cú sốc ở $u$ ảnh hưởng tới những gì'', được trả lời bằng hàm \texttt{downstream}: duyệt
theo chiều rộng từ $u$ tối đa 3 bước, mỗi thực thể nhận dấu của đường ngắn nhất tới nó. Hình~\ref{fig:kg_shock} và
Bảng~\ref{tab:shock} là kết quả cho cú sốc DXY tăng: mười thực thể bị ảnh hưởng. FPT là cổ phiếu duy nhất có lợi
(qua tỷ giá); VN30, HPG, VIC và cả bốn nút vàng đều chịu chiều giảm; VCB không nằm trong danh sách vì KG không có
kênh nối VCB với tỷ giá.

\maNguon{kg_finance__downstream}{Lan truyền cú sốc theo chiều rộng, giữ dấu của đường ngắn nhất}

\hinhpdf{kg_shock}{0.92}{Lan truyền cú sốc ``DXY tăng'' trong \texttt{VietnamFinanceKG}: xanh là tăng, đỏ là
giảm}{kg_shock}

\bang{tab_shock}{Chiều tác động lên các thực thể khi DXY tăng, theo đường ngắn nhất}{shock}

\subsection{Kiểm chứng dấu cạnh bằng dữ liệu}\label{sec:kiemchung}
Một KG tài chính viết tay dễ chứa những ``hiểu biết thông thường'' sai. Bài kiểm chứng mọi cạnh nhân quả có chuỗi
giá ở cả hai đầu: tính tương quan Pearson giữa lợi suất 20 phiên của biến nguồn và biến đích trên toàn giai đoạn, rồi
so dấu với dấu lý thuyết. Kết quả (Bảng~\ref{tab:edgeev}, Hình~\ref{fig:mk_edges}): \SL{dem/edges_agree} trên
\SL{dem/edges_checked} cạnh khớp dấu. Cạnh \textsf{XAU\_USD} $\to$ \textsf{Vang\_Quy\_Doi} gần như tất định
(\SL{market/corr_20d/XAU/GOLD_VND}); cạnh DXY $\to$ XAU mạnh và đúng chiều; các cạnh tỷ giá $\to$ HPG, VIC yếu
nhưng đúng dấu.

\bang{tab_edges_evidence}{Đối chiếu dấu lý thuyết của cạnh KG với tương quan thực nghiệm}{edgeev}

\hinhpdf{mk_edges}{1.0}{Tương quan thực nghiệm (thanh) và dấu lý thuyết trong KG (hình thoi) của bảy cạnh có chuỗi
giá ở hai đầu}{mk_edges}

Cạnh không khớp là USD/VND $\to$ FPT: lý thuyết cho dấu dương, dữ liệu cho tương quan cùng kỳ
\SL{market/corr_20d/USDVND/FPT}. Sai lệch này không có nghĩa cạnh sai, mà cho thấy hai cơ chế chạy ở hai thang thời
gian khác nhau. Trong ngắn hạn, tỷ giá tăng nhanh đi cùng dòng vốn ngoại rút khỏi thị trường và tâm lý né rủi ro, kéo
mọi cổ phiếu lớn giảm cùng lúc, kể cả FPT. Lợi ích cơ bản của doanh thu ngoại tệ chỉ hiện ra khi doanh nghiệp công
bố kết quả kinh doanh, tức là vài tháng sau. Tương quan trễ (Hình~\ref{fig:mk_lag}) xác nhận cách giải thích này:
tương quan giữa lợi suất 20 phiên của tỷ giá hôm nay và của FPT $k$ phiên sau tăng dần theo $k$, đổi dấu quanh
$k = 20$ và đạt cực đại \SL{dem/lag_corr_FPT} ở độ trễ \SL{dem/lag_peak_FPT} phiên, trong khi với HPG, đường tương
quan chỉ tiến về 0 mà không đổi dấu đáng kể.

\hinhpdf{mk_lag}{1.0}{Tương quan giữa lợi suất 20 phiên của USD/VND và của cổ phiếu sau $k$ phiên}{mk_lag}

\nhanxet{Cạnh KG mã hóa cơ chế cơ bản, còn tương quan cùng kỳ đo phản ứng của dòng tiền. Bài giữ nguyên cạnh
USD/VND $\to$ FPT vì có căn cứ văn bản; đây là ví dụ cho thấy KG tài chính cần được đối chiếu với dữ liệu thay vì chỉ
dựa vào hiểu biết thông thường, và là lý do Chương~\ref{chap:ketluan} đề xuất gắn tương quan thực nghiệm cùng độ trễ
vào thuộc tính của từng cạnh. Một lưu ý về phương pháp: lợi suất 20 phiên tính trên cửa sổ chồng lấn nên các quan sát liền nhau phụ
thuộc nhau; độ lớn tương quan dùng được để so chiều, nhưng không nên đọc như kiểm định ý nghĩa thống kê.}

% ---------------------------------------------------------------------
\section{Dự báo chuỗi thời gian bằng BiLSTM}\label{sec:bilstm}
\subsection{Phát biểu bài toán}
Mô hình không dự báo trực tiếp giá mà dự báo \emph{lợi suất log} phiên kế tiếp, $r_{t+1} = \ln(P_{t+1}/P_t)$,
rồi suy ra giá $\hat P_{t+1} = P_t\, e^{\hat r_{t+1}}$. Chuỗi giá không dừng: một mô hình hồi quy trên giá đạt $R^2$
rất cao chỉ bằng cách chép giá hôm qua, và con số đó không nói gì về khả năng dự báo. Lợi suất gần dừng hơn, và cho
phép so sánh với mốc khó nhất: \emph{dự báo ngây thơ} ``giá không đổi'', tức $\hat r_{t+1} = 0$. Với mốc này,
$\mathrm{RMSE}_{\text{ngây thơ}} = \sqrt{\tfrac{1}{n}\sum r_{t+1}^2}$, và chỉ số chính của bài là tỉ số
$\mathrm{RMSE}/\mathrm{RMSE}_{\text{ngây thơ}}$: dưới 1 mới là có năng lực dự báo.

Mỗi phiên có 7 đặc trưng với cổ phiếu (6 với vàng): lợi suất log của chính tài sản; lợi suất của bốn biến ngoại sinh
USD/VND, XAU, DXY, ETF VN30 (bỏ XAU khi dự báo vàng quy đổi vì XAU là thành phần của chính nó); độ lệch của giá so
với trung bình trượt 20 phiên; và độ biến động 20 phiên. Mỗi mẫu là một cửa sổ
\SL{forecast/config/window} phiên liên tiếp, tức một tenxơ $30 \times 7$.

\maNguon{data__forecast_features}{Đặc trưng dự báo: lợi suất của đích và biến ngoại sinh, độ lệch MA20, độ biến
động 20 phiên; \texttt{keep\_last} giữ phiên mới nhất khi suy luận}

\maNguon{data__make_windows}{Cắt chuỗi đặc trưng thành các cửa sổ trượt cho mô hình tuần tự}

Dữ liệu được chia \emph{theo thời gian}, không xáo trộn (Hình~\ref{fig:t5_split}): huấn luyện trước
\SL{forecast/config/val_start}, kiểm định năm 2024, kiểm thử từ \SL{forecast/config/test_start} tới hết dữ liệu.
Trung bình và độ lệch chuẩn để chuẩn hóa chỉ tính trên tập huấn luyện. Với FPT, ba tập có \SL{forecast/FPT/n/train},
\SL{forecast/FPT/n/val} và \SL{forecast/FPT/n/test} cửa sổ.

\hinhtikz{t5_split}{Chia dữ liệu theo thời gian và cửa sổ trượt 30 phiên: mô hình không bao giờ nhìn thấy tương lai
của mẫu nó dự báo}{t5_split}

\subsection{Kiến trúc LSTM hai chiều}
LSTM~\cite{hochreiter1997lstm} giải quyết bài toán gradient biến mất của RNN bằng một ô nhớ $c_t$ và ba cổng điều
khiển luồng thông tin. Với đầu vào $x_t$ và trạng thái ẩn trước $h_{t-1}$:

\begin{align}
  f_t &= \sigma\!\left(W_f [h_{t-1}, x_t] + b_f\right) & &\text{cổng quên} \label{eq:lstm-f}\\
  i_t &= \sigma\!\left(W_i [h_{t-1}, x_t] + b_i\right), \quad
  \tilde c_t = \tanh\!\left(W_c [h_{t-1}, x_t] + b_c\right) & &\text{cổng vào, ứng viên} \\
  c_t &= f_t \odot c_{t-1} + i_t \odot \tilde c_t & &\text{cập nhật ô nhớ} \\
  o_t &= \sigma\!\left(W_o [h_{t-1}, x_t] + b_o\right), \quad h_t = o_t \odot \tanh(c_t) & &\text{cổng ra} \label{eq:lstm-o}
\end{align}

LSTM hai chiều~\cite{schuster1997brnn} chạy thêm một LSTM ngược thời gian trên \emph{cùng cửa sổ} và ghép hai trạng
thái, $h_t = [\overrightarrow{h}_t ; \overleftarrow{h}_t]$. Vì cửa sổ chỉ gồm các phiên đã qua, chiều ngược đọc từ
phiên $t$ lùi về phiên $t-29$ và không làm rò rỉ tương lai; tác dụng của nó là để trạng thái cuối mang thông tin của
cả đầu lẫn cuối cửa sổ. Mạng gồm hai tầng BiLSTM 64 chiều mỗi hướng, lấy trạng thái ở bước cuối (128 chiều) qua
Dropout và hai lớp tuyến tính $128 \to 32 \to 1$; tổng cộng 140.865 tham số với 7 đặc trưng, so với 54.081 tham số
của LSTM một chiều cùng cấu hình.

\hinhtikz{t5_bilstm3d}{Bộ dự báo BiLSTM: tenxơ cửa sổ $30 \times 7$, hai tầng LSTM xuôi và ngược thời gian, đầu
tuyến tính cho lợi suất phiên kế tiếp và MC Dropout cho khoảng dự báo}{t5_bilstm3d}

\maNguon{dl_models__SeqForecaster}{Mô hình tuần tự dùng chung cho BiLSTM và LSTM một chiều}

\bang{tab_fc_config}{Cấu hình bộ dự báo}{fcconfig}

\subsection{Huấn luyện}
Hàm mất mát là Huber~\cite{huber1964robust} với $\delta = 1$ trên lợi suất đã chuẩn hóa:
\begin{equation}
  \mathcal{L}_\delta(e) = \begin{cases} \tfrac12 e^2 & |e| \le \delta \\ \delta\left(|e| - \tfrac12\delta\right) &
  |e| > \delta \end{cases}
\end{equation}
Lợi suất có đuôi dày: vài phiên giảm sàn hay tăng trần có thể chi phối MSE, trong khi Huber chỉ phạt tuyến tính các
sai số lớn. Tối ưu bằng AdamW~\cite{loshchilov2019adamw} (tốc độ học $10^{-3}$, suy giảm trọng số $10^{-4}$), lô 128,
cắt chuẩn gradient ở 1,0 để tránh bùng nổ gradient trong LSTM, tối đa 80 epoch và dừng sớm khi mất mát kiểm định
không giảm sau 10 epoch. Huấn luyện chạy trên GPU, mỗi mô hình mất vài giây
(Bảng~\ref{tab:forecast}). Đối chứng gồm LSTM một chiều cùng cấu hình, hồi quy Ridge trên 30 trễ của mọi đặc
trưng, và dự báo ngây thơ.

\maDai{train_forecast__train_rnn}{Vòng huấn luyện: Huber, AdamW, cắt gradient, dừng sớm theo mất mát kiểm định}

\maNguon{train_forecast__evaluate}{Đánh giá trên tập kiểm thử: RMSE lợi suất, tỉ số với dự báo ngây thơ, MAE giá,
độ chính xác hướng}

\hinhpdf{fc_history}{1.0}{Đường cong học của BiLSTM cho FPT và vàng quy đổi}{fc_history}

Đường cong học (Hình~\ref{fig:fc_history}) cho thấy dấu hiệu đầu tiên của kết quả: mất mát huấn luyện giảm đều,
còn mất mát kiểm định gần như đi ngang hoặc tăng ngay từ những epoch đầu. Với HPG, VIC và vàng quy đổi, trọng số tốt
nhất là của epoch thứ nhất và quá trình dừng ở epoch 11; với FPT, mất mát kiểm định tốt nhất (epoch 22) chỉ thấp hơn
epoch đầu khoảng 2\%. Phần mô hình học thêm từ tập huấn luyện gần như không tổng quát hóa sang năm 2024.

\subsection{Kết quả và diễn giải}
\bang{tab_forecast}{Kết quả dự báo trên tập kiểm thử của năm tài sản và bốn mô hình}{forecast}

\haihinh{fc_ratio}{Tỉ số RMSE so với dự báo ngây thơ}{fc_diracc}{Tỉ lệ đoán đúng chiều tăng/giảm}{So sánh BiLSTM,
LSTM và Ridge trên năm tài sản}{fc_compare}

Kết quả cần được nói thẳng. Chỉ \SL{dem/fc_beat} trên \SL{dem/fc_models} mô hình có tỉ số RMSE dưới 1, và mô hình
tốt nhất (LSTM trên vàng quy đổi) chỉ giảm sai số \SL{dem/fc_best_gainPct}\% so với dự báo ngây thơ, một mức không có
ý nghĩa thực tế. BiLSTM của FPT, mô hình phức tạp nhất trên tài sản biến động mạnh nhất trong nhóm, kém mốc ngây thơ
tới tỉ số \SL{dem/fc_worst_ratio}. Ridge trên 30 trễ kém dự báo ngây thơ ở mọi tài sản: với hơn 200 đặc trưng tuyến
tính, mô hình khớp nhiễu của tập huấn luyện. Tỉ lệ đoán đúng chiều dao động quanh 50\%, cao nhất là LSTM trên vàng
(\SL{forecast/GOLD_VND/lstm/dir_accPct}\%) và thấp nhất là hai mô hình LSTM trên VIC (khoảng 41\%, tức tệ hơn tung
đồng xu). Dự báo ngây thơ không có cột đúng chiều
vì lợi suất 0 không có dấu.

\hinhpdf{fc_test}{1.0}{Giá thực tế và giá dự báo một phiên của BiLSTM trên tập kiểm thử}{fc_test}

Hình~\ref{fig:fc_test} dễ gây hiểu lầm nếu đọc thiếu cẩn thận: đường dự báo bám sát đường thực tế, nhưng đó là vì
$\hat P_{t+1} = P_t e^{\hat r_{t+1}}$ và lợi suất một phiên chỉ vài phần trăm, nên bất kỳ dự báo nào gần giá hôm qua
cũng trông khớp. Hình~\ref{fig:fc_scatter} đặt lợi suất dự báo cạnh lợi suất thực tế và cho thấy hai kiểu thất bại
khác nhau. Với vàng quy đổi, BiLSTM \emph{co hẳn về 0}: độ lệch chuẩn của dự báo chỉ \SL{dem/fc_std_pred_GOLD_VND}\% so với
\SL{dem/fc_std_true_GOLD_VND}\% của thực tế, tức mô hình tự học ra dự báo ngây thơ, và tỉ số RMSE vì vậy sát 1. Với
FPT, BiLSTM làm điều ngược lại: dự báo dao động mạnh (độ lệch chuẩn \SL{dem/fc_std_pred_FPT}\% so với
\SL{dem/fc_std_true_FPT}\%) và lệch dương trung bình \SL{dem/fc_mean_pred_FPT}\% mỗi phiên, nhưng các
dao động đó không trùng chiều với thực tế, nên mỗi lần đoán mạnh là một lần cộng thêm sai số; đó là lý do tỉ số RMSE
của nó tệ nhất bảng. Cả hai kiểu đều dẫn tới cùng kết luận: lợi suất ngày mai gần như không dự đoán được từ thông tin
công khai của 30 phiên trước, và mô hình nào ``dám'' dự báo xa khỏi 0 thì bị phạt.

\hinhpdf{fc_scatter}{1.0}{Lợi suất dự báo và lợi suất thực tế trên tập kiểm thử: FPT dự báo mạnh nhưng sai chiều,
vàng quy đổi co về 0}{fc_scatter}

Kết quả này phù hợp với giả thuyết thị trường hiệu quả dạng yếu~\cite{fama1970efficient}: nếu giá đã phản ánh mọi
thông tin trong lịch sử giá, không mô hình nào chỉ dùng lịch sử giá có thể dự báo lợi suất tốt hơn bước ngẫu nhiên.
Thị trường Việt Nam được xem là kém hiệu quả hơn thị trường phát triển, nhưng ở tầm một phiên và với các cổ phiếu
thanh khoản cao nhất VN30, phần dự báo được, nếu có, nhỏ hơn nhiễu. Bài không coi đây là thất bại cần che đi, mà là
một kết quả thực nghiệm quyết định cách tác tử dùng mô hình.

\subsection{Khoảng dự báo cho tác tử}
Công cụ \texttt{du\_bao\_bilstm} của tác tử không chỉ trả một con số. Lớp \texttt{Forecaster} lấy cửa sổ 30 phiên
mới nhất và chạy 200 lượt với Dropout bật, mỗi lượt là một mẫu từ phân phối hậu nghiệm xấp xỉ của trọng số theo MC
Dropout~\cite{gal2016dropout}. Tuy nhiên MC Dropout chỉ đo \emph{bất định của mô hình}; với chuỗi lợi suất, phần này
rất nhỏ so với nhiễu vốn có của giá. Phiên bản đầu của công cụ dùng riêng MC Dropout và báo ``xác suất tăng 100\%''
với khoảng 90\% rộng chưa tới $\pm 0{,}2\%$ quanh giá hiện tại, một kết luận tự tin sai hoàn toàn. Bản cuối cộng vào
mỗi mẫu một nhiễu phần dư có độ lệch chuẩn bằng RMSE trên tập kiểm thử:

\begin{equation}\label{eq:khoangdubao}
  r^{(s)} = f_{\theta^{(s)}}(\mathbf{x}_{t-29:t}) + \varepsilon^{(s)}, \quad
  \varepsilon^{(s)} \sim \mathcal{N}\!\left(0, \mathrm{RMSE}^2_{\text{kiểm thử}}\right), \quad s = 1, \dots, 200,
\end{equation}

\noindent rồi lấy phân vị 5\% và 95\% của $P_t e^{r^{(s)}}$ làm khoảng dự báo 90\% và tỉ lệ $r^{(s)} > 0$ làm xác suất
tăng. Bảng~\ref{tab:fclive} đặt hai nguồn bất định cạnh nhau: với FPT, độ lệch chuẩn của MC Dropout chỉ
\SL{forecast_live/FPT/mc_std_pct}\% trong khi RMSE kiểm thử là \SL{forecast_live/FPT/rmse_pct}\%, tức nhiễu của dữ
liệu lớn hơn bất định của mô hình hơn hai mươi lần. Xác suất tăng sau khi sửa nằm trong khoảng 50 -- 60\%, đúng với
mức hiểu biết thật của mô hình.

\maNguon{dl_models__Forecaster.predict}{Công cụ dự báo của tác tử: MC Dropout cộng nhiễu phần dư cho khoảng dự báo
90\% và xác suất tăng}

\bang{tab_fc_live}{Dự báo phiên tới của năm tài sản: bất định mô hình so với sai số dự báo thực tế}{fclive}

Văn bản công cụ trả cho tác tử luôn đi kèm hai chỉ số kiểm thử: tỉ lệ đoán đúng chiều và tỉ số RMSE so với dự báo
ngây thơ, cùng câu ``không vượt dự báo giữ nguyên giá'' khi tỉ số không dưới 1. Như vậy câu trả lời cuối cùng của tác
tử tự mang theo mức độ tin cậy của chính nó.

% ---------------------------------------------------------------------
\section{Tình huống minh họa: tác tử trả lời nhà đầu tư}\label{sec:tc-tinhhuong}
\subsection{Câu hỏi tác động vĩ mô}
Câu hỏi ``Tỷ giá USD/VND tăng mạnh thì cổ phiếu FPT và giá vàng SJC ảnh hưởng thế nào?'' cần cả ba lớp. Liên kết
thực thể nhận ra \textsf{USD\_VND}, \textsf{FPT} và \textsf{Vang\_SJC}; bộ lập kế hoạch gọi \texttt{kg\_tai\_chinh}
để suy luận hai cặp tác động, \texttt{du\_lieu\_thi\_truong} để lấy giá FPT gần nhất, \texttt{quy\_doi\_vang} để tính
giá quy đổi theo~\eqref{eq:vangquydoi}, và \texttt{truy\_xuat\_rag} để lấy đoạn văn bản giải thích cơ chế. Khung
dưới đây là vết suy luận thật của tác tử, do \texttt{build.py} sinh từ \texttt{agent.json}.

\begin{vetsuyluan}{Vết suy luận ReAct thật cho câu hỏi tài chính (kịch bản 9)}
\input{generated/trace_finance}
\end{vetsuyluan}

\begin{vetsuyluan}{Câu trả lời cuối của tác tử (bản tổng hợp theo khuôn, có trích dẫn [n])}
\input{generated/answer_finance}
\end{vetsuyluan}

Bản nháp của Qwen2.5-1.5B cho câu hỏi này không sai số liệu nào (trung thực số học 1,00) nhưng không ghi số hiệu
bằng chứng [n] nào, nên bộ kiểm chứng từ chối và trả bản tổng hợp theo khuôn; Mục~\ref{sec:agent-ketqua} phân tích
hiện tượng này trên toàn bộ kịch bản. Bản cuối đặt cạnh nhau hai loại bằng chứng: kết luận suy luận từ KG (FPT và
vàng SJC cùng chịu tác động cùng chiều của tỷ giá, mỗi kết luận kèm kênh truyền dẫn và cơ chế) và số liệu tính từ dữ
liệu (giá FPT, giá vàng quy đổi), mỗi dòng mang số hiệu [n] trỏ về bằng chứng gốc. Hình~\ref{fig:ui_finance} là cùng
câu hỏi trên Web App.

\anhUI{ui_finance}{Trợ lý tài chính trên Web App: câu hỏi tác động tỷ giá, vết suy luận và câu trả lời có trích
dẫn}{ui_finance}

\subsection{Câu hỏi dự báo}
Với ``Dự báo giá cổ phiếu FPT phiên tới?'', bộ lập kế hoạch nhận ra ý định dự báo và gọi \texttt{du\_bao\_bilstm} sau
\texttt{du\_lieu\_thi\_truong} (Hình~\ref{fig:ui_finance_forecast}). Câu trả lời nêu giá dự báo, khoảng dự báo 90\%,
xác suất tăng, và ngay sau đó là hai chỉ số kiểm thử cho thấy mô hình không vượt dự báo ngây thơ. Một nhà đầu tư đọc
câu trả lời này biết chính xác nên tin con số tới đâu, điều mà một câu trả lời chỉ có ``giá dự kiến 63.315 đồng''
không làm được.

\anhUI{ui_finance_forecast}{Câu hỏi dự báo FPT: tác tử gọi BiLSTM và trả kèm khoảng dự báo, xác suất tăng và độ tin
cậy trên tập kiểm thử}{ui_finance_forecast}

\subsection{Khám phá đồ thị tri thức trên giao diện}
Trang ``Đồ thị tri thức'' của Web App vẽ \texttt{VietnamFinanceKG} bằng SVG thuần, tô màu nút theo lớp, cạnh ngược
chiều màu đỏ (Hình~\ref{fig:ui_kg_finance}). Người dùng nhấn vào một nút để xem lớp của nó và mọi cạnh vào ra kèm
tài liệu nguồn; đây là cách nhanh nhất để kiểm tra vì sao tác tử kết luận một chiều tác động.

\anhUI{ui_kg_finance}{Trang khám phá \texttt{VietnamFinanceKG} trên Web App}{ui_kg_finance}

% ---------------------------------------------------------------------
\section{Tổng kết chương}
\begin{itemize}
  \item KG tài chính với \SL{kg/finance/nodes} thực thể và \SL{kg/finance/edges} cạnh có dấu, có nguồn, đủ để suy
        luận đa bước và lan truyền cú sốc; khi các kênh mâu thuẫn, bộ suy luận trả ``hỗn hợp'' thay vì ép một chiều.
  \item \SL{dem/edges_agree} trên \SL{dem/edges_checked} cạnh kiểm chứng được khớp dấu với dữ liệu; cạnh tỷ giá
        $\to$ FPT chỉ khớp khi xét độ trễ khoảng \SL{dem/lag_peak_FPT} phiên.
  \item BiLSTM không thắng có ý nghĩa dự báo ngây thơ ở tầm một phiên, phù hợp giả thuyết thị trường hiệu quả dạng
        yếu; giá trị của nó với tác tử nằm ở khoảng dự báo trung thực và các chỉ số tin cậy đi kèm, không ở con số
        điểm.
  \item Khoảng dự báo chỉ dựa trên MC Dropout tự tin sai; cộng nhiễu phần dư theo~\eqref{eq:khoangdubao} đưa xác
        suất tăng về đúng mức 50 -- 60\%.
\end{itemize}

% =====================================================================
%  CHƯƠNG 6 — TRIỂN KHAI HỆ THỐNG TÁC TỬ THÔNG MINH
% =====================================================================
\chapter{Triển khai hệ thống tác tử thông minh}\label{chap:agent}

Hai chương trước dựng các thành phần riêng lẻ: KG y tế và MLP nguy cơ, KG tài chính và BiLSTM, cùng kho văn bản và
bộ truy xuất của Chương~\ref{chap:rag}. Chương này ghép chúng thành một tác tử duy nhất nhận câu hỏi tiếng Việt tự
do, tự quyết định miền và công cụ cần dùng, ghi lại toàn bộ vết suy luận, rồi viết câu trả lời có trích dẫn và tự
kiểm chứng trước khi trả cho người dùng. Phần cuối chương đánh giá tác tử trên \SL{dem/agent_scen} kịch bản và so
với chính LLM đó khi trả lời một mình.

% ---------------------------------------------------------------------
\section{Mô hình ReAct}\label{sec:react}
\subsection{Suy luận đan xen hành động}
ReAct~\cite{yao2023react} kết hợp hai dòng nghiên cứu: suy luận từng bước bằng chuỗi ý nghĩ~\cite{wei2022cot} và mô
hình ngôn ngữ biết gọi công cụ~\cite{schick2023toolformer}. Thay vì trả lời ngay, tác tử lặp ba bước:
\textbf{Thought} (nói ra mình cần biết gì), \textbf{Action} (gọi một công cụ với tham số cụ thể) và
\textbf{Observation} (đọc kết quả công cụ trả về). Vết Thought -- Action -- Observation vừa là bộ nhớ làm việc của
tác tử, vừa là bản giải trình cho người dùng: mỗi kết luận cuối cùng đều truy được về một quan sát cụ thể.

\hinhtikz{t6_react}{Vòng lặp ReAct: Thought $\to$ Action $\to$ Observation lặp quanh bộ nhớ vết suy luận, kết thúc
bằng bước tổng hợp và kiểm chứng}{t6_react}

Đối với bài toán y tế và tài chính, ưu điểm quan trọng nhất của ReAct là \emph{tách tri thức khỏi mô hình ngôn
ngữ}. LLM không cần nhớ ngưỡng eGFR của metformin hay hệ số 1,2057; nó chỉ cần đọc các con số đó trong quan sát. Điều
này trực tiếp giảm hiện tượng ảo giác~\cite{ji2023hallucination}, vốn nguy hiểm nhất đúng ở hai miền này, như
Mục~\ref{sec:agent-ketqua} sẽ cho thấy bằng ví dụ thật.

\subsection{Bộ công cụ}
Tác tử có \SL{dem/tools} công cụ thuộc bốn loại (Bảng~\ref{tab:agenttools}): học sâu (DL), đồ thị tri thức (KG),
truy xuất văn bản (RAG) và tính toán xác định (CALC). Mỗi công cụ trả về hai thứ: một đối tượng quan sát có cấu trúc
cho bộ lập kế hoạch và giao diện, và một danh sách bằng chứng dạng \texttt{(nguồn, văn bản)} để LLM trích dẫn. Nhãn
nguồn ghi rõ loại và gốc của bằng chứng: \texttt{KG:f05-fpt}, \texttt{RAG:h03-ruou-bia}, \texttt{BiLSTM:FPT},
\texttt{Dữ liệu:FPT}.

\bang{tab_agent_tools}{Bảy công cụ của tác tử, số lần được gọi và độ trễ trung vị trên \SL{dem/agent_scen}
kịch bản}{agenttools}

\maNguon{agent__Tool}{Đặc tả một công cụ: tên, loại, mô tả và hàm thực thi}

% ---------------------------------------------------------------------
\section{Kiến trúc hệ thống}\label{sec:kientruc-agent}
\subsection{Năm tầng}
Hệ thống được tổ chức thành năm tầng (Hình~\ref{fig:t6_layers}). Tầng giao diện (Web App) chỉ gửi câu hỏi và hồ sơ
qua HTTP. Tầng điều phối (\texttt{ReActAgent}) phân tích câu hỏi, lập kế hoạch và chạy vòng lặp. Tầng công cụ bọc các
mô hình và đồ thị thành bảy hàm có cùng giao diện. Tầng tri thức gồm hai KG, kho văn bản đã chia đoạn và chỉ mục
truy xuất. Tầng mô hình gồm MLP, BiLSTM, bộ nhúng e5, bộ xếp hạng lại bge và LLM Qwen2.5-1.5B~\cite{qwen2024qwen25},
tất cả chạy cục bộ, không gọi dịch vụ ngoài.

\hinhtikz{t6_layers}{Kiến trúc năm tầng của tác tử và lớp, hàm hiện thực tương ứng trong gói \texttt{tl2}}{t6_layers}

\subsection{Luồng xử lý một câu hỏi}
Hình~\ref{fig:t6_sequence} là biểu đồ tuần tự của một câu hỏi y tế có kèm hồ sơ xét nghiệm. Năm bước chính:

\begin{enumerate}
  \item \textbf{Liên kết thực thể}: \texttt{EntityLinker} (Chương~\ref{chap:rag}) quét câu hỏi bằng bí danh của mọi
        thực thể trong hai KG, ưu tiên bí danh dài nhất và so khớp cả dạng bỏ dấu, nên ``tieu duong'' và ``tiểu
        đường'' đều được nhận.
  \item \textbf{Nhận miền}: miền là KG có nhiều thực thể được liên kết hơn; câu hỏi không liên kết được thực thể nào
        mặc định là y tế.
  \item \textbf{Lập kế hoạch}: chọn chuỗi công cụ theo loại thực thể và ý định (Mục~\ref{sec:laplkh}).
  \item \textbf{Thực thi}: gọi lần lượt từng công cụ, mỗi lần ghi một cặp Thought -- Action và một Observation.
  \item \textbf{Tổng hợp và kiểm chứng}: LLM viết câu trả lời từ bằng chứng; bộ kiểm chứng quyết định dùng bản của
        LLM hay bản tổng hợp theo khuôn (Mục~\ref{sec:kiemchung-agent}).
\end{enumerate}

\hinhtikz{t6_sequence}{Biểu đồ tuần tự cho một câu hỏi y tế có hồ sơ xét nghiệm: trình duyệt, Flask, tác tử, các
công cụ, LLM và bộ kiểm chứng}{t6_sequence}

% ---------------------------------------------------------------------
\section{Hiện thực tác tử}\label{sec:hienthuc-agent}
\subsection{Bộ lập kế hoạch tất định}\label{sec:laplkh}
Một quyết định thiết kế quan trọng là \emph{không} để LLM tự chọn công cụ. Với 1,5 tỉ tham số, Qwen2.5 sinh lời gọi
công cụ không ổn định: sai tên, sai tham số hoặc bỏ qua công cụ cần thiết. Thay vào đó bộ lập kế hoạch dùng luật trên
kết quả liên kết thực thể và từ khóa ý định (Đoạn mã~\ref{code:agent__ReActAgent.plan}):

\begin{itemize}
  \item \textbf{Y tế}: có chỉ số xét nghiệm (HbA1c, glucose, BMI, trích từ câu hỏi bằng biểu thức chính quy hoặc gửi
        kèm từ giao diện) thì gọi \texttt{chan\_doan\_nguy\_co}; có thực thể y tế thì gọi \texttt{kg\_y\_te}.
  \item \textbf{Tài chính}: có biến vĩ mô (tỷ giá, lãi suất, nghị định) thì gọi \texttt{kg\_tai\_chinh}; có mã cổ
        phiếu thì gọi \texttt{du\_lieu\_thi\_truong}; hỏi mức giá vàng thì gọi \texttt{quy\_doi\_vang}; có từ khóa
        xu hướng (``dự báo'', ``phiên tới'', ``nên mua'') thì gọi \texttt{du\_bao\_bilstm}.
  \item \textbf{Luôn luôn}: kết thúc bằng \texttt{truy\_xuat\_rag} để có văn bản gốc làm căn cứ trích dẫn.
\end{itemize}

Đây là biến thể ``lập kế hoạch rồi thực thi'' của ReAct: chuỗi hành động được quyết định một lần từ đầu thay vì sau
mỗi quan sát. Đổi lại tính linh hoạt, bộ lập kế hoạch đạt độ chính xác chọn công cụ gần tuyệt đối
(Mục~\ref{sec:agent-ketqua}) và hoàn toàn tái lập được, điều cần thiết khi đánh giá.

\maDai{agent__ReActAgent.plan}{Bộ lập kế hoạch: nhận miền từ thực thể liên kết, chọn chuỗi công cụ theo loại thực
thể và ý định}

\maNguon{agent__ReActAgent.parse_profile}{Trích hồ sơ xét nghiệm từ câu hỏi tiếng Việt tự do}

\subsection{Vòng lặp và các công cụ}
Hàm \texttt{run} thực thi kế hoạch, đo thời gian mỗi công cụ và ghi vết (Đoạn mã~\ref{code:agent__ReActAgent.run}).
Công cụ phức tạp nhất là \texttt{kg\_tai\_chinh}: nó tách thực thể thành \emph{biến nguồn} (yếu tố vĩ mô, chính
sách, sự kiện) và \emph{biến đích} (cổ phiếu, vàng, chỉ số), gọi \texttt{impact\_paths} cho mọi cặp và diễn đạt
đường ngắn nhất thành một câu có kênh truyền dẫn và cơ chế (Đoạn mã~\ref{code:agent__ReActAgent._t_fkg}). Khi câu
hỏi chỉ nêu biến nguồn, công cụ xét mặc định bốn cổ phiếu và vàng SJC.

\maNguon{agent__ReActAgent.run}{Vòng lặp ReAct: mỗi bước một Thought, một Action đo thời gian và một Observation}

\maNguon{agent__ReActAgent._t_fkg}{Công cụ \texttt{kg\_tai\_chinh}: suy luận đa bước cho mọi cặp nguồn -- đích và
diễn đạt thành bằng chứng có nguồn}

\subsection{Tổng hợp có căn cứ và bộ kiểm chứng}\label{sec:kiemchung-agent}
Bằng chứng từ mọi công cụ được đánh số $[1], \dots, [n]$ và đưa vào lời nhắc của LLM cùng chỉ dẫn hệ thống: chỉ dùng
thông tin trong bằng chứng, mỗi ý quan trọng phải kèm số hiệu, không bịa số liệu
(Đoạn mã~\ref{code:llm__SYSTEM_GROUNDED}, \ref{code:llm__grounded_prompt}). Đây là cùng lời nhắc đã dùng cho
GraphRAG ở Chương~\ref{chap:rag}, nên kết quả hai chương so sánh được với nhau. LLM sinh tham lam (greedy) để mọi lần
chạy cho cùng một câu trả lời.

\maNguon{llm__SYSTEM_GROUNDED}{Chỉ dẫn hệ thống cho LLM: chỉ dùng bằng chứng, trích dẫn [n], không bịa số liệu}

\maNguon{llm__grounded_prompt}{Lời nhắc có căn cứ: bằng chứng đánh số, câu hỏi, yêu cầu trích dẫn}

Mô hình nhỏ không phải lúc nào cũng tuân thủ chỉ dẫn, vì vậy câu trả lời đi qua bộ kiểm chứng trước khi tới người
dùng (Hình~\ref{fig:t6_verifier}). Bộ kiểm chứng kiểm ba điều kiện:

\begin{enumerate}
  \item có ít nhất một trích dẫn $[n]$;
  \item mọi trích dẫn trỏ tới bằng chứng có thật ($1 \le n \le$ số bằng chứng);
  \item \emph{độ trung thực số học} không dưới 0,8, với
        \begin{equation}\label{eq:faith}
          \mathrm{Faith} = \frac{\left|\{\text{con số trong câu trả lời có trong bằng chứng hoặc câu hỏi}\}\right|}
                                {\left|\{\text{con số trong câu trả lời}\}\right|}.
        \end{equation}
\end{enumerate}

Con số được chuẩn hóa trước khi so (bỏ dấu ngăn nghìn và dấu thập phân, nên 1.234,5 và 1,234.5 là một), số một
chữ số và số hiệu trích dẫn không tính. Con số người dùng tự nêu (liều 500 mg, HbA1c 7,3) được coi là có căn cứ, vì
nhắc lại dữ kiện của câu hỏi không phải bịa. Không đạt một trong ba điều kiện, tác tử trả bản tổng hợp theo khuôn:
liệt kê nguyên văn từng bằng chứng kèm số hiệu. Bản khuôn kém tự nhiên nhưng đúng theo cấu tạo, và vết suy luận ghi
rõ lý do bản LLM bị từ chối.

\hinhtikz{t6_verifier}{Bộ kiểm chứng: trích dẫn phải trỏ tới bằng chứng có thật và con số phải có trong bằng chứng;
không đạt thì dùng bản tổng hợp theo khuôn}{t6_verifier}

\maNguon{agent__ReActAgent.synthesize}{Tổng hợp: LLM viết từ bằng chứng, bộ kiểm chứng chọn bản LLM hoặc bản
khuôn}

\maNguon{agent__verify}{Bộ kiểm chứng: trích dẫn hợp lệ và độ trung thực số học~\eqref{eq:faith}}

% ---------------------------------------------------------------------
\section{Đánh giá tác tử}\label{sec:agent-ketqua}
\subsection{Thiết kế thí nghiệm}
Bộ đánh giá gồm \SL{dem/agent_scen} kịch bản, \SL{dem/agent_health} y tế và \SL{dem/agent_finance} tài chính, bao
phủ mọi công cụ và các loại câu hỏi đã phân tích ở hai chương trước. Mỗi kịch bản có ba thứ được viết trước khi chạy:
câu hỏi, \emph{tập công cụ kỳ vọng} và một danh sách \emph{dữ kiện bắt buộc}, mỗi dữ kiện là một nhóm cụm từ đồng
nghĩa (chẳng hạn \{``tránh'', ``không nên'', ``không uống''\}). Các chỉ số:

\begin{itemize}
  \item \textbf{Precision và recall công cụ}: so tập công cụ đã gọi với tập kỳ vọng.
  \item \textbf{Độ phủ dữ kiện}: tỉ lệ nhóm dữ kiện bắt buộc có ít nhất một cụm từ xuất hiện trong câu trả lời cuối
        (so khớp không dấu), cùng hàm đã dùng cho GraphRAG.
  \item \textbf{Trung thực số học}~\eqref{eq:faith} và \textbf{tỉ lệ bản LLM qua kiểm chứng}.
  \item \textbf{Đối chứng ``chỉ LLM''}: cùng Qwen2.5-1.5B trả lời cùng câu hỏi, không có công cụ, với chỉ dẫn ``trả
        lời ngắn gọn, chính xác''.
\end{itemize}

\subsection{Kết quả tổng hợp}
\bang{tab_agent_agg}{Kết quả tổng hợp của tác tử theo miền}{agentagg}

\hinhpdf{agent_eval}{1.0}{Các chỉ số của tác tử theo miền}{agent_eval}

Bảng~\ref{tab:agentagg} và Hình~\ref{fig:agent_eval} cho ba kết luận. Thứ nhất, bộ lập kế hoạch tất định gọi đủ mọi
công cụ cần thiết (recall \SL{agent/agg/all/tool_recall}) và nhận đúng miền ở mọi kịch bản. Precision tài chính dưới
1 vì hai kịch bản (9 và 15) gọi thêm \texttt{quy\_doi\_vang} khi câu hỏi nhắc tới giá vàng mà tập kỳ vọng không có; lời gọi
thừa này vô hại và thực ra bổ sung thông tin. Thứ hai, độ phủ dữ kiện của tác tử là \SL{agent/agg/all/coveragePct}\%,
so với \SL{agent/agg/all/coverage_llm_onlyPct}\% khi LLM trả lời một mình: công cụ nâng độ phủ thêm
\SL{dem/agent_cov_gainPct} điểm phần trăm. Thứ ba, và đáng bàn nhất, chỉ \SL{agent/agg/all/llm_passPct}\% bản nháp
của LLM qua kiểm chứng.

\bang{tab_agent_scen}{Kết quả từng kịch bản}{agentscen}

\subsection{Vì sao LLM trả lời một mình là nguy hiểm}
Độ phủ thấp của đối chứng ``chỉ LLM'' không chỉ là thiếu thông tin. Đọc từng câu trả lời cho thấy các lỗi có hậu quả
thật:

\begin{itemize}
  \item Kịch bản 1, câu hỏi về bia khi đang dùng metformin: LLM trả lời ``Dùng Metformin không ảnh hưởng đến việc
        uống bia'', bỏ qua hoàn toàn nguy cơ nhiễm toan lactic mà Quyết định 5481/QĐ-BYT cảnh báo.
  \item Kịch bản 8, đái tháo đường kèm suy tim: LLM liệt kê đầu tiên nhóm thiazolidinedione (pioglitazone), là nhóm
        thuốc gây giữ nước và \emph{chống chỉ định} ở người suy tim, và không nhắc tới nhóm SGLT2 được ưu tiên. Tác tử
        trả lời đúng ``ưu tiên SGLT2 (empagliflozin)'' với một trích dẫn tới KG.
  \item Kịch bản 10, giá vàng quy đổi hôm nay: LLM tự nhận thời điểm là ``tháng 10/2023'' và trả ``645 -- 650
        USD/ounce, tương đương 21,89 -- 22,17 triệu đồng mỗi lượng 1 gram'', sai cả mức giá (giá thật
        \SL{market/gold/last/xau} USD/ounce), cả đơn vị (lượng là 37,5 g). Mô hình không có cách nào biết giá hôm
        nay; tác tử thì tính đúng \SL{market/gold/last/gold_vnd} triệu đồng/lượng từ dữ liệu mới nhất.
  \item Kịch bản 13, lãi suất điều hành và VCB, VIC: LLM gọi VIC là ``ngân hàng thương mại''.
\end{itemize}

\hinhpdf{agent_vs_llm}{1.0}{Độ phủ dữ kiện theo từng kịch bản: tác tử có công cụ so với LLM trả lời một
mình}{agent_vs_llm}

Trung thực số học của đối chứng (\SL{agent/agg/all/faith_llm_only}) trông không thấp, nhưng phải đọc đúng cách: 11
trên 16 câu trả lời chỉ LLM không có con số hai chữ số nào, và theo định nghĩa~\eqref{eq:faith} một câu trả lời
không có con số được tính trung thực tuyệt đối. Khi LLM thật sự đưa ra con số, như ở kịch bản 10, tỉ lệ rơi xuống
0,17. Đây là ví dụ cho thấy một chỉ số trung thực phải được đọc cùng độ phủ, không đọc riêng.

\subsection{Vì sao bản nháp LLM bị từ chối}
Trong \SL{dem/agent_scen} kịch bản, chỉ \SL{dem/agent_llm} bản nháp qua kiểm chứng (kịch bản 6 và 8, cả hai đều ngắn,
một ý, một trích dẫn). \SL{dem/agent_nocite} bản nháp bị từ chối vì \emph{không có trích dẫn nào}, dù hầu hết trung
thực về số liệu (chỉ một bản dưới 1,0) và nhiều bản đúng về nội dung. Một bản nháp còn lại (kịch bản 7) có ba trích dẫn nhưng một phần ba số con số
không tìm thấy trong bằng chứng. Nguyên nhân chính vì vậy không phải LLM bịa số, mà là mô hình 1,5B tham số không
tuân thủ yêu cầu định dạng ``kèm số hiệu [n]'' khi bằng chứng dài và nhiều nguồn.

\begin{vetsuyluan}{Bản nháp của Qwen2.5-1.5B cho kịch bản 9 (bị từ chối vì không có trích dẫn)}
\input{generated/llmdraft_finance}
\end{vetsuyluan}

Bản nháp trên minh họa thêm một giới hạn của bộ kiểm chứng: nó không có con số nào nên đạt trung thực số học tuyệt
đối, nhưng câu ``giá vàng miếng SJC cũng tăng do giá quy đổi tăng nhanh hơn chi phí'' đã trộn cơ chế của FPT vào vàng,
và câu cuối viện dẫn ``giá nguyên vật liệu, công nghệ, thị trường xuất khẩu'' không có trong bằng chứng nào. Kiểm
chứng con số bắt được lỗi bịa số liệu, nhưng không bắt được lỗi trộn ngữ nghĩa; việc thiếu trích dẫn vô tình chặn
được bản nháp này. Một bộ kiểm chứng mạnh hơn cần đối chiếu từng câu với bằng chứng được trích, chẳng hạn bằng mô
hình suy diễn ngôn ngữ tự nhiên (NLI); Chương~\ref{chap:ketluan} bàn hướng này.

\subsection{Ba kịch bản chưa đạt độ phủ đầy đủ}
Ba kịch bản có độ phủ 0,50, mỗi kịch bản chỉ ra một giới hạn khác nhau:

\begin{itemize}
  \item \textbf{Kịch bản 4} (HbA1c 5,9\%): câu trả lời không có cụm ``tiền đái tháo đường''. KG đã có tri thức cần
        thiết (cạnh \textsf{HbA1c\_Tien} $\xrightarrow{\textsf{TIEU\_CHUAN\_CHAN\_DOAN}}$ \textsf{Tien\_DTD} theo
        Bảng~\ref{tab:chandoan}), nhưng công cụ \texttt{kg\_y\_te} chỉ tra theo thực thể có tên trong câu hỏi, không
        đối chiếu chỉ số xét nghiệm của hồ sơ với ngưỡng chẩn đoán. Đây là một bước nối giữa tầng học sâu và tầng đồ
        thị nên được bổ sung.
  \item \textbf{Kịch bản 6} (suy thận nặng): bản LLM qua kiểm chứng, nói đúng ``không nên dùng metformin khi eGFR
        $<$ 30'', nhưng dùng ``không nên'' thay cho ``chống chỉ định'', nên bị tính thiếu một dữ kiện. Đây là giới hạn
        của độ phủ đo bằng so khớp cụm từ.
  \item \textbf{Kịch bản 15} (Fed và tỷ giá): \textsf{USD\_VND} là \textsf{MacroFactor} nên bộ lập kế hoạch xếp nó vào
        biến nguồn, không xét cặp Fed $\to$ USD/VND. Một biến có thể vừa là nguyên nhân vừa là kết quả; bộ lập kế
        hoạch cần xét vai trò theo vị trí trong câu, không chỉ theo lớp của thực thể.
\end{itemize}

\subsection{Độ trễ}
Tổng thời gian trung bình một câu hỏi là \SL{agent/agg/all/total_ms} ms, trong đó LLM chiếm phần lớn
(\SL{agent/agg/all/llm_s} giây cho bước tổng hợp). Hình~\ref{fig:agent_tools} so sánh các công cụ: truy vấn KG và
tính toán chưa tới 1 ms, MLP khoảng \SL{agent/agg/tool_ms/chan_doan_nguy_co} ms, BiLSTM với 200 mẫu MC Dropout khoảng
\SL{agent/agg/tool_ms/du_bao_bilstm} ms, và truy xuất RAG khoảng \SL{agent/agg/tool_ms/truy_xuat_rag} ms vì phải
nhúng câu hỏi và chạy bộ xếp hạng lại trên 30 ứng viên. Nếu bỏ LLM và chỉ dùng bản khuôn, tác tử trả lời trong
chưa tới một giây.

\hinhpdf{agent_tools}{1.0}{Số lần gọi và độ trễ trung vị của từng công cụ}{agent_tools}

% ---------------------------------------------------------------------
\section{Ứng dụng Web}\label{sec:webapp}
\subsection{Kiến trúc triển khai}
Web App gồm một máy chủ Flask và một trang HTML với JavaScript thuần, không dùng thư viện giao diện ngoài; đồ thị tri
thức và biểu đồ giá được vẽ bằng SVG ngay trên trình duyệt. Máy chủ nạp tác tử một lần khi khởi động và cung cấp bốn
API (Hình~\ref{fig:t6_arch}):

\begin{itemize}
  \item \texttt{POST /api/agent}: chạy tác tử, trả vết suy luận, bằng chứng, câu trả lời và kết quả kiểm chứng
        (Đoạn mã~\ref{code:app__api_agent});
  \item \texttt{GET /api/kg/<tên>}: trả đồ thị kèm tọa độ Kamada-Kawai tính sẵn ở máy chủ
        (Đoạn mã~\ref{code:app__kg_payload});
  \item \texttt{POST /api/rag}: chạy ba bộ truy xuất trên cùng câu hỏi để so sánh;
  \item \texttt{GET /api/series/<mã>}: trả 250 phiên giá gần nhất cho biểu đồ.
\end{itemize}

\hinhtikz{t6_arch}{Kiến trúc triển khai: trình duyệt, máy chủ Flask với bốn API, tác tử và phân bổ CPU/GPU cho
từng thành phần}{t6_arch}

Mọi mô hình học sâu chạy trên GPU trong cùng tiến trình (Bảng~\ref{tab:moitruong}): LLM bf16, bộ nhúng e5 fp16, bộ
xếp hạng lại bf16, MLP và BiLSTM; BM25, KG và các công cụ tính toán chạy trên CPU.

\maNguon{app__api_agent}{API tác tử: nhận câu hỏi và hồ sơ tùy chọn, trả vết suy luận và câu trả lời đã kiểm
chứng}

\maNguon{app__kg_payload}{API đồ thị tri thức: bố cục Kamada-Kawai tính sẵn, cạnh diễn đạt thành câu kèm nguồn}

\subsection{Giao diện}
Trang có bốn thẻ: Trợ lý Y tế, Trợ lý Tài chính, Khám phá KG và Tra cứu RAG (Hình~\ref{fig:ui_overview}). Hai thẻ
trợ lý có cùng bố cục: khung câu hỏi với các câu mẫu, khung vết suy luận tô màu theo loại bước, khung câu trả lời
ghi rõ bản LLM hay bản khuôn cùng thời gian xử lý, và danh sách bằng chứng được trích. Thẻ y tế có thêm biểu mẫu hồ sơ
xét nghiệm và đồng hồ nguy cơ; thẻ tài chính có biểu đồ giá 250 phiên. Ảnh chụp của hai thẻ trợ lý đã được trình bày
ở Hình~\ref{fig:ui_health}, \ref{fig:ui_health_food}, \ref{fig:ui_finance} và~\ref{fig:ui_finance_forecast}.

\anhUI{ui_overview}{Màn hình chính của Web App: bốn thẻ chức năng và trợ lý y tế với hồ sơ xét nghiệm}{ui_overview}

Thẻ Khám phá KG (Hình~\ref{fig:ui_kg_health}) cho phép chọn một trong hai KG, nhấn vào nút để xem các cạnh vào ra
kèm nguồn. Thẻ Tra cứu RAG (Hình~\ref{fig:ui_rag}) chạy BM25, Dense e5-base và Hybrid + Rerank trên cùng câu hỏi,
hiển thị ba đoạn đứng đầu của mỗi bộ cùng điểm số và thời gian, để người dùng tự thấy kết quả của
Chương~\ref{chap:rag}: BM25 nhạy với từ khóa, dense hiểu đồng nghĩa, và kết hợp có xếp hạng lại cho kết quả tốt
nhất.

\anhUI{ui_kg_health}{Thẻ Khám phá KG với \texttt{VietnamDiabetesKG}: nút tô theo lớp, cạnh nguy cơ màu đỏ,
khung bên liệt kê các cạnh của nút đang chọn}{ui_kg_health}

\anhUI{ui_rag}{Thẻ Tra cứu RAG: ba bộ truy xuất trên cùng một câu hỏi}{ui_rag}

% ---------------------------------------------------------------------
\section{Tổng kết chương}
\begin{itemize}
  \item Tác tử ReAct với bộ lập kế hoạch tất định gọi đủ công cụ ở mọi kịch bản và nâng độ phủ dữ kiện từ
        \SL{agent/agg/all/coverage_llm_onlyPct}\% lên \SL{agent/agg/all/coveragePct}\% so với cùng LLM trả lời một
        mình.
  \item LLM trả lời một mình mắc lỗi nguy hiểm ở cả hai miền (khẳng định metformin không ảnh hưởng khi uống bia, đề
        xuất thuốc chống chỉ định cho người suy tim); công cụ và trích dẫn loại bỏ được các lỗi này.
  \item Bộ kiểm chứng từ chối phần lớn bản nháp vì thiếu trích dẫn; bản khuôn giữ cho câu trả lời luôn truy vết được,
        nhưng kiểm chứng con số không bắt được lỗi trộn ngữ nghĩa.
  \item Ba kịch bản thiếu dữ kiện chỉ ra ba hướng cải tiến cụ thể: luật tiền đái tháo đường trong KG, đo độ phủ theo
        ngữ nghĩa, và lập kế hoạch theo vai trò của thực thể trong câu.
\end{itemize}

% =====================================================================
%  CHƯƠNG 7 — KẾT LUẬN VÀ HƯỚNG PHÁT TRIỂN
% =====================================================================
\chapter{Kết luận và hướng phát triển}\label{chap:ketluan}

\section{Kết luận}
Tiểu luận đã nghiên cứu đồ thị tri thức, sinh tăng cường truy xuất và tác tử thông minh, rồi hiện thực hóa toàn bộ
thành một hệ thống chạy được cho hai bài toán của Việt Nam: tư vấn cho người bệnh đái tháo đường típ 2 và phân tích
tác động vĩ mô lên cổ phiếu VN30 và vàng SJC. Mọi con số trong báo cáo được sinh tự động từ kết quả thí nghiệm, và mọi
đoạn mã được trích từ chính mã nguồn đã chạy ra các con số đó. Kết quả chính:

\begin{enumerate}
  \item \textbf{Đồ thị tri thức.} Hai KG có bản thể học riêng, ràng buộc miền -- miền giá trị và nguồn văn bản cho
        từng cạnh: \texttt{VietnamDiabetesKG} (\SL{kg/health/nodes} thực thể, \SL{kg/health/edges} quan hệ) và
        \texttt{VietnamFinanceKG} (\SL{kg/finance/nodes} thực thể, \SL{kg/finance/edges} quan hệ). Suy diễn RDFS làm
        \SL{dem/sparql} truy vấn SPARQL trả đúng kết quả, trong đó có truy vấn chỉ có kết quả \emph{sau} suy diễn. Với
        câu hỏi có điều kiện số trên dữ liệu bảng, truy vấn KG trả đúng
        \SL{kg/tabular/numeric_query/kg_hits}/\SL{kg/tabular/numeric_query/truth} hồ sơ, trong khi tìm kiếm văn bản
        TF-IDF chỉ đưa \SL{kg/tabular/numeric_query/tfidf_top20_correct} hồ sơ đúng vào top-20.
  \item \textbf{Truy xuất.} Trên \SL{retrieval_grid/n_questions} câu hỏi, chia đoạn theo cấu trúc kết hợp BM25,
        e5 và xếp hạng lại đạt Hit@1 \SL{retrieval_grid/grid/structure/hybrid_rerank/hit@1Pct}\%, so với
        \SL{retrieval_grid/grid/fixed/dense/hit@1Pct}\% của chia đoạn cố định với dense thuần. Hai kỹ thuật RAG nâng
        cao phổ biến lại \emph{làm giảm} chất lượng khi dùng LLM 1,5B: HyDE hạ Hit@1 từ
        \SL{advanced/agg/dense/hit@1Pct}\% xuống \SL{advanced/agg/hyde/hit@1Pct}\%, Multi-Query xuống
        \SL{advanced/agg/multiquery/hit@1Pct}\%. GraphRAG cục bộ nâng độ phủ đáp án từ
        \SL{graphrag/agg/naive/ans_covPct}\% lên \SL{graphrag/agg/graph_local/ans_covPct}\%.
  \item \textbf{Học sâu.} MLP nguy cơ đái tháo đường đạt ROC-AUC \SL{diabetes/mlp/roc_aucPct}\% và PR-AUC
        \SL{diabetes/mlp/pr_aucPct}\% trên tập kiểm thử. BiLSTM dự báo lợi suất không thắng có ý nghĩa dự báo ngây thơ
        (\SL{dem/fc_beat}/\SL{dem/fc_models} mô hình có tỉ số RMSE dưới 1, tốt nhất chỉ giảm
        \SL{dem/fc_best_gainPct}\% sai số), phù hợp giả thuyết thị trường hiệu quả dạng yếu.
  \item \textbf{Tác tử.} Tác tử ReAct với \SL{dem/tools} công cụ gọi đủ công cụ cần thiết ở mọi kịch bản và nâng độ
        phủ dữ kiện từ \SL{agent/agg/all/coverage_llm_onlyPct}\% (LLM trả lời một mình) lên
        \SL{agent/agg/all/coveragePct}\%, đồng thời loại bỏ các lỗi nguy hiểm của LLM như khẳng định metformin không
        ảnh hưởng khi uống bia hay đề xuất thuốc chống chỉ định cho người suy tim.
  \item \textbf{Ứng dụng.} Web App Flask chạy toàn bộ mô hình cục bộ trên một GPU 8 GB, hiển thị vết suy luận, câu
        trả lời có trích dẫn, đồ thị tri thức tương tác và so sánh ba bộ truy xuất.
\end{enumerate}

Bài học chung của tiểu luận là các thành phần bổ trợ nhau đúng như lý thuyết dự đoán: KG cho ràng buộc chính xác và
suy luận nhiều bước, RAG cho văn bản gốc để giải thích và trích dẫn, mô hình học sâu cho đánh giá định lượng từ dữ
liệu của từng người, còn LLM chỉ làm việc nó làm tốt nhất là diễn đạt. Khi LLM được giao thêm việc nhớ tri thức, tự
viết truy vấn giả định (HyDE) hay tự chọn công cụ, chất lượng giảm rõ rệt với mô hình nhỏ.

\section{Hạn chế}
\begin{itemize}
  \item \textbf{LLM nhỏ không tuân thủ định dạng trích dẫn.} Chỉ \SL{agent/agg/all/llm_passPct}\% bản nháp qua
        kiểm chứng; \SL{dem/agent_nocite} bản bị từ chối vì không có số hiệu [n]. Câu trả lời cuối vì vậy thường là
        bản tổng hợp theo khuôn: đúng và truy vết được, nhưng kém tự nhiên.
  \item \textbf{Bộ kiểm chứng chỉ kiểm con số và trích dẫn.} Nó không phát hiện được câu trộn cơ chế hay viện dẫn
        ngoài bằng chứng khi câu đó không chứa con số (Mục~\ref{sec:agent-ketqua}).
  \item \textbf{Xác suất của MLP chưa hiệu chỉnh.} Trọng số lớp dương làm xác suất bị đẩy cao so với tần suất thực;
        con số ``xác suất mắc bệnh'' chỉ nên đọc như điểm xếp hạng nguy cơ.
  \item \textbf{KG viết tay và nhỏ.} Hai KG có vài chục thực thể, cạnh được viết từ kho văn bản thay vì trích xuất
        tự động; dấu cạnh tài chính cố định trong khi tương quan thực tế đổi dấu theo giai đoạn, như cạnh USD/VND $\to$
        FPT chỉ khớp dữ liệu khi xét độ trễ.
  \item \textbf{Bộ lập kế hoạch dựa trên luật.} Nó tái lập được và chính xác trên tập kịch bản, nhưng không thích
        nghi với câu hỏi ngoài mẫu, và xét vai trò nguyên nhân -- kết quả theo lớp của thực thể thay vì theo câu.
  \item \textbf{Quy mô đánh giá.} \SL{retrieval_grid/n_questions} câu hỏi truy xuất, \SL{dem/grag_q} câu hỏi GraphRAG
        và \SL{dem/agent_scen} kịch bản tác tử đủ để so sánh các phương án trong bài, nhưng chưa đủ để kết luận có ý
        nghĩa thống kê cho những chênh lệch nhỏ.
\end{itemize}

\section{Hướng phát triển}
\begin{itemize}
  \item \textbf{LLM tuân thủ trích dẫn tốt hơn}: dùng mô hình 7B lượng tử hóa 4-bit vừa với GPU 8 GB, hoặc tinh chỉnh
        LoRA mô hình 1,5B trên các cặp (bằng chứng, câu trả lời có trích dẫn) sinh từ chính bản tổng hợp theo khuôn.
  \item \textbf{Kiểm chứng theo ngữ nghĩa}: đối chiếu từng câu của câu trả lời với bằng chứng được trích bằng một mô
        hình suy diễn ngôn ngữ tự nhiên (NLI) đa ngữ, thay cho chỉ kiểm con số.
  \item \textbf{Hiệu chỉnh xác suất}: áp dụng temperature scaling hoặc hồi quy isotonic trên tập kiểm định để xác
        suất của MLP phản ánh đúng tần suất mắc bệnh.
  \item \textbf{Nối tầng học sâu với KG}: đối chiếu chỉ số xét nghiệm của hồ sơ với các ngưỡng chẩn đoán đã có trong
        KG, để tác tử tự kết luận ``tiền đái tháo đường'' thay vì chỉ báo điểm nguy cơ.
  \item \textbf{KG tài chính có bằng chứng thực nghiệm}: gắn tương quan, độ trễ và khoảng thời gian hiệu lực vào thuộc
        tính của từng cạnh, cập nhật định kỳ từ dữ liệu, để tác tử nói ``tác động cùng chiều thể hiện sau khoảng hai
        tháng'' thay vì một dấu cố định.
  \item \textbf{Mở rộng KG tự động}: dùng LLM trích xuất bộ ba từ văn bản mới của Bộ Y tế, NHNN và báo cáo doanh
        nghiệp, kiểm tra bằng ràng buộc bản thể học trước khi nạp, và chuyển lưu trữ từ NetworkX sang cơ sở dữ liệu đồ
        thị như Neo4j khi quy mô tăng.
  \item \textbf{Bộ lập kế hoạch lai}: giữ luật làm mặc định, nhưng cho LLM đề xuất thêm công cụ khi luật không khớp,
        với bộ kiểm tra tên và tham số trước khi thực thi.
\end{itemize}


\cleardoublepage
\phantomsection
\addcontentsline{toc}{chapter}{Tài liệu tham khảo}
{\sloppy\printbibliography[title={Tài liệu tham khảo}]}

\end{document}
